\documentclass[11pt]{extarticle}
\usepackage{fontspec}
\setmainfont{DejaVu Serif}
\setsansfont{DejaVu Sans}
\setmonofont{DejaVu Sans Mono}
\defaultfontfeatures{Ligatures=TeX}
\usepackage{amsmath}
\usepackage{unicode-math}
\setmathfont{DejaVu Math TeX Gyre}
\renewcommand{\contentsname}{Оглавление}
\usepackage[a4paper,margin=2.5cm]{geometry}
\usepackage{graphicx}
\usepackage{longtable}
\usepackage{xltabular}
\usepackage{booktabs}
\usepackage{array}
\usepackage{calc}
\usepackage{hyperref}
\hypersetup{colorlinks=true,linkcolor=blue,citecolor=blue,urlcolor=blue}
\usepackage[normalem]{ulem}
\usepackage{setspace}
\singlespacing
\usepackage{titlesec}
\setcounter{secnumdepth}{-2}
\usepackage{indentfirst}
\setlength{\parindent}{1.25cm}
\usepackage{enumitem}
\setlist{nosep}
\providecommand{\tightlist}{\setlength{\itemsep}{0pt}\setlength{\parskip}{0pt}}
\usepackage{ragged2e}

% --- Manual Russian hyphenation (container's language.dat.lua/polyglossia is broken) ---
\input{ru_hyphenation_setup.tex}
\newlanguage\russianlang
\language\russianlang
\lefthyphenmin=2
\righthyphenmin=2
\patterns{}
\input{/usr/share/texlive/texmf-dist/tex/generic/hyph-utf8/patterns/tex/hyph-ru.tex}
\language=\russianlang
\AtBeginDocument{\language=\russianlang}

\tolerance=2000
\emergencystretch=2em
\hbadness=6000
\sloppy

\title{Математическое моделирование осаждения парафинов, асфальтенов и других тяжёлых компонентов нефти при неизотермической фильтрации в пористой среде\\[0.5em]\large Систематический обзор литературы}
\author{}
\date{2026}

\begin{document}
\maketitle
\tableofcontents
\newpage

\section{Введение}\label{ux432ux432ux435ux434ux435ux43dux438ux435}

\subsection{Постановка проблемы}\label{ux43fux43eux441ux442ux430ux43dux43eux432ux43aux430-ux43fux440ux43eux431ux43bux435ux43cux44b}

Добыча нефти из пластов с высоким содержанием парафинов, смол и асфальтенов повсеместно осложняется выпадением твёрдой фазы тяжёлых углеводородных компонентов непосредственно в поровом пространстве коллектора, в стволе скважины и в промысловых трубопроводах. Явление имеет единую физическую природу --- потерю термодинамической устойчивости раствора тяжёлых компонентов в нефти при изменении давления, температуры и состава, --- но проявляется в разных геометриях (пласт, скважина, труба) через разный набор доминирующих механизмов переноса и разный масштаб времени и пространства. В нефтяном пласте ключевая особенность задачи состоит в том, что снижение температуры при закачке холодной воды или газа, дросселирование потока при движении к забою и локальное истощение лёгких фракций при фильтрации многокомпонентной смеси происходят одновременно с самим течением --- то есть фильтрация принципиально неизотермична и многокомпонентна, а тепловое, гидродинамическое и фазовое поля взаимно сопряжены.

Экономические и технологические последствия этого явления значительны: снижение проницаемости призабойной зоны пласта (формейшн дэмидж) приводит к падению продуктивности скважин, рост отложений в стволе и лифтовых трубах требует периодической механической или химической депарафинизации, а неверный прогноз глубины и профиля отложений ведёт либо к избыточным (и дорогостоящим) профилактическим мероприятиям, либо к незапланированным остановкам добычи. Именно поэтому на протяжении более чем полувека --- от первых инженерных корреляций 1960--1980-х годов до современных гибридных физически-информированных моделей --- построение количественных, физически обоснованных математических моделей осаждения тяжёлых компонентов нефти остаётся одной из центральных задач нефтепромысловой термодинамики и подземной гидромеханики.

\subsection{Физическая природа осаждения тяжёлых компонентов: единство и различие механизмов}\label{ux444ux438ux437ux438ux447ux435ux441ux43aux430ux44f-ux43fux440ux438ux440ux43eux434ux430-ux43eux441ux430ux436ux434ux435ux43dux438ux44f-ux442ux44fux436ux451ux43bux44bux445-ux43aux43eux43cux43fux43eux43dux435ux43dux442ux43eux432-ux435ux434ux438ux43dux441ux442ux432ux43e-ux438-ux440ux430ux437ux43bux438ux447ux438ux435-ux43cux435ux445ux430ux43dux438ux437ux43cux43eux432}

Тяжёлые компоненты нефти, ответственные за твердофазные отложения, представляют собой не одно вещество, а совокупность структурно различных классов углеводородов и гетероатомных соединений, для которых, тем не менее, можно построить единое компонентное описание через гомологические ряды и структурно-групповой анализ. Парафины (нормальные и слаборазветвлённые алканы) описываются общей формулой \(C_nH_{2n+2}\) и кристаллизуются из раствора при охлаждении ниже температуры начала кристаллизации (Wax Appearance Temperature, WAT) по классическому механизму твердофазного равновесия между упорядоченной (орторомбической или гексагональной) кристаллической решёткой длинноцепочечных алканов и жидкой нефтяной матрицей. Наиболее длинноцепочечные парафины (\(n \gtrsim 20\)--\(25\)) наименее растворимы и кристаллизуются первыми, что определяет фракционный, а не одноступенчатый характер выпадения твёрдой фазы. Нафтены (\(C_nH_{2n}\), насыщенные циклические структуры) и ароматические углеводороды (\(C_nH_{2n-6}\) и их полициклические аналоги) сами по себе не образуют кристаллическую твёрдую фазу при пластовых условиях, но выступают растворителем по отношению к парафинам и, в существенно большей степени, стабилизирующей средой по отношению к асфальтенам. Смолы и асфальтены --- не индивидуальные соединения, а операционно определяемые (по растворимости в стандартных растворителях) классы наиболее полярных и высокомолекулярных компонентов нефти: асфальтены нерастворимы в лёгких \emph{н}-алканах и растворимы в ароматических растворителях, их молекулярная архитектура представляет собой конденсированные полициклические ароматические ядра с алифатическими заместителями и гетероатомами серы, азота, кислорода, а также следовыми металлопорфиринами никеля и ванадия. В отличие от парафинов, асфальтены в исходном состоянии находятся в нефти не в виде истинного молекулярного раствора, а в виде коллоидной системы --- наноагрегатов и кластеров наноагрегатов, стабилизированных смолами; их дестабилизация носит принципиально иную физическую природу, чем кристаллизация парафинов, и описывается не теорией твердофазного равновесия, а коллоидной и статистической термодинамикой ассоциирующих систем.

Это структурное различие определяет и различие доминирующих термодинамических подходов, систематизированных в настоящем обзоре: для парафинов адекватным и вычислительно экономичным описанием чаще всего служит модель твёрдого раствора или модель множественной твёрдой фазы (multi-solid model) в сочетании с представлением жидкой фазы как регулярного или полимерного (по Флори--Хаггинс) раствора; для асфальтенов требуется либо термодинамика коллоидной дестабилизации (модель Йена--Маллинса, теория мицеллизации), либо статистическая теория ассоциирующих флюидов (SAFT, PC-SAFT), учитывающая полидисперсность и ассоциативные взаимодействия явно. При этом оба класса моделей должны быть выражены через один и тот же набор псевдокомпонентов нефти --- по числу атомов углерода (Single Carbon Number, SCN) или по структурно-групповому PNA/PIONA-разложению --- чтобы термодинамический блок мог быть согласованно встроен в единую композиционную модель многофазной многокомпонентной фильтрации.

Помимо термодинамического равновесия (то есть самого факта существования и количества твёрдой фазы при данных \(P\), \(T\) и составе), критически важна кинетика перехода вещества из мобильной (растворённой или взвешенной) в иммобилизованную (осевшую на скелете породы или стенке канала) форму. Именно кинетический, а не равновесно-термодинамический аспект задачи определяет пространственно-временной профиль отложений: скорость диффузионного переноса растворённых молекул парафина к переохлаждённой стенке, скорость доставки асфальтеновых наноагрегатов к поверхности поры за счёт сдвиговой дисперсии, броуновской диффузии и гравитационного осаждения, а также обратный процесс --- механический вынос (эрозию) уже осевшего материала при повышении локальной скорости или напряжения сдвига. Наконец, накопление осевшей твёрдой фазы снижает пористость и, в существенно большей степени за счёт нелинейной зависимости, проницаемость породы, что через закон Дарси меняет локальное поле скоростей и, через уравнение энергии, температурное поле --- а значит, и сами термодинамические условия дестабилизации в соседних точках пласта. Таким образом, полная математическая модель является принципиально сопряжённой (coupled) системой: термодинамика фазового равновесия → кинетика захвата/выноса → изменение проницаемости и пористости → изменение поля течения и температуры → обратное влияние на термодинамику. Именно эта многоуровневая обратная связь и составляет главный предмет настоящего обзора.

\subsection{Цель, охват и структура обзора}\label{ux446ux435ux43bux44c-ux43eux445ux432ux430ux442-ux438-ux441ux442ux440ux443ux43aux442ux443ux440ux430-ux43eux431ux437ux43eux440ux430}

Целью настоящего обзора является систематизация математических моделей осаждения тяжёлых компонентов нефти (парафинов, асфальтенов и смол) при неизотермической многокомпонентной фильтрации в пористой среде, с акцентом на физико-химическое обоснование каждого класса моделей, а не на общее феноменологическое описание проблемы. В отличие от многих обзоров, ориентированных на промысловые аспекты предотвращения отложений (ингибиторы, тепловая изоляция, механическая очистка), здесь рассматриваются исключительно модельные представления --- от термодинамики фазового равновесия до численных схем и промышленных симуляторов, --- а также их экспериментальная валидация и современные тенденции развития области, включая гибридные физически-информированные подходы.

Обзор построен по следующей логике. Раздел 1 излагает термодинамические основы осаждения тяжёлых компонентов: модели твёрдо-жидкостного равновесия для парафинов (модель твёрдого раствора, модель множественной твёрдой фазы, полимерная теория Флори--Хаггинса), кубические и ассоциативные уравнения состояния, а также принципы компонентной характеризации нефти (SCN, PNA/PIONA, SARA). Раздел 2 переходит к кинетике и механизмам осаждения --- диффузии Фика, сдвиговой дисперсии, броуновскому и гравитационному переносу, а также моделям роста и старения отложений. Раздел 3 формулирует полную сопряжённую систему уравнений неизотермической многофазной фильтрации (неразрывность, закон Дарси, уравнение энергии с учётом скрытой теплоты фазового перехода) в пористой среде. Раздел 4 посвящён обратной связи между осаждением и изменением фильтрационно-ёмкостных свойств породы --- модели Козени--Кармана и её модификациям, а также замкнутой формулировке системы «осаждение → повреждение пласта → течение». Раздел 5 систематизирует специфичную для асфальтенов и смол физику --- коллоидную модель Йена--Маллинса, термодинамику SAFT/PC-SAFT и модели захвата/выноса частиц в поровом пространстве. Раздел 6 обсуждает численные методы решения полученных систем уравнений (конечно-объёмные и конечно-элементные схемы, многомасштабные методы, операторно-ориентированные архитектуры) и даёт сравнительный обзор исследовательских и промышленных симуляторов. Раздел 7 систематизирует экспериментальные методики определения входных параметров моделей (WAT, SARA, кинетические константы) и опубликованные сопоставления модельных прогнозов с лабораторными и промысловыми данными. Наконец, раздел 8 обсуждает современные (2019--2026 гг.) тенденции области --- от детализации компонентного состава до машинного обучения и физически-информированных нейронных сетей --- и формулирует открытые проблемы.

Такая структура отражает естественную иерархию физических масштабов и уровней описания задачи: от молекулярной термодинамики фазового равновесия (разделы 1, 5) через кинетику переноса на масштабе поры и кристалла (разделы 2, 5.4) к континуальному описанию сопряжённой фильтрации на масштабе пласта (разделы 3--4) и, наконец, к практической реализации --- численным методам, симуляторам и экспериментальной проверке (разделы 6--7) --- с заключительным взглядом на то, куда движется область в последние годы (раздел 8).

\newpage

\section{Раздел 1. Термодинамика осаждения тяжёлых компонентов нефти (парафины, асфальтены, смолы)}\label{ux440ux430ux437ux434ux435ux43b-1.-ux442ux435ux440ux43cux43eux434ux438ux43dux430ux43cux438ux43aux430-ux43eux441ux430ux436ux434ux435ux43dux438ux44f-ux442ux44fux436ux451ux43bux44bux445-ux43aux43eux43cux43fux43eux43dux435ux43dux442ux43eux432-ux43dux435ux444ux442ux438-ux43fux430ux440ux430ux444ux438ux43dux44b-ux430ux441ux444ux430ux43bux44cux442ux435ux43dux44b-ux441ux43cux43eux43bux44b}

\subsection{1.1. Физическая природа фазовых переходов в многокомпонентных углеводородных смесях}\label{ux444ux438ux437ux438ux447ux435ux441ux43aux430ux44f-ux43fux440ux438ux440ux43eux434ux430-ux444ux430ux437ux43eux432ux44bux445-ux43fux435ux440ux435ux445ux43eux434ux43eux432-ux432-ux43cux43dux43eux433ux43eux43aux43eux43cux43fux43eux43dux435ux43dux442ux43dux44bux445-ux443ux433ux43bux435ux432ux43eux434ux43eux440ux43eux434ux43dux44bux445-ux441ux43cux435ux441ux44fux445}

Нефть представляет собой многокомпонентную смесь углеводородов различных гомологических рядов и гетероатомных соединений, находящуюся в пластовых условиях в состоянии, близком к термодинамическому равновесию. При снижении температуры и/или давления в ходе неизотермической фильтрации к скважине или при движении по трубопроводу происходят два принципиально различных типа фазовых переходов:

\begin{itemize}
\tightlist
\item
  \textbf{твёрдое--жидкое (SLE, solid--liquid equilibrium)} --- кристаллизация высокомолекулярных н-алканов (парафинов) с образованием твёрдой воскообразной фазы при охлаждении ниже температуры начала кристаллизации парафинов (WAT, Wax Appearance Temperature);
\item
  \textbf{жидкость--жидкое (LLE) / коллоидное разделение фаз} --- дестабилизация асфальтеновых нано- и микроагрегатов, стабилизированных смолами в объёме нефти как лиофильной коллоидной системы, и их последующая флокуляция и осаждение при изменении давления, температуры или состава (закачка растворителя, разгазирование).
\end{itemize}

Оба процесса управляются условием равенства химических потенциалов (летучестей) компонента \(i\) в сосуществующих фазах. Термодинамическое согласование этих переходов с детальным химическим составом нефти --- ключевая методологическая задача, решаемая через явное описание гомологических рядов и построение псевдокомпонентов.

\subsection{1.2. Химический состав нефти: генерализованные формулы гомологических рядов}\label{ux445ux438ux43cux438ux447ux435ux441ux43aux438ux439-ux441ux43eux441ux442ux430ux432-ux43dux435ux444ux442ux438-ux433ux435ux43dux435ux440ux430ux43bux438ux437ux43eux432ux430ux43dux43dux44bux435-ux444ux43eux440ux43cux443ux43bux44b-ux433ux43eux43cux43eux43bux43eux433ux438ux447ux435ux441ux43aux438ux445-ux440ux44fux434ux43eux432}

Нефть состоит из углеводородов нескольких структурных классов, каждый из которых описывается общей брутто-формулой в зависимости от числа атомов углерода \(n\) и степени цикличности/ненасыщенности:

\begin{itemize}
\tightlist
\item
  \textbf{Парафины (алканы)}, \(\mathrm{C}_n\mathrm{H}_{2n+2}\) --- линейные и разветвлённые насыщенные углеводороды; именно нормальные (н-) парафины с \(n \gtrsim 18\)--20 ответственны за образование воскового осадка при охлаждении;
\item
  \textbf{Нафтены (циклоалканы)}, \(\mathrm{C}_n\mathrm{H}_{2n}\) для моноциклических структур, с понижением водородного числа на 2 за каждое дополнительное кольцо (\(\mathrm{C}_n\mathrm{H}_{2n-2}\), \(\mathrm{C}_n\mathrm{H}_{2n-4}\), \ldots{} для би-, трициклических нафтенов);
\item
  \textbf{Ароматические углеводороды} --- моноциклические (алкилбензолы) \(\mathrm{C}_n\mathrm{H}_{2n-6}\), далее конденсированные полициклические ароматические ряды \(\mathrm{C}_n\mathrm{H}_{2n-x}\) с ростом \(x\) по мере увеличения числа сопряжённых ароматических колец;
\item
  \textbf{Смолы и асфальтены} --- не описываются простой брутто-формулой гомологического ряда, а представляют собой полидисперсные конденсированные полициклические ароматические структуры («средние молекулы» по модели Йена--Маллинса) с алифатическими заместителями и гетероатомами (S, N, O), а также следовыми металлопорфириновыми комплексами Ni и V.
\end{itemize}

Такая PNA-классификация (Paraffins--Naphthenes--Aromatics; в расширенном варианте PIONA добавляет изопарафины и олефины) лежит в основе структурно-группового анализа нефтяных фракций и позволяет описывать неопределённую (плюс-)фракцию нефти конечным набором псевдокомпонентов с заданными молекулярной массой, плотностью и критическими свойствами, оцениваемыми по корреляциям Риази--Даубера, связывающим среднюю температуру кипения и относительную плотность фракции с молекулярной массой и критическими параметрами (Riazi, Daubert, 1987).

\subsection{1.3. SARA- и SCN-характеризация как база построения псевдокомпонентов}\label{sara--ux438-scn-ux445ux430ux440ux430ux43aux442ux435ux440ux438ux437ux430ux446ux438ux44f-ux43aux430ux43a-ux431ux430ux437ux430-ux43fux43eux441ux442ux440ux43eux435ux43dux438ux44f-ux43fux441ux435ux432ux434ux43eux43aux43eux43cux43fux43eux43dux435ux43dux442ux43eux432}

Прямой покомпонентный анализ тяжёлой части нефти (фракция \(C_{7+}\) и выше) невозможен из-за огромного числа изомеров, поэтому на практике применяются два взаимодополняющих подхода характеризации.

\textbf{SARA-фракционирование} (Saturates--Aromatics--Resins--Asphaltenes) разделяет нефть по растворимости/полярности на насыщенные углеводороды, ароматику, смолы и асфальтены (нерастворимый в н-алкане, растворимый в толуоле остаток). Это разделение задаёт макроструктуру состава, необходимую для последующего построения термодинамической модели: насыщенная фракция поставляет н-парафины, кристаллизующиеся при охлаждении, а смолисто-асфальтеновая часть определяет коллоидную стабильность нефти как раствора.

\textbf{Метод одиночных углеродных чисел (Single Carbon Number, SCN)}, предложенный Кацем и Фирузабади, группирует компоненты плюс-фракции конденсатных и нефтяных систем по числу атомов углерода в псевдокомпоненты \(C_7, C_8, \dots, C_{n+}\), каждому из которых приписываются усреднённые молекулярная масса, плотность и температура кипения по данным дистилляции, что позволяет использовать их непосредственно как компоненты в уравнении состояния при расчёте PVT-поведения смеси (Katz, Firoozabadi, 1978). Уитсон формализовал процедуру разбиения плюс-фракции на SCN-группы и последующего их укрупнения (lumping) в ограниченное число псевдокомпонентов для практических расчётов (Whitson, 1983).

Комбинация SARA (макросостав по полярности) и SCN/PNA (молекулярно-массовое и структурно-групповое распределение внутри каждой фракции) даёт полный набор псевдокомпонентов \(\{i\}\) с приписанными им долями \(z_i\), молярными массами \(M_i\), нормальными температурами кипения \(T_{b,i}\) и, через корреляции Риази--Даубера, критическими параметрами \(T_{c,i}, P_{c,i}, \omega_i\), необходимыми для дальнейших термодинамических расчётов равновесия.

\subsection{1.4. Условия фазового равновесия: химический потенциал и летучесть}\label{ux443ux441ux43bux43eux432ux438ux44f-ux444ux430ux437ux43eux432ux43eux433ux43e-ux440ux430ux432ux43dux43eux432ux435ux441ux438ux44f-ux445ux438ux43cux438ux447ux435ux441ux43aux438ux439-ux43fux43eux442ux435ux43dux446ux438ux430ux43b-ux438-ux43bux435ux442ux443ux447ux435ux441ux442ux44c}

Для системы, находящейся в термодинамическом равновесии при постоянных \(T\) и \(P\), необходимым и достаточным условием сосуществования фаз \(\alpha\) (например, жидкой) и \(\beta\) (твёрдой или второй жидкой) является равенство химических потенциалов каждого компонента \(i\) во всех фазах:

\[
\mu_i^{\alpha}(T,P,\mathbf{x}^\alpha) = \mu_i^{\beta}(T,P,\mathbf{x}^\beta), \qquad i = 1,\dots,N.
\]

Вводя летучесть \(f_i\) через соотношение \(\mu_i = \mu_i^{\circ} + RT\ln(f_i/f_i^{\circ})\), условие равновесия сводится к равенству летучестей:

\[
f_i^{\alpha}(T,P,\mathbf{x}^\alpha) = f_i^{\beta}(T,P,\mathbf{x}^\beta).
\]

Для твёрдо-жидкого равновесия н-парафинов летучесть компонента \(i\) в твёрдой фазе обычно выражают через летучесть чистого переохлаждённого твёрдого вещества \(f_i^{S,\text{pure}}\) и коэффициент активности \(\gamma_i^S\) в твёрдом растворе (или активность \(a_i^S = \gamma_i^S x_i^S\) для случая многофазного твёрдого раствора), а в жидкой --- через летучесть чистой переохлаждённой жидкости и коэффициент активности \(\gamma_i^L\):

\[
x_i^{S}\,\gamma_i^{S}\,f_i^{S,\text{pure}}(T,P) = x_i^{L}\,\gamma_i^{L}\,f_i^{L,\text{pure}}(T,P).
\]

Отношение летучестей чистого переохлаждённого твёрдого вещества к чистой переохлаждённой жидкости при температуре ниже точки плавления компонента \(T_{m,i}\) выражается через термодинамический цикл плавления и обычно записывается в виде (с учётом энтальпии плавления \(\Delta H_{m,i}\), температуры плавления \(T_{m,i}\), теплоты и температуры твердо-твёрдого перехода \(\Delta H_{tr,i}\), \(T_{tr,i}\), а также разности теплоёмкостей \(\Delta C_{p,i}\) между жидкой и твёрдой фазами):

\[
\ln\frac{f_i^{S,\text{pure}}}{f_i^{L,\text{pure}}} = -\frac{\Delta H_{m,i}}{RT}\left(1-\frac{T}{T_{m,i}}\right) - \frac{\Delta H_{tr,i}}{RT}\left(1-\frac{T}{T_{tr,i}}\right) + \frac{\Delta C_{p,i}}{R}\left(\ln\frac{T_{m,i}}{T} - 1 + \frac{T_{m,i}}{T}\right).
\]

Именно такое соотношение (или его модификации) лежит в основе практически всех термодинамических моделей осаждения парафинов, различающихся между собой в основном способом описания коэффициентов активности \(\gamma_i^{S}, \gamma_i^{L}\), то есть неидеальности твёрдого и жидкого растворов.

\subsection{1.5. Multi-solid, solid-solution и полимерные модели растворов}\label{multi-solid-solid-solution-ux438-ux43fux43eux43bux438ux43cux435ux440ux43dux44bux435-ux43cux43eux434ux435ux43bux438-ux440ux430ux441ux442ux432ux43eux440ux43eux432}

\subsubsection{1.5.1. Multi-solid модель}\label{multi-solid-ux43cux43eux434ux435ux43bux44c}

В подходе Лира-Галеаны, Фирузабади и Праусница каждый твёрдый н-парафин рассматривается как отдельная чистая (или квазичистая) твёрдая фаза, не образующая единого раствора с другими твёрдыми парафинами: экспериментально подтверждённое допущение состоит в том, что осаждённый воск состоит из нескольких твёрдых фаз, каждая из которых представлена чистым компонентом или псевдокомпонентом, не смешивающимся с другими твёрдыми фазами (Lira-Galeana, Firoozabadi, Prausnitz, 1996). При этом жидкофазные свойства (летучести) рассчитываются из уравнения состояния, а не из модели активности твёрдого раствора. Это допущение резко упрощает вычисления (активность твёрдой фазы \(a_i^S = 1\) для присутствующих твёрдых компонентов), и модель хорошо согласуется с экспериментальными данными по количеству выпавшего воска для бинарных и многокомпонентных смесей, включая реальную нефть (Lira-Galeana, Firoozabadi, Prausnitz, 1996). Multi-solid подход физически оправдан тем, что молекулярные размеры и кристаллические структуры разных н-алканов существенно различаются, что затрудняет образование единого твёрдого раствора --- отсюда экспериментально наблюдаемое расслоение осадка на несколько твёрдых фаз с разной температурой плавления.

\subsubsection{1.5.2. Solid-solution (модель твёрдого раствора)}\label{solid-solution-ux43cux43eux434ux435ux43bux44c-ux442ux432ux451ux440ux434ux43eux433ux43e-ux440ux430ux441ux442ux432ux43eux440ux430}

Альтернативный и исторически более ранний подход, развитый Воном, трактует осаждённый воск как единый неидеальный твёрдый раствор, находящийся в равновесии с жидкой (и, при необходимости, паровой) фазой; развита теория твёрдо-жидко-парового равновесия и применена к расчёту образования восковой фазы из тяжёлых углеводородных смесей (Won, 1986). Неидеальность как твёрдого, так и жидкого растворов описывается регулярной теорией раствора (regular solution theory) через параметры растворимости \(\delta_i\):

\[
\ln \gamma_i = \frac{V_i}{RT}\left(\delta_i - \bar\delta\right)^2, \qquad \bar\delta = \sum_j \phi_j \delta_j,
\]

где \(V_i\) --- молярный объём компонента, \(\phi_j\) --- объёмная доля компонента \(j\) в соответствующей фазе. Хансен с соавторами и впоследствии Педерсен усовершенствовали этот подход, применив его на практике для прогноза точки помутнения (cloud point) и количества осаждённого воска в реальных нефтях с использованием SCN-характеризации плюс-фракции (Hansen, Fredenslund, Pedersen, Rønningsen, 1988; Pedersen, 1995) непосредственно по результатам газохроматографического анализа нефти. Пан и др. дополнительно исследовали влияние давления и состава на массу выпавшего воска, сопоставив экспериментальные данные с расчётами по модели твёрдого раствора (Pan, Firoozabadi, Fotland, 1997). Нитита с соавторами распространили аналогичный формализм на газоконденсатные системы, где важна одновременная работа с паровой фазой (Nichita, Goual, Firoozabadi, 2001).

Ключевое отличие solid-solution от multi-solid модели --- в предположении о взаимной растворимости твёрдых парафинов, что для узких по составу восков (близкие цепи \(n\)) физически обоснованно, но становится менее точным для широких по молекулярно-массовому распределению смесей, где наблюдается сегрегация на несколько твёрдых фаз.

\subsubsection{1.5.3. Локально-композиционные модели (Вильсон) для твёрдого раствора}\label{ux43bux43eux43aux430ux43bux44cux43dux43e-ux43aux43eux43cux43fux43eux437ux438ux446ux438ux43eux43dux43dux44bux435-ux43cux43eux434ux435ux43bux438-ux432ux438ux43bux44cux441ux43eux43d-ux434ux43bux44f-ux442ux432ux451ux440ux434ux43eux433ux43e-ux440ux430ux441ux442ux432ux43eux440ux430}

Кутиньо и Стенби развили более строгий подход к описанию неидеальности твёрдого раствора н-алканов, применив локально-композиционное уравнение Вильсона (изначально разработанное для жидких растворов) к описанию энергетики упаковки молекул в кристаллической решётке твёрдого воска --- предложены предиктивные модели локального состава для твёрдо-жидкого равновесия в системах н-алканов, основанные на уравнении Вильсона для многокомпонентных систем (Coutinho, Stenby, 1996). Такой подход учитывает энергетическое различие взаимодействий между соседними цепями разной длины в кристаллической решётке и позволяет предсказывать не только температуру начала кристаллизации, но и относительное распределение количества выпавшего парафина по углеродным номерам без подгоночных параметров, оценённых по каждой конкретной нефти.

\subsubsection{1.5.4. Полимерные модели растворов (Флори--Хаггинса) и коллоидная термодинамика асфальтенов}\label{ux43fux43eux43bux438ux43cux435ux440ux43dux44bux435-ux43cux43eux434ux435ux43bux438-ux440ux430ux441ux442ux432ux43eux440ux43eux432-ux444ux43bux43eux440ux438ux445ux430ux433ux433ux438ux43dux441ux430-ux438-ux43aux43eux43bux43bux43eux438ux434ux43dux430ux44f-ux442ux435ux440ux43cux43eux434ux438ux43dux430ux43cux438ux43aux430-ux430ux441ux444ux430ux43bux44cux442ux435ux43dux43eux432}

Для смол и асфальтенов прямое применение теории регулярных растворов Гильдебранда некорректно из-за огромного различия молярных объёмов дисперсной (полиароматические кластеры) и дисперсионной (лёгкие и средние алканы) сред. Поэтому асфальтеновая термодинамика развивалась по двум направлениям.

\textbf{Молекулярно-термодинамические (полимерно-растворные) модели.} Ву, Праусниц и Фирузабади разработали молекулярно-термодинамическую структуру для равновесия асфальтен--нефть (Wu, Prausnitz, Firoozabadi, 1998), в которой асфальтен рассматривается как полидисперсный полимероподобный растворённый компонент в смеси растворителей (мальтенах), а свободная энергия смешения записывается в духе теории Флори--Хаггинса, дополненной вкладом от ассоциации (водородные связи, π--π стекинг ароматических плоскостей). Впоследствии эта модель была дополнена явным учётом равновесия ассоциации асфальтеновых молекул при расчёте осаждения асфальтенов в пластовых флюидах (Wu, Prausnitz, Firoozabadi, 2000). Свободная энергия Гиббса смешения по Флори--Хаггинсу для полимерного (асфальтенового) компонента записывается как

\[
\frac{\Delta G_{mix}}{RT} = \sum_i n_i \ln\phi_i + \chi\, n_1 \phi_2,
\]

где \(\phi_i\) --- объёмные доли, \(n_i\) --- число молей, \(\chi\) --- параметр Флори--Хаггинса взаимодействия между асфальтеном (индекс 2) и растворителем (индекс 1), связанный с параметрами растворимости \(\delta_1, \delta_2\) соотношением \(\chi = V_1(\delta_1-\delta_2)^2/RT\). Отсюда следует, что осаждение асфальтенов термодинамически провоцируется падением параметра растворимости смеси \(\delta_1\) (например, при добавлении лёгкого н-алкана-осадителя или при разгазировании нефти со снижением давления ниже давления насыщения, когда лёгкие компоненты покидают жидкую фазу).

\textbf{Коллоидно-термодинамические модели.} Леонтаритис и Мансури предложили альтернативную концепцию, где асфальтен в нефти существует в виде коллоидных мицелл, стабилизированных адсорбированными на их поверхности молекулами смол; выпадение асфальтена интерпретируется как флокуляция коллоида при десорбции стабилизирующих смол с поверхности мицелл вследствие изменения термобарических условий или состава --- предложена термодинамико-коллоидная модель флокуляции асфальтенов при добыче и переработке нефти (Leontaritis, Mansoori, 1987). Викторов и Фирузабади формализовали эту идею в виде строгой модели мицеллизации, вводя равновесие агрегации мономер ⇌ мицелла и рассчитывая распределение агрегатов по размерам через баланс энергии агрегации и энтропии смешения (Victorov, Firoozabadi, 1996). Современное развитие этой линии --- модель агрегационной термодинамики Ванга и др., где явно вводится ступенчатое равновесие нанооагрегации и кластеризации асфальтеновых молекул (в духе иерархической структуры Йена--Маллинса: молекула → наноагрегат → кластер), количественно связывающая термодинамику агрегации с наблюдаемым распределением размеров агрегатов (Wang, Hao, Islam, Chen, 2016). Это согласуется с результатами прямых спектроскопических наблюдений кластеров асфальтеновых наноагрегатов в пластовых условиях по данным даунхол-градиента состава (Mullins, Seifert, Zuo, Zeybek, 2012).

Буэнростро-Гонсалес с соавторами систематически сопоставили теоретические (термодинамико-коллоидные) и экспериментальные аспекты осаждения асфальтенов при добавлении н-алкана-осадителя, подтвердив применимость модели равновесия жидкость--жидкое (с асфальтеновой фазой как отдельной жидкой/твёрдой псевдофазой) к количественному прогнозу порога осаждения (Buenrostro-González, Lira-Galeana et al., 2004).

\subsection{1.6. Уравнения состояния: кубические EOS и PC-SAFT для детального компонентного состава}\label{ux443ux440ux430ux432ux43dux435ux43dux438ux44f-ux441ux43eux441ux442ux43eux44fux43dux438ux44f-ux43aux443ux431ux438ux447ux435ux441ux43aux438ux435-eos-ux438-pc-saft-ux434ux43bux44f-ux434ux435ux442ux430ux43bux44cux43dux43eux433ux43e-ux43aux43eux43cux43fux43eux43dux435ux43dux442ux43dux43eux433ux43e-ux441ux43eux441ux442ux430ux432ux430}

Жидкофазные (а для газоконденсатных систем --- и паровые) летучести компонентов необходимо рассчитывать из уравнения состояния (EOS), согласованного с принятой SARA/SCN-характеризацией.

\textbf{Кубические EOS} (Пенг--Робинсон, Соаве--Редлиха--Квонга) остаются рабочим инструментом в большинстве multi-solid и solid-solution реализаций благодаря вычислительной эффективности и совместимости с существующими PVT-моделями пластовых флюидов; они требуют лишь критических параметров псевдокомпонентов (\(T_c, P_c, \omega\)), получаемых из корреляций Риази--Даубера по молекулярной массе и плотности SCN-фракции. Общая методология расчёта фазовых равновесий флюидов через кубические уравнения состояния рассмотрена, в частности, у Ву, Праусница и Фирузабади (Wu, Prausnitz, Firoozabadi --- обзор EOS, 2000).

\textbf{PC-SAFT (Perturbed-Chain Statistical Associating Fluid Theory)}, развитое Гроссом и Садовски на основе исходной SAFT-теории Чапмена и др., описывает молекулы как цепочки сегментов с явным учётом формы (числа сегментов \(m\), диаметра сегмента \(\sigma\)), дисперсионных взаимодействий и, при необходимости, ассоциации (водородные связи) --- предложено уравнение состояния Perturbed-Chain SAFT, основанное на теории возмущений для цепных молекул (Gross, Sadowski, 2001), на базе исходной модели раствора SAFT для ассоциирующих флюидов (Chapman, Gubbins, Jackson, Radosz, 1989). Для тяжёлых нефтяных компонентов, особенно асфальтенов, преимущество PC-SAFT состоит в физически прозрачном разделении вкладов твёрдой сферы, цепочечности, дисперсии и ассоциации в свободную энергию, что позволяет напрямую связать параметры EOS с молекулярной структурой псевдокомпонента (числом ароматических колец, длиной алкильных цепей), оценённой из SARA/PNA-анализа, а не только с эмпирическими корреляциями критических свойств. Это делает PC-SAFT предпочтительным инструментом там, где детальная химическая структура компонента (в частности --- степень ароматичности и способность к самоассоциации асфальтенов через π-стекинг и водородные связи гетероатомных групп) существенна для термодинамики равновесия.

\subsection{1.7. Прогнозирование температуры начала кристаллизации парафинов (WAT)}\label{ux43fux440ux43eux433ux43dux43eux437ux438ux440ux43eux432ux430ux43dux438ux435-ux442ux435ux43cux43fux435ux440ux430ux442ux443ux440ux44b-ux43dux430ux447ux430ux43bux430-ux43aux440ux438ux441ux442ux430ux43bux43bux438ux437ux430ux446ux438ux438-ux43fux430ux440ux430ux444ux438ux43dux43eux432-wat}

Температура появления воска (WAT) определяется как температура, при которой мольная доля выпавшего твёрдого парафина становится (формально) отличной от нуля, то есть точка бифуркации кривой \(w_S(T)\) количества твёрдой фазы от температуры. Во всех рассмотренных термодинамических моделях (multi-solid, solid-solution, локально-композиционные) WAT находится как температура, при которой суммарное условие равновесия

\[
\sum_i x_i^L \gamma_i^L \exp\!\left[-\frac{\Delta H_{m,i}}{RT}\Big(1-\frac{T}{T_{m,i}}\Big)+\dots\right] = 1
\]

впервые выполняется при понижении температуры от начального однофазного (жидкого) состояния --- то есть когда сумма «эффективных активностей» компонентов в предположении их полного растворения в жидкости достигает единицы. Точность прогноза WAT сильно зависит от качества характеризации хвостовой (высокомолекулярной) части распределения н-парафинов по SCN, поскольку именно самые тяжёлые компоненты первыми достигают насыщения. Экспериментальная валидация подобных моделей против измерений WAT, точки помутнения и количества выпавшего воска подтверждена, в частности, для многокомпонентных модельных смесей и реальных нефтей у Хансена и др. (1988) и Педерсен (1995), а для газоконденсатных систем --- у Нититы и др. (2001).

\subsection{1.8. Резюме раздела}\label{ux440ux435ux437ux44eux43cux435-ux440ux430ux437ux434ux435ux43bux430}

Термодинамическое описание осаждения тяжёлых компонентов нефти строится по единой логической схеме: (1) детальная химическая характеризация нефти через SARA-фракционирование и SCN/PNA-анализ, дающая набор псевдокомпонентов с явной привязкой к гомологическим рядам (\(\mathrm{C}_n\mathrm{H}_{2n+2}\) для парафинов и т.д.); (2) выбор модели неидеальности сосуществующих фаз --- multi-solid (Лира-Галеана и др.), solid-solution (Вон; Хансен и др.; Педерсен) или локально-композиционная (Кутиньо--Стенби) для парафинов, полимерно-растворная Флори--Хаггинс-подобная (Ву--Праусниц--Фирузабади) или коллоидно-мицеллярная (Леонтаритис--Мансури; Викторов--Фирузабади; Ван и др.) для асфальтенов; (3) расчёт летучестей через кубическое EOS или PC-SAFT (Гросс--Садовски на основе SAFT Чапмена и др.), согласованный с принятой характеризацией. Этот комплекс моделей формирует термодинамическую подсистему для замыкания уравнений неизотермической многофазной фильтрации в пласте, обсуждаемых далее с кинетической точки зрения в Разделе 2.

\newpage

\section{Раздел 2. Кинетика и механизмы осаждения тяжёлых компонентов нефти}\label{ux440ux430ux437ux434ux435ux43b-2.-ux43aux438ux43dux435ux442ux438ux43aux430-ux438-ux43cux435ux445ux430ux43dux438ux437ux43cux44b-ux43eux441ux430ux436ux434ux435ux43dux438ux44f-ux442ux44fux436ux451ux43bux44bux445-ux43aux43eux43cux43fux43eux43dux435ux43dux442ux43eux432-ux43dux435ux444ux442ux438}

\subsection{2.1. От термодинамики к кинетике: постановка задачи}\label{ux43eux442-ux442ux435ux440ux43cux43eux434ux438ux43dux430ux43cux438ux43aux438-ux43a-ux43aux438ux43dux435ux442ux438ux43aux435-ux43fux43eux441ux442ux430ux43dux43eux432ux43aux430-ux437ux430ux434ux430ux447ux438}

Термодинамическое равновесие (Раздел 1) определяет лишь предельное, асимптотическое состояние системы --- какая доля тяжёлого компонента \emph{может} выпасть в твёрдую/дисперсную фазу при заданных \(T, P, z_i\). Однако в условиях неизотермической фильтрации в пласте, в стволе скважины или в трубопроводе система не находится в состоянии локального равновесия по всему объёму: у стенки поры (или трубы), где температура ниже, локальная концентрация растворённого парафина или устойчивость асфальтеновых агрегатов ниже равновесного значения в объёме потока, что создаёт движущую силу для переноса вещества к стенке и последующего роста отложений. Кинетическая часть модели должна описать: (i) механизм переноса растворённых/диспергированных тяжёлых частиц к поверхности отложения; (ii) кинетику собственно кристаллизации/агрегации (не мгновенную, в отличие от предположения локального равновесия); (iii) рост и последующую эволюцию (старение) слоя отложений.

\subsection{2.2. Физика транспортных механизмов переноса к стенке}\label{ux444ux438ux437ux438ux43aux430-ux442ux440ux430ux43dux441ux43fux43eux440ux442ux43dux44bux445-ux43cux435ux445ux430ux43dux438ux437ux43cux43eux432-ux43fux435ux440ux435ux43dux43eux441ux430-ux43a-ux441ux442ux435ux43dux43aux435}

Экспериментальные исследования отложения парафина в трубопроводах (в первую очередь классическая работа по трансаляскинскому трубопроводу) выделили несколько конкурирующих механизмов латерального переноса растворённых и взвешенных частиц воска к стенке: молекулярную диффузию, сдвиговую дисперсию и броуновскую диффузию --- предполагается, что осаждение происходит в результате латерального переноса за счёт диффузии, сдвиговой дисперсии и броуновской диффузии (Burger, Perkins, Striegler, 1981). Позднее этот список был расширен добавлением гравитационного оседания и сдвигового сдирания (shear stripping) отложений набегающим потоком. Ниже рассмотрен физический смысл и математическая формулировка каждого механизма.

\subsubsection{2.2.1. Молекулярная диффузия (закон Фика)}\label{ux43cux43eux43bux435ux43aux443ux43bux44fux440ux43dux430ux44f-ux434ux438ux444ux444ux443ux437ux438ux44f-ux437ux430ux43aux43eux43d-ux444ux438ux43aux430}

Молекулярная диффузия признаётся практически всеми исследователями доминирующим механизмом переноса растворённых парафинов к охлаждённой стенке в условиях ламинарного и турбулентного однофазного потока: молекулярная диффузия является основной причиной осаждения молекул воска и главной теоретической опорой для построения моделей осаждения воска. Физический механизм: при радиальном градиенте температуры \(dT/dr\) в потоке возникает радиальный градиент равновесной растворимости парафина в жидкой фазе \(dC_{eq}/dr = (dC_{eq}/dT)(dT/dr)\); поскольку растворимость н-парафинов резко падает с понижением температуры, у стенки жидкость оказывается локально пересыщенной относительно равновесной концентрации при данной температуре стенки, и растворённый парафин диффундирует к стенке по закону Фика:

\[
J_i = -D_i \frac{\partial C_i}{\partial r}\bigg|_{r=R} \approx -D_i \frac{dC_{eq}}{dT}\frac{dT}{dr}\bigg|_{r=R},
\]

где \(J_i\) --- диффузионный поток массы к стенке (кг/(м²·с)), \(D_i\) --- коэффициент молекулярной диффузии парафина в нефти, \(R\) --- радиус трубы/поры. Массовая скорость роста отложения на единицу площади стенки в квазистационарном приближении принимается равной этому диффузионному потоку. Такая формулировка (диффузионная модель Фика, управляемая температурным градиентом через производную кривой растворимости) лежит в основе большинства инженерных моделей отложения парафина, включая работу Уайнгартена и Ойхнера, посвящённую методам прогноза осаждения и отложения воска (Weingarten, Euchner, 1988), и экспериментально-теоретическое исследование кинетики отложения парафина Брауна, Нисена и Эриксона: измерение и прогноз кинетики отложения парафина (Brown, Niesen, Erickson, 1993).

\subsubsection{2.2.2. Сдвиговая дисперсия (shear dispersion)}\label{ux441ux434ux432ux438ux433ux43eux432ux430ux44f-ux434ux438ux441ux43fux435ux440ux441ux438ux44f-shear-dispersion}

Второй механизм --- перенос уже выпавших в объёме потока мелких кристаллов твёрдого парафина к стенке под действием сдвигового течения (аналог механизма Тейлора-дисперсии применительно к твёрдым частицам в сдвиговом поле). Наличие радиального градиента скорости сдвига \(\dot\gamma(r)\) приводит к латеральной миграции взвешенных частиц из области высокого сдвига (центр трубы) в область низкого сдвига (пристеночная зона) --- эффект, количественно описываемый через сдвиговый коэффициент диффузии \(D_{shear} \propto \dot\gamma\, d_p^2\), где \(d_p\) --- характерный размер частицы. Бюргер, Перкинс и Стриглер экспериментально показали значимость этого механизма наряду с молекулярной и броуновской диффузией применительно к отложению воска в Трансаляскинском трубопроводе (Burger, Perkins, Striegler, 1981). Последующие исследования (включая обзоры механизмов отложения) в целом сходятся на том, что для однофазного ламинарного и умеренно турбулентного потока в прямых трубах доминирует именно молекулярная диффузия, тогда как сдвиговая дисперсия становится значимой при повышенных концентрациях твёрдой фазы и в изогнутых участках/при высоких скоростях сдвига.

\subsubsection{2.2.3. Броуновское движение}\label{ux431ux440ux43eux443ux43dux43eux432ux441ux43aux43eux435-ux434ux432ux438ux436ux435ux43dux438ux435}

Броуновская диффузия описывает случайное блуждание малых (субмикронных) кристаллов парафина или частиц асфальтена под действием теплового движения молекул растворителя, приводящее к эффективному диффузионному потоку с коэффициентом Стокса--Эйнштейна

\[
D_B = \frac{k_B T}{6\pi \mu r_p},
\]

где \(k_B\) --- постоянная Больцмана, \(\mu\) --- динамическая вязкость нефти, \(r_p\) --- радиус частицы. Для типичных размеров частиц воска (\(r_p \sim 1\)--10 мкм) и вязкости нефти этот механизм даёт существенно меньший поток массы по сравнению с молекулярной диффузией растворённых молекул парафина, поэтому в большинстве инженерных моделей вклад броуновской диффузии на отложение воска в прямых трубах признаётся пренебрежимо малым: влияние броуновского движения и гравитационного оседания на отложение воска в целом может не учитываться. Для более крупных и плотных асфальтеновых агрегатов и флоккул, однако, броуновская диффузия и коагуляция агрегатов остаются существенным механизмом, определяющим кинетику роста частиц перед их осаждением на стенке.

\subsubsection{2.2.4. Гравитационное оседание}\label{ux433ux440ux430ux432ux438ux442ux430ux446ux438ux43eux43dux43dux43eux435-ux43eux441ux435ux434ux430ux43dux438ux435}

Гравитационное оседание описывает вертикальное движение твёрдых частиц или тяжёлой (асфальтеновой/смолистой) дисперсной фазы под действием силы тяжести за вычетом плавучести, уравновешиваемой силой вязкого сопротивления (закон Стокса для малых чисел Рейнольдса частицы):

\[
u_{settle} = \frac{2 r_p^2 (\rho_p - \rho_f) g}{9\mu},
\]

где \(\rho_p, \rho_f\) --- плотности частицы и флюида. Для отложения парафина в горизонтальных трубах этот механизм, как правило, второстепенен по сравнению с диффузионным переносом в прямых участках трубы, что было экспериментально подтверждено переключением ориентации испытательного участка трубы в исследовании Бюргера с соавторами. Для асфальтеновых частиц ситуация иная: в вертикальных и наклонных стволах скважин, а также при формировании достаточно крупных флоккул, гравитационное оседание (наряду с термофорезом и силой сопротивления) может играть значимую роль в переносе частиц к стенке, что отражено в комплексных эйлеровых моделях переноса асфальтенов, объединяющих молекулярную и турбулентную диффузию, турбофорез, термофорез и гравитационное оседание.

\subsubsection{2.2.5. Сдвиговое сдирание (shear stripping) и суммарный баланс массы у стенки}\label{ux441ux434ux432ux438ux433ux43eux432ux43eux435-ux441ux434ux438ux440ux430ux43dux438ux435-shear-stripping-ux438-ux441ux443ux43cux43cux430ux440ux43dux44bux439-ux431ux430ux43bux430ux43dux441-ux43cux430ux441ux441ux44b-ux443-ux441ux442ux435ux43dux43aux438}

Помимо механизмов переноса \emph{к} стенке, необходимо учитывать противоположно направленный процесс --- механическое сдирание (эрозию) уже отложившегося, но ещё не «созревшего» мягкого геля потоком нефти. Это эффективно снижает нетто-скорость роста отложения при высоких скоростях потока/напряжениях сдвига на стенке. Общий баланс массы на единицу длины трубы/поровой стенки для толщины отложения \(\delta(t)\) записывается как разность потока осаждения и потока сдирания:

\[
\rho_{dep}\frac{d\delta}{dt} = J_{deposition}(T, \dot\gamma, C) - J_{stripping}(\tau_w, \delta),
\]

где \(\rho_{dep}\) --- плотность (пористого) отложения, \(\tau_w\) --- напряжение сдвига на стенке. Сочетание конкурирующих механизмов переноса и сдирания объясняет экспериментально наблюдаемый максимум толщины отложения при некоторой промежуточной скорости потока --- при малых скоростях слаб конвективный подвод тепла и вещества, при высоких --- превалирует механическое сдирание.

\subsection{2.3. Кинетические модели роста слоя отложений}\label{ux43aux438ux43dux435ux442ux438ux447ux435ux441ux43aux438ux435-ux43cux43eux434ux435ux43bux438-ux440ux43eux441ux442ux430-ux441ux43bux43eux44f-ux43eux442ux43bux43eux436ux435ux43dux438ux439}

\subsubsection{2.3.1. Модель мгновенного (равновесного) осаждения vs.~кинетическая модель}\label{ux43cux43eux434ux435ux43bux44c-ux43cux433ux43dux43eux432ux435ux43dux43dux43eux433ux43e-ux440ux430ux432ux43dux43eux432ux435ux441ux43dux43eux433ux43e-ux43eux441ux430ux436ux434ux435ux43dux438ux44f-vs.-ux43aux438ux43dux435ux442ux438ux447ux435ux441ux43aux430ux44f-ux43cux43eux434ux435ux43bux44c}

Историческая эволюция моделей отложения воска в трубах прошла путь от предположения о мгновенном достижении термодинамического равновесия по всей толщине отложения («solubility model», где вся разность концентраций между объёмом потока и линией равновесия при температуре стенки немедленно кристаллизуется на стенке) к модели, явно учитывающей конечную скорость кристаллизации парафина в объёме жидкости («kinetic model»). Хуанг, Ли, Сенра и Фоглер сформулировали фундаментальную модель, показавшую, что оба предельных подхода (независимая модель теплообмена/массообмена без учёта кристаллизации, и модель мгновенной растворимости) соответственно недооценивают либо переоценивают толщину отложения, тогда как учёт кинетики кристаллизации воска в потоке нефти как конкурирующего с отложением на стенке процесса даёт точные предсказания для лабораторных и пилотных экспериментов на трёх различных нефтях (Huang, Lee, Senra, Fogler, 2011). Формально кинетика кристаллизации записывается как уравнение релаксации к локальному равновесию:

\[
\frac{dC_{S}}{dt} = k_{cryst}\left(C_{eq}(T) - C_L\right)_+,
\]

где \(C_S, C_L\) --- концентрации выпавшего твёрдого и растворённого парафина, \(k_{cryst}\) --- эффективная константа скорости кристаллизации, \(C_{eq}(T)\) --- равновесная растворимость по термодинамической модели раздела 1.

\subsubsection{2.3.2. Уравнение теплопереноса и сопряжённая тепловая задача}\label{ux443ux440ux430ux432ux43dux435ux43dux438ux435-ux442ux435ux43fux43bux43eux43fux435ux440ux435ux43dux43eux441ux430-ux438-ux441ux43eux43fux440ux44fux436ux451ux43dux43dux430ux44f-ux442ux435ux43fux43bux43eux432ux430ux44f-ux437ux430ux434ux430ux447ux430}

Поскольку интенсивность отложения определяется в первую очередь радиальным градиентом температуры (через производную кривой равновесной растворимости \(dC_{eq}/dT\)), кинетическая модель роста отложения решается совместно с уравнением энергии в потоке (конвективно-кондуктивный теплоперенос в трубе/поре с учётом теплового сопротивления растущего слоя отложения, обладающего низкой теплопроводностью по сравнению с металлом стенки). Рост толщины отложения \(\delta(t)\) в квазистационарном приближении для диффузионно-контролируемого режима записывается как

\[
\frac{d\delta}{dt} = \frac{1}{\rho_{dep}\,C_{dep}}\, D_{eff}\, \frac{dC_{eq}}{dT}\,\frac{\partial T}{\partial r}\bigg|_{r=R(t)},
\]

где \(C_{dep}\) --- массовая доля парафина в самом отложении (определяемая, в свою очередь, эволюционным процессом старения, см. §2.4), а \(D_{eff}\) --- эффективный коэффициент диффузии, учитывающий как молекулярную, так и (при необходимости) сдвиговую составляющие. Такая сопряжённая формулировка положена в основу практически всех современных промысловых моделей прогноза отложения парафина в добывающих скважинах и трубопроводах, начиная с эмпирико-теоретических корреляций Уайнгартена и Ойхнера (Weingarten, Euchner, 1988) и заканчивая современными фундаментальными моделями (Huang, Lee, Senra, Fogler, 2011).

\subsubsection{2.3.3. Кинетика осаждения асфальтенов: агрегация и перенос частиц}\label{ux43aux438ux43dux435ux442ux438ux43aux430-ux43eux441ux430ux436ux434ux435ux43dux438ux44f-ux430ux441ux444ux430ux43bux44cux442ux435ux43dux43eux432-ux430ux433ux440ux435ux433ux430ux446ux438ux44f-ux438-ux43fux435ux440ux435ux43dux43eux441-ux447ux430ux441ux442ux438ux446}

Для асфальтенов кинетическая цепочка «осаждение → перенос → отложение» дополнительно включает стадию агрегации первично выпавших наноагрегатов в более крупные флоккулы, кинетика которой определяется режимом столкновений частиц (диффузионно-лимитируемая агрегация, DLA, или реакционно-лимитируемая агрегация, RLA) и явно входит в комплексные модели транспорта асфальтенов к стенке скважины/трубопровода, где молекулярная диффузия и сдвиговое сдирание рассматриваются как конкурирующие механизмы, определяющие радиальную скорость диффузии и скорость отложения асфальтеновых частиц на поверхности ствола скважины. Радиальная скорость диффузии при этом связывается с радиальным температурным градиентом, индуцирующим радиальный градиент концентрации в жидкой фазе, --- то есть формально аналогично уравнению Фика для парафинов, но с константой скорости, определяемой не кристаллизацией, а кинетикой агрегации частиц (зависящей от степени пересыщения по асфальтену и локальной температуры/давления).

\subsection{2.4. Модели старения (ageing) отложений}\label{ux43cux43eux434ux435ux43bux438-ux441ux442ux430ux440ux435ux43dux438ux44f-ageing-ux43eux442ux43bux43eux436ux435ux43dux438ux439}

\subsubsection{2.4.1. Физическая природа процесса старения}\label{ux444ux438ux437ux438ux447ux435ux441ux43aux430ux44f-ux43fux440ux438ux440ux43eux434ux430-ux43fux440ux43eux446ux435ux441ux441ux430-ux441ux442ux430ux440ux435ux43dux438ux44f}

Отложения парафина, образующиеся на охлаждённой стенке, изначально представляют собой рыхлый гель --- трёхмерную кристаллическую сетку н-алканов, в порах которой захвачено значительное количество жидкой нефти (мальтеновой фазы вместе с растворёнными более лёгкими парафинами). Сингх, Венкатесан, Фоглер и Нагараджан провели фундаментальное исследование образования и старения такого тонкого начального гель-слоя, показав, что после формирования начального геля молекулы воска продолжают диффундировать в его структуру, увеличивая содержание твёрдой фазы в отложении (Singh, Venkatesan, Fogler, Nagarajan, 2000). Это медленное дозревание --- «старение» --- приводит к двум наблюдаемым эффектам: (i) увеличению массовой доли твёрдого парафина в отложении со временем при почти неизменной внешней толщине отложения; (ii) соответствующему увеличению механической твёрдости (прочности на сдвиг, температуры плавления) отложения, поскольку доля более тяжёлых, высокоплавких компонентов в застывшем геле растёт.

\subsubsection{2.4.2. Механизм встречной диффузии (counter-diffusion)}\label{ux43cux435ux445ux430ux43dux438ux437ux43c-ux432ux441ux442ux440ux435ux447ux43dux43eux439-ux434ux438ux444ux444ux443ux437ux438ux438-counter-diffusion}

Ключевой физический механизм старения, установленный экспериментально Сингхом с соавторами, --- это встречная (counter-diffusion) диффузия компонентов различного углеродного числа: процесс старения представляет собой явление встречной диффузии с критическим углеродным числом, выше которого молекулы воска диффундируют внутрь геля, а ниже которого молекулы масла (более лёгкие компоненты) диффундируют из отложения наружу (Singh, Venkatesan, Fogler, Nagarajan, 2000). Иными словами, гель-отложение выступает как избирательно проницаемая среда: тяжёлые парафины (с температурой плавления выше локальной температуры геля) термодинамически выгодно кристаллизуются внутри пористой структуры и диффундируют туда из окружающей нефти, тогда как относительно лёгкие компоненты (не кристаллизующиеся при данной температуре) вымываются диффузией в обратном направлении в объём потока. Математически это описывается многокомпонентной системой уравнений диффузии Фика для каждого псевдокомпонента \(i\) внутри пористой структуры геля:

\[
\frac{\partial C_i^{gel}}{\partial t} = D_i^{eff}\frac{1}{r}\frac{\partial}{\partial r}\left(r\frac{\partial C_i^{gel}}{\partial r}\right), \qquad i = 1,\dots,N,
\]

с граничным условием на внешней (жидкостной) стороне геля, задаваемым равновесной концентрацией компонента \(i\) в объёмной нефти при локальной температуре, и с зависимостью эффективного коэффициента диффузии \(D_i^{eff}\) от локальной пористости геля (доли захваченной жидкой нефти), которая сама эволюционирует со временем по мере кристаллизации.

Авторы также показали, что скорость старения геля зависит от расхода потока нефти (через теплообмен и, соответственно, температуру стенки) и температуры стенки --- то есть кинетика встречной диффузии сопряжена с той же тепловой задачей, что и рост толщины отложения (§2.3.2). Сингх, Венкатесан, Фоглер и Нагараджан впоследствии исследовали морфологическую эволюцию уже сформировавшихся толстых восковых отложений при длительном старении, показав качественные изменения кристаллической структуры и распределения захваченной нефти по толщине слоя со временем эксплуатации (Singh, Venkatesan, Fogler, Nagarajan, 2001).

\subsubsection{2.4.3. Следствия для эксплуатации: механическая прочность и удаление отложений}\label{ux441ux43bux435ux434ux441ux442ux432ux438ux44f-ux434ux43bux44f-ux44dux43aux441ux43fux43bux443ux430ux442ux430ux446ux438ux438-ux43cux435ux445ux430ux43dux438ux447ux435ux441ux43aux430ux44f-ux43fux440ux43eux447ux43dux43eux441ux442ux44c-ux438-ux443ux434ux430ux43bux435ux43dux438ux435-ux43eux442ux43bux43eux436ux435ux43dux438ux439}

Практическое значение процесса старения состоит в том, что молодое (недавно образовавшееся) отложение существенно мягче и легче удаляется механическими или химическими методами (скребкование, растворители), чем «созревший» гель того же возраста в трубопроводе, эксплуатирующемся длительное время без очистки. Это создаёт критическую зависимость эффективности эксплуатационных мероприятий (периодичности очистки скребком, дозировки ингибиторов) не только от скорости роста толщины отложения, но и от кинетики его внутреннего дозревания, что требует явного включения уравнений встречной диффузии (§2.4.2) в промысловые расчётные модели, а не только уравнения роста толщины (§2.3.2).

\subsubsection{2.4.4. Старение асфальтеновых и смолистых отложений}\label{ux441ux442ux430ux440ux435ux43dux438ux435-ux430ux441ux444ux430ux43bux44cux442ux435ux43dux43eux432ux44bux445-ux438-ux441ux43cux43eux43bux438ux441ux442ux44bux445-ux43eux442ux43bux43eux436ux435ux43dux438ux439}

Для асфальтеновых отложений аналогичный по духу, но химически иной процесс старения связан с постепенной ароматизацией/конденсацией отложившегося материала и десорбцией стабилизирующих смол в окружающую фазу, что повышает степень агрегации и снижает растворимость осевшего материала со временем, затрудняя последующее растворение отложений органическими растворителями. Количественное описание этого процесса менее развито по сравнению с парафиновыми гелями и по большей части опирается на перенос концепций встречной диффузии и термодинамики агрегации (см. §1.5.4), адаптированных к пористой структуре осевшего асфальтенового слоя.

\subsection{2.5. Резюме раздела}\label{ux440ux435ux437ux44eux43cux435-ux440ux430ux437ux434ux435ux43bux430-1}

Кинетическое описание осаждения тяжёлых компонентов нефти дополняет термодинамическое равновесие (Раздел 1) четырьмя элементами: (1) физикой транспорта растворённых молекул и дисперсных частиц к охлаждённой стенке поры/трубы --- молекулярной диффузией по Фику (доминирующий механизм для растворённых парафинов, Burger et al.~1981; Weingarten, Euchner, 1988; Brown, Niesen, Erickson, 1993), сдвиговой дисперсией, броуновской диффузией и гравитационным оседанием (существенны для дисперсных частиц воска и особенно асфальтенов); (2) кинетикой собственно кристаллизации/агрегации, конкурирующей по времени с транспортом и корректирующей идеализированное равновесное предположение (Huang, Lee, Senra, Fogler, 2011); (3) сопряжённой тепловой задачей, задающей движущую силу диффузионного переноса через радиальный градиент температуры; (4) моделью старения отложений как процесса встречной диффузии компонентов различного углеродного числа внутри пористого гель-каркаса, приводящего к отвердеванию и изменению состава отложения со временем (Singh, Venkatesan, Fogler, Nagarajan, 2000, 2001). Совокупность термодинамических (Раздел 1) и кинетических (Раздел 2) моделей образует физико-химическую базу для построения замкнутой математической модели неизотермической многофазной фильтрации нефти в пласте с учётом осаждения парафинов, асфальтенов и смол.

\newpage

\section{Раздел 3. Связанные неизотермические модели фильтрации в пористой среде}\label{ux440ux430ux437ux434ux435ux43b-3.-ux441ux432ux44fux437ux430ux43dux43dux44bux435-ux43dux435ux438ux437ux43eux442ux435ux440ux43cux438ux447ux435ux441ux43aux438ux435-ux43cux43eux434ux435ux43bux438-ux444ux438ux43bux44cux442ux440ux430ux446ux438ux438-ux432-ux43fux43eux440ux438ux441ux442ux43eux439-ux441ux440ux435ux434ux435}

\subsection{3.1. Постановка задачи и общая структура модели}\label{ux43fux43eux441ux442ux430ux43dux43eux432ux43aux430-ux437ux430ux434ux430ux447ux438-ux438-ux43eux431ux449ux430ux44f-ux441ux442ux440ux443ux43aux442ux443ux440ux430-ux43cux43eux434ux435ux43bux438}

Математическое описание неизотермической фильтрации многокомпонентной нефти в пласте, сопровождающейся кристаллизацией и осаждением тяжёлых компонентов (парафинов, асфальтенов, смол), строится как связанная (coupled) система дифференциальных уравнений в частных производных, включающая как минимум четыре блока: (i) уравнения материального баланса (неразрывности) для каждой фазы и компонента; (ii) закон фильтрации (обобщённый закон Дарси, многофазный, с относительными фазовыми проницаемостями); (iii) уравнение сохранения энергии, учитывающее конвективный и кондуктивный перенос тепла и источники/стоки, связанные с фазовыми переходами; (iv) замыкающие соотношения фазового равновесия (термодинамическая модель растворимости/кристаллизации), определяющие количество и состав выпадающей твёрдой фазы как функцию локальных давления, температуры и состава. Общий методологический каркас такого рода связанных «тепло-массоперенос + фазовый переход» моделей в пористых средах подробно систематизирован для родственного класса задач --- неизотермической фильтрации с разложением газовых гидратов, где авторы явно указывают, что базовые уравнения сохранения массы, закон Дарси и уравнение сохранения энергии дополняются отдельными соотношениями для скрытой теплоты фазового перехода и параметров неизотермической фильтрации как это систематически показано для родственного класса задач --- неизотермической фильтрации с разложением газовых гидратов, где базовые уравнения сохранения массы, закон Дарси и уравнение сохранения энергии дополняются отдельными соотношениями для скрытой теплоты фазового перехода (Бородин, Мусакаев, 2022) --- структурно эта схема прямо переносится на задачу осаждения парафинов/асфальтенов, поскольку в обоих случаях речь идёт о движущемся межфазном фронте кристаллизации внутри порового пространства с локальным тепловыделением.

\subsection{3.2. Уравнения неразрывности (материального баланса)}\label{ux443ux440ux430ux432ux43dux435ux43dux438ux44f-ux43dux435ux440ux430ux437ux440ux44bux432ux43dux43eux441ux442ux438-ux43cux430ux442ux435ux440ux438ux430ux43bux44cux43dux43eux433ux43e-ux431ux430ux43bux430ux43dux441ux430}

Для многофазной (нефть--вода--газ, при необходимости) и многокомпонентной системы уравнение неразрывности записывается для каждого компонента \(i\) (парафин/асфальтен рассматривается как псевдокомпонент или сумма фракций с разными температурами плавления) в форме

\[
\frac{\partial}{\partial t}\left(\phi \sum_{\alpha} \rho_\alpha S_\alpha \omega_{i\alpha}\right) + \nabla\cdot\left(\sum_\alpha \rho_\alpha \omega_{i\alpha}\, \mathbf{u}_\alpha\right) = -\dot{m}_i^{dep} + q_i,
\]

где \(\phi\) --- пористость, \(S_\alpha\) --- насыщенность фазы \(\alpha\) (нефть, вода, газ, а также твёрдая осаждённая фаза, если она выделяется в отдельный ``квази-компонент''), \(\rho_\alpha\) --- плотность фазы, \(\omega_{i\alpha}\) --- массовая доля компонента \(i\) в фазе \(\alpha\), \(\mathbf{u}_\alpha\) --- скорость фильтрации фазы, \(\dot{m}_i^{dep}\) --- массовый сток компонента \(i\) из подвижной фазы в твёрдые отложения (кинетика осаждения), \(q_i\) --- источник/сток на скважине. Для системы в целом накладывается условие \(\sum_\alpha S_\alpha + S_{dep}=1\) (с учётом объёма, занятого отложениями \(S_{dep}\), которая уменьшает эффективную пористость). Такая обобщённая многокомпонентная запись баланса массы согласуется с формулировками, используемыми при моделировании осаждения асфальтенов в пористой среде, где твёрдая осаждённая фаза выделяется явным образом как отдельное неподвижное «депонированное» состояние с собственным уравнением баланса и кинетикой обмена с подвижной нефтяной фазой (Monteagudo, Rajagopal, Lage, 2002) --- именно в этой работе явное разделение асфальтена на растворённый, флоккулированный (переносимый) и осаждённый (неподвижный) компоненты вводится через отдельные уравнения баланса массы, связанные кинетическими членами захвата/отрыва частиц.

\subsection{3.3. Закон Дарси (многофазная фильтрация)}\label{ux437ux430ux43aux43eux43d-ux434ux430ux440ux441ux438-ux43cux43dux43eux433ux43eux444ux430ux437ux43dux430ux44f-ux444ux438ux43bux44cux442ux440ux430ux446ux438ux44f}

Движение каждой подвижной фазы описывается обобщённым законом Дарси с учётом относительных фазовых проницаемостей:

\[
\mathbf{u}_\alpha = -\frac{k\, k_{r\alpha}(S_w, S_o, S_g)}{\mu_\alpha(T, p, \mathbf{\omega})}\left(\nabla p_\alpha - \rho_\alpha \mathbf{g}\right),
\]

где \(k\) --- абсолютная проницаемость пласта (зависящая от степени повреждения --- см. Раздел 4), \(k_{r\alpha}\) --- относительная фазовая проницаемость, \(\mu_\alpha\) --- вязкость фазы, существенно зависящая от температуры (для парафинистых нефтей --- резко возрастающая при приближении к температуре начала кристаллизации), \(p_\alpha\) --- давление фазы (с учётом капиллярного давления между фазами), \(\mathbf{g}\) --- ускорение свободного падения. Капиллярные давления \(p_{cow}=p_o-p_w\), \(p_{cog}=p_g-p_o\) замыкают систему давлений. В условиях неизотермической фильтрации ключевая особенность состоит в сильной температурной зависимости \(\mu_o(T)\) --- вязкость парафинистой нефти может возрастать на порядки при охлаждении ниже температуры насыщения парафином (WAT, wax appearance temperature), что делает систему существенно нелинейной и требует итерационного согласования полей давления, температуры и насыщенностей на каждом временном шаге.

\subsection{3.4. Уравнение энергии}\label{ux443ux440ux430ux432ux43dux435ux43dux438ux435-ux44dux43dux435ux440ux433ux438ux438}

Уравнение сохранения энергии для пористой среды в предположении локального термического равновесия между фазами (одна общая температура \(T\) для скелета и насыщающих фаз) записывается в виде

\[
\frac{\partial}{\partial t}\Big[(1-\phi)\rho_s c_{ps} T + \phi\sum_\alpha \rho_\alpha S_\alpha c_{p\alpha} T\Big] + \nabla\cdot\Big(\sum_\alpha \rho_\alpha c_{p\alpha} T\, \mathbf{u}_\alpha\Big) = \nabla\cdot(\lambda_{eff}\nabla T) + Q_{dep} + Q_{well} ,
\]

где \(\rho_s, c_{ps}\) --- плотность и теплоёмкость твёрдого скелета породы, \(\lambda_{eff}\) --- эффективная теплопроводность насыщенной пористой среды (обычно берётся как взвешенная по объёмным долям сумма теплопроводностей скелета и флюидов, либо по более сложным смесевым правилам), \(Q_{well}\) --- источник/сток тепла на скважине (например, при закачке холодной воды). Левая часть содержит нестационарный член накопления внутренней энергии и конвективный перенос тепла движущимися фазами; правая часть --- кондуктивный перенос тепла (закон Фурье) и объёмные источники. Именно такая балансовая структура --- накопление + конвекция + кондукция + источники фазового перехода --- является общей чертой всех неизотермических моделей фильтрации с фазовыми переходами, что подчёркивается в обзорной работе по неизотермической фильтрации с разложением газогидратов, где отдельно выводятся члены уравнения сохранения энергии и вклад скрытой теплоты фазового перехода как показано в обзорной работе по неизотермической фильтрации с разложением газогидратов, где отдельно выводятся члены уравнения сохранения энергии и вклад скрытой теплоты фазового перехода (Бородин, Мусакаев, 2022)

Ключевой член \(Q_{dep}\) отвечает за тепловыделение при кристаллизации/осаждении тяжёлых компонентов и рассматривается отдельно ниже.

\subsection{3.5. Учёт скрытой теплоты кристаллизации и осаждения тяжёлых компонентов}\label{ux443ux447ux451ux442-ux441ux43aux440ux44bux442ux43eux439-ux442ux435ux43fux43bux43eux442ux44b-ux43aux440ux438ux441ux442ux430ux43bux43bux438ux437ux430ux446ux438ux438-ux438-ux43eux441ux430ux436ux434ux435ux43dux438ux44f-ux442ux44fux436ux451ux43bux44bux445-ux43aux43eux43cux43fux43eux43dux435ux43dux442ux43eux432}

Кристаллизация парафинов (и, с оговорками, флоккуляция/осаждение асфальтенов) представляет собой экзотермический фазовый переход ``жидкость → твёрдое тело'', сопровождающийся выделением скрытой теплоты плавления/кристаллизации \(\Delta H_f\). Наиболее методически проработанный подход к включению этого эффекта в уравнение энергии --- \textbf{энтальпийно-пористый (enthalpy--porosity) метод}, заимствованный из теории затвердевания сплавов и адаптированный для описания отложения парафина в трубах и, по аналогии, в пористых средах. В этом подходе полная удельная энтальпия смеси представляется как сумма чувствительной (sensible) составляющей и скрытой теплоты, взвешенной по локальной доле твёрдой фазы \(f_s(T)\):

\[
H(T) = \int_{T_{ref}}^{T} c_p\, dT' + f_s(T)\,\Delta H_f, \qquad
Q_{dep} = -\rho \frac{\partial (f_s \Delta H_f)}{\partial t},
\]

где \(f_s(T)\) --- локальная доля твёрдой (кристаллизовавшейся) фазы, определяемая термодинамическим равновесием (см. ниже), плавно меняющаяся от 0 (выше температуры начала кристаллизации, WAT) до некоторого предельного значения при глубоком охлаждении. Такая формулировка, объединяющая уравнение энергии с необратимой термодинамикой фазового перехода, детально разработана для трубопроводного транспорта парафинистой нефти (Banki, Hoteit, Firoozabadi, 2008) и по своей структуре напрямую переносима на фильтрационные задачи в пласте, где роль ``твёрдой стенки трубы'' играет скелет породы, а роль пограничного слоя расплава --- прилегающая к зерну порового канала плёнка нефти.

Количество выпадающего твёрдого компонента (а значит, и величина скрытого тепловыделения) определяется многокомпонентным термодинамическим равновесием ``жидкость--твёрдое тело''. Для парафинов классической и наиболее широко применяемой основой служит \textbf{мультитвердофазная (multisolid-phase) термодинамическая модель}, в которой каждая парафиновая фракция (по числу атомов углерода в цепи) может образовывать собственную твёрдую фазу с индивидуальной температурой и энтальпией плавления, а равновесие между жидкой и множеством твёрдых фаз находится из условия равенства химических потенциалов (фугитивностей) компонента в жидкой и твёрдой фазах (Lira-Galeana, Firoozabadi, Prausnitz, 1996). Для асфальтенов аналогичную роль играют модели коллоидной устойчивости и жидкостно-жидкостного/твёрдо-жидкостного равновесия, восходящие к описанию асфальтена как коллоидных частиц, стабилизированных смолами и дестабилизирующихся при изменении давления, температуры и состава (падение давления ниже давления насыщения, смешение с растворителями, изменение температуры) --- систематический обзор полевых наблюдений и подходов к моделированию даёт классическая работа (Leontaritis, Mansoori, 1988). В обоих случаях выход твёрдой фазы \(f_s\) (или концентрация выпавшего асфальтена) вычисляется как функция локальных \(T, p\) и полного компонентного состава, что и подставляется в источниковый член \(Q_{dep}\) уравнения энергии и в стоковый член \(\dot m_i^{dep}\) уравнений неразрывности.

Более поздние обзоры моделирования отложения асфальтенов систематизируют весь спектр подходов --- от чисто термодинамических (равновесных) моделей выпадения до кинетических моделей осаждения частиц на скелете породы, --- подчёркивая, что полная модель обязана включать оба звена: (1) термодинамику фазового равновесия, определяющую \emph{сколько} твёрдой фазы образуется, и (2) кинетику массопереноса/захвата частиц, определяющую \emph{где} и \emph{с какой скоростью} эта твёрдая фаза оседает на стенках порового пространства (Alhosani, Ravichandran, Daraboina, 2020). Аналогичным образом всесторонний обзор механизмов, теории и практических рекомендаций по отложению асфальтенов в пластовых условиях объединяет термодинамические, коллоидные и гидродинамические аспекты задачи (Alimohammadi, Zendehboudi, James, 2019).

\subsection{3.6. Кинетика осаждения и обратная связь с полем скоростей}\label{ux43aux438ux43dux435ux442ux438ux43aux430-ux43eux441ux430ux436ux434ux435ux43dux438ux44f-ux438-ux43eux431ux440ux430ux442ux43dux430ux44f-ux441ux432ux44fux437ux44c-ux441-ux43fux43eux43bux435ux43c-ux441ux43aux43eux440ux43eux441ux442ux435ux439}

Помимо равновесного количества выпавшей твёрдой фазы, необходима кинетическая модель, описывающая скорость её реального осаждения на стенках порового канала (в отличие от простого мгновенного равновесия). Стандартный подход --- линейная (или более сложная, зависящая от локальной скорости фильтрации и диффузии) кинетика захвата частиц по типу модели глубинной фильтрации:

\[
\dot m_i^{dep} = k_{dep}(T, |\mathbf{u}|)\, C_i - k_{ent}(|\mathbf{u}|)\, \sigma_i ,
\]

где \(C_i\) --- концентрация переносимого (взвешенного/растворённого) компонента \(i\), \(\sigma_i\) --- концентрация уже осевшего материала (объём отложений на единицу объёма породы), \(k_{dep}\) --- коэффициент осаждения (может зависеть от температуры через вязкость и от локальной скорости через число Пекле диффузионно-конвективного переноса к стенке поры), \(k_{ent}\) --- коэффициент увлечения/отрыва частиц потоком (entrainment), заметно возрастающий при высоких скоростях фильтрации. Именно такая форма --- совместное уравнение баланса ``взвешенная фаза ⇄ осевшая фаза'' с конкурирующими членами захвата и увлечения --- была впервые систематически обоснована для мигрирующих в пористой среде мелких частиц (Gruesbeck, Collins, 1982) и впоследствии адаптирована для термически и химически индуцированного осаждения парафинов и асфальтенов. Позднее обновлённые формулировки кинетики захвата/удаления частиц в насыщенных пористых средах, учитывающие немонотонную зависимость коэффициентов от локальной скорости и ионной силы раствора, представлены в работе (Civan, 2015).

Для парафиновых отложений в пласте отдельная строка развития --- модели, явно учитывающие связанный тепло-массоперенос при закачке холодной воды (типичный сценарий поддержания пластового давления заводнением), когда фронт охлаждения продвигается впереди фронта вытеснения и провоцирует локальную кристаллизацию парафина именно в охлаждённой зоне. Такая система уравнений (неразрывность нефти и воды + закон Дарси для двух фаз + уравнение энергии с конвекцией от закачиваемой холодной воды + кинетика осаждения парафина, зависящая от локального переохлаждения относительно температуры кристаллизации) детально построена и решена численно в работе (Никифоров, Садовников, Никифоров, 2022), что представляет собой один из немногих примеров полностью связанной неизотермической многофазной модели именно для \emph{пластовой} (а не трубопроводной) задачи. Более ранняя базовая модель отложения парафина непосредственно в пластовых условиях с учётом теплопереноса и кинетики кристаллизации в радиальном течении к скважине предложена в работе (Ring, Wattenbarger, Keating, 1994) --- она остаётся одной из первых полноценных инженерных моделей связанного тепло-массопереноса с осаждением твёрдой фазы именно для околоскважинной зоны пласта.

\subsection{3.7. Типичные упрощающие допущения}\label{ux442ux438ux43fux438ux447ux43dux44bux435-ux443ux43fux440ux43eux449ux430ux44eux449ux438ux435-ux434ux43eux43fux443ux449ux435ux43dux438ux44f}

При построении и численной реализации связанных неизотермических моделей систематически используется ряд упрощений, продиктованных как физической разумностью, так и вычислительной сложностью полной системы:

\begin{enumerate}
\def\labelenumi{\arabic{enumi}.}
\tightlist
\item
  \textbf{Локальное термодинамическое (и часто термическое) равновесие.} Предполагается, что на масштабе элементарного объёма (представительного элементарного объёма, REV) фазовое равновесие ``жидкость--твёрдое тело'' устанавливается мгновенно по сравнению с характерным временем конвективного переноса, так что состав и количество выпавшей твёрдой фазы в каждой точке определяются исключительно локальными \(T, p, \omega\) без учёта запаздывания. Это допущение оправдано для медленной фильтрации, но становится сомнительным при быстрых температурных фронтах (например, при интенсивной закачке холодной воды), где вводятся конечные кинетические константы \(k_{dep}\) (см. п. 3.6) как компромисс между полным равновесием и полной кинетикой.
\item
  \textbf{Локальное термическое равновесие фаз со скелетом} (single-temperature model) --- как правило, принимается единая температура для твёрдого скелета, нефти, воды и осевшей твёрдой фазы в каждой точке, что упрощает уравнение энергии до одного уравнения (см. п. 3.4) вместо системы связанных уравнений для каждой фазы отдельно; двухтемпературные (LTNE, local thermal non-equilibrium) модели используются существенно реже и оправданы лишь при малых временах или высокой скорости фильтрации.
\item
  \textbf{Пренебрежение вязкой диссипацией и работой сил давления} в уравнении энергии --- при типичных скоростях пластовой фильтрации эти члены на порядки меньше конвективного и кондуктивного переноса и обычно опускаются.
\item
  \textbf{Квазистационарность поля давления} на каждом временном шаге теплопереноса (расщепление по физическим процессам, operator splitting) --- учитывая, что характерное время установления поля давления на порядки меньше характерного времени распространения температурного фронта, во многих реализациях уравнение фильтрации решается как квазистационарное на ``замороженном'' (текущем) поле проницаемости и вязкости, тогда как поля температуры и насыщенностей эволюционируют на более медленном временном шаге.
\item
  \textbf{Идеализация геометрии порового пространства} для целей вывода эффективных коэффициентов (эффективной теплопроводности, эффективной кинетики осаждения) --- обычно используется представление порового пространства в виде пучка параллельных капилляров или гранулярной упаковки сфер, что и лежит в основе моделей типа Козени--Кармана (см. Раздел 4).
\item
  \textbf{Аддитивность вкладов компонентов в скрытую теплоту} --- при многокомпонентном описании парафиновых фракций часто принимается, что суммарная скрытая теплота осаждения есть простая сумма вкладов отдельных углеродных фракций, без учёта возможных эффектов смешения между образующимися кристаллическими фазами.
\end{enumerate}

\subsection{3.8. Граничные и начальные условия для скважинных задач}\label{ux433ux440ux430ux43dux438ux447ux43dux44bux435-ux438-ux43dux430ux447ux430ux43bux44cux43dux44bux435-ux443ux441ux43bux43eux432ux438ux44f-ux434ux43bux44f-ux441ux43aux432ux430ux436ux438ux43dux43dux44bux445-ux437ux430ux434ux430ux447}

Специфика скважинных (околоскважинных) задач диктует характерный набор граничных и начальных условий:

\begin{itemize}
\tightlist
\item
  \textbf{Начальные условия}: однородное начальное распределение давления \(p(r,0)=p_i\) (пластовое давление), температуры \(T(r,0)=T_i\) (геотермическая температура пласта на данной глубине) и насыщенностей \(S_\alpha(r,0)=S_{\alpha,i}\) (начальная нефте-, водо-, газонасыщенность), при этом обычно считается, что весь тяжёлый компонент изначально растворён (осевшая фаза отсутствует: \(\sigma_i(r,0)=0\)), поскольку пластовая температура выше температуры кристаллизации в невозмущённой зоне.
\item
  \textbf{Условия на скважине (внутренняя граница, \(r=r_w\))}: либо заданный дебит (условие Неймана для потока массы каждой фазы), либо заданное забойное давление (условие Дирихле для \(p\)); для температуры --- либо заданная температура закачиваемого флюида (при заводнении холодной водой, что и провоцирует кристаллизацию парафина в прискважинной зоне --- сценарий, детально смоделированный в работе Никифорова с соавторами), либо условие адиабатической (теплоизолированной) стенки скважины, либо условие теплообмена с окружающей средой через насосно-компрессорные трубы (более характерно для трубопроводных, а не пластовых моделей).
\item
  \textbf{Условия на внешней границе дренирования (\(r=r_e\) или на бесконечности)}: непроницаемая граница (условие Неймана нулевого потока, \(\partial p/\partial r=0\)) для замкнутого элемента разработки, либо условие постоянного давления \(p(r_e,t)=p_e\) на контуре питания; для температуры --- либо постоянная невозмущённая геотермическая температура \(T(r_e,t)=T_i\), либо условие нулевого теплового потока при достаточно большом радиусе дренирования, за пределами которого тепловое возмущение от скважины не успевает распространиться за время расчёта.
\item
  \textbf{Условия сопряжения на движущемся фронте кристаллизации} (для моделей с явным выделением фронта, а не с ``размазанной'' по объёму enthalpy-porosity формулировкой) --- условие Стефана, связывающее скачок теплового потока на границе раздела ``переохлаждённая/чистая'' зоны со скоростью продвижения фронта и скрытой теплотой кристаллизации; на практике для пластовых задач чаще используется именно сквозной (enthalpy-porosity) метод без явного выделения фронта, поскольку положение и форма фронта кристаллизации в трёхмерной неоднородной пористой среде плохо определены геометрически.
\item
  \textbf{Симметрия задачи}: подавляющее большинство скважинных моделей строится в предположении радиальной (1D в направлении \(r\)) или радиально-вертикальной (2D, \(r\)--\(z\)) симметрии течения к вертикальной скважине, что существенно упрощает численную реализацию связанной системы, хотя и не учитывает возможную анизотропию проницаемости и азимутальную неоднородность реального пласта.
\end{itemize}

Таким образом, полная связанная неизотермическая модель фильтрации многокомпонентной нефти с осаждением тяжёлых компонентов представляет собой систему из уравнений неразрывности по компонентам (с явным выделением стока в осевшую фазу), многофазного закона Дарси с температурно-зависимой вязкностью, уравнения энергии с источником скрытой теплоты фазового перехода и замыкающих термодинамических (равновесие ``жидкость--твёрдое тело'') и кинетических (скорость осаждения/увлечения) соотношений --- при этом ключевая обратная связь этой системы с изменяющимся полем проницаемости и пористости пласта рассматривается отдельно в Разделе 4.

\newpage

\section{Раздел 4. Повреждение пласта и обратная связь}\label{ux440ux430ux437ux434ux435ux43b-4.-ux43fux43eux432ux440ux435ux436ux434ux435ux43dux438ux435-ux43fux43bux430ux441ux442ux430-ux438-ux43eux431ux440ux430ux442ux43dux430ux44f-ux441ux432ux44fux437ux44c}

\subsection{4.1. Физика снижения проницаемости и пористости при отложении твёрдой фазы}\label{ux444ux438ux437ux438ux43aux430-ux441ux43dux438ux436ux435ux43dux438ux44f-ux43fux440ux43eux43dux438ux446ux430ux435ux43cux43eux441ux442ux438-ux438-ux43fux43eux440ux438ux441ux442ux43eux441ux442ux438-ux43fux440ux438-ux43eux442ux43bux43eux436ux435ux43dux438ux438-ux442ux432ux451ux440ux434ux43eux439-ux444ux430ux437ux44b}

Отложение парафинов, асфальтенов и смол на стенках поровых каналов представляет собой процесс постепенного заполнения твёрдым материалом части порового пространства, физически эквивалентный явлению \textbf{повреждения пласта (formation damage)} --- необратимому или частично обратимому снижению эффективной проницаемости призабойной зоны или всего коллектора вследствие уменьшения полезного (доступного для течения) объёма пор и сужения поровых горл. Систематическое описание механизмов такого рода повреждения --- миграция и осаждение мелких частиц, набухание глин, выпадение неорганических солей и органических тяжёлых компонентов --- было впервые формализовано применительно к мигрирующим мелкодисперсным частицам в модели, вводящей раздельные (взаимодействующие через кинетику захвата/увлечения) концентрации взвешенной и осевшей фаз (Gruesbeck, Collins, 1982); впоследствии этот формализм стал стандартным каркасом и для органических отложений.

Ключевой физический механизм следующий: по мере охлаждения нефти ниже температуры кристаллизации парафина (WAT) или при снижении давления ниже давления начала выпадения асфальтенов, твёрдая фаза зарождается в объёме порового флюида (гомогенная или гетерогенная нуклеация на поверхности зерна) и затем либо остаётся во взвешенном/коллоидном состоянии и переносится потоком, либо необратимо прилипает (adheres) к стенке поровых каналов, уменьшая их эффективный радиус. Экспериментальные исследования отложения асфальтенов в модельных пористых средах непосредственно демонстрируют, что даже небольшое по массе отложение твёрдой фазы способно резко и нелинейно снижать проницаемость, поскольку отложение преимущественно происходит в наиболее узких поровых горлах, которые дают основной вклад в фильтрационное сопротивление (Taheri-Shakib, Shekarifard, Naderi, 2018). Прямая визуализация методом рентгеновской микротомографии совместно с флуоресцентной микроскопией подтверждает предпочтительное отложение асфальтенов в узких поровых горлах и на контактах зёрен, что и объясняет непропорционально сильное снижение проницаемости относительно объёмной доли отложений (Feng, Zeng и др., 2017).

Таким образом, физически различают два конкурирующих эффекта повреждения:
1. \textbf{Объёмный (bulk) эффект} --- уменьшение доступной для течения пористости на величину объёма отложений: \(\phi = \phi_0 - \sigma/\rho_{dep}\), где \(\sigma\) --- масса отложений на единицу объёма породы, \(\rho_{dep}\) --- плотность отложенного материала.
2. \textbf{Структурный (pore-throat) эффект} --- непропорционально сильное снижение проницаемости относительно снижения пористости, поскольку сужение поровых горл (где сосредоточено основное гидравлическое сопротивление по закону Пуазёйля, \(\propto r^4\)) даёт квадратичный или более сильный эффект на проницаемость по сравнению с линейным эффектом на пористость.

\subsection{4.2. Модель Козени--Кармана и её модификации}\label{ux43cux43eux434ux435ux43bux44c-ux43aux43eux437ux435ux43dux438ux43aux430ux440ux43cux430ux43dux430-ux438-ux435ux451-ux43cux43eux434ux438ux444ux438ux43aux430ux446ux438ux438}

Классическая \textbf{модель Козени--Кармана} связывает абсолютную проницаемость пласта \(k\) с пористостью \(\phi\) и удельной поверхностью порового пространства через представление пористой среды как пучка извилистых капилляров:

\[
k = \frac{\phi^3}{c_{KC}\, \tau^2\, S_v^2\,(1-\phi)^2},
\]

где \(c_{KC}\) --- константа Козени (обычно \(\approx 5\) для гранулярных сред), \(\tau\) --- извилистость (tortuosity), \(S_v\) --- удельная поверхность твёрдой фазы на единицу объёма скелета. Для целей моделирования повреждения пласта отложениями удобнее использовать относительную (нормированную на начальное состояние) форму:

\[
\frac{k}{k_0} = \left(\frac{\phi}{\phi_0}\right)^{3} \left(\frac{1-\phi_0}{1-\phi}\right)^{2},
\]

которая при \(\phi \to \phi_0\) (отсутствие отложений) даёт \(k\to k_0\), а при уменьшении \(\phi\) вследствие отложений --- резкое, куб-квадратичное падение \(k\). Однако классическая форма Козени--Кармана, выведенная для однородной по размеру зернистой упаковки, систематически завышает проницаемость при значительном заполнении порового пространства мелкодисперсным материалом (частицами, кристаллами), поскольку не учитывает избирательную закупорку именно узких поровых горл. Это было продемонстрировано и скорректировано в модифицированной модели, вводящей эмпирическую степенную поправку с настраиваемым показателем, специально калиброванную по данным о снижении проницаемости от отложения мигрирующих мелких частиц (fines) в чистых неконсолидированных пористых средах (Krauss, Mays, 2013). В таких модификациях типичная форма записывается как

\[
\frac{k}{k_0} = \left(\frac{\phi}{\phi_0}\right)^{n},
\]

где показатель степени \(n\) подбирается эмпирически (часто \(n \gg 3\), вплоть до \(n\sim 8\)--\(19\) в зависимости от типа отложений и структуры порового пространства) и явно отражает то, что реальное снижение проницаемости при отложении твёрдой фазы существенно сильнее, чем предсказывает классическая кубическая зависимость Козени--Кармана.

Для многокомпонентных отложений (одновременное присутствие парафиновой, асфальтеновой и смолистой фракций с разными механизмами прилипания --- кристаллизация на охлаждённой стенке для парафина против коллоидной флоккуляции и адсорбции для асфальтена) естественным обобщением служит взвешенная по вкладам компонентов эффективная удельная поверхность \(S_v^{eff} = S_v^{(0)} + \sum_j \alpha_j \sigma_j\), где \(\alpha_j\) --- эмпирический коэффициент ``эффективности'' отложения компонента \(j\) в снижении гидравлического радиуса пор, что позволяет учесть разную морфологию отложений (плёночную для парафина против агрегатной/кластерной для асфальтена) в рамках единой модификации формулы Козени--Кармана.

Помимо модификаций собственно формулы Козени--Кармана, широко используется и альтернативный, кинетически ориентированный класс моделей --- уравнения баланса взвешенной/осевшей фазы совместно с эмпирической функцией повреждения \(f_d(\sigma)\), связывающей текущую проницаемость с накопленной концентрацией отложений напрямую, без промежуточного пересчёта через пористость:

\[
\frac{k}{k_0} = f_d(\sigma) = \left(1 - \beta\,\frac{\sigma}{\sigma_{max}}\right)^{\gamma},
\]

где \(\beta, \gamma\) --- эмпирические параметры повреждения, \(\sigma_{max}\) --- предельная (насыщающая) концентрация отложений. Именно такой класс моделей --- связывающий динамику захвата/увлечения частиц (см. Раздел 3, п. 3.6) с прямой эмпирической функцией повреждения проницаемости --- систематически развивается и обновляется применительно к формулировкам кинетики осаждения и удаления частиц в насыщенных пористых средах (Civan, 2015), где отдельно подчёркивается необходимость согласования кинетических параметров захвата с эмпирическими корреляциями ``повреждение -- объём отложений'', ранее откалиброванными на прямых лабораторных экспериментах по фильтрации через реальные образцы керна.

Применительно именно к асфальтеновым отложениям --- которые, в отличие от парафина, часто откладываются не сплошной плёнкой, а в виде агрегатов, частично блокирующих поровые горла и меняющих смачиваемость поверхности --- была построена модель повреждения, явно связывающая скорость осаждения асфальтена (через параметр захвата, зависящий от локальной скорости течения и концентрации) с изменением как пористости, так и, отдельно, коэффициента извилистости и удельной поверхности в формуле проницаемости, что позволило воспроизвести наблюдаемое в экспериментах существенно более резкое (по сравнению с чисто объёмным вычетом) снижение продуктивности скважин при асфальтеновом повреждении (Soulgani, Tohidi, Jamialahmadi и др., 2011). Более общий обзорный анализ методов моделирования отложения асфальтенов, объединяющий термодинамические, кинетические и структурно-петрофизические (проницаемость-пористость) блоки модели в единую расчётную схему, дан в работе (Alhosani, Ravichandran, Daraboina, 2020), а всесторонний обзор именно реакции пласта (permeability impairment) на отложение асфальтенов, включая обсуждение применимости и ограничений различных модификаций Козени--Кармана, представлен в работе (Alimohammadi, Zendehboudi, James, 2019).

\subsection{4.3. Обратная связь (feedback loop): отложения → проницаемость → скорость и температура → условия осаждения}\label{ux43eux431ux440ux430ux442ux43dux430ux44f-ux441ux432ux44fux437ux44c-feedback-loop-ux43eux442ux43bux43eux436ux435ux43dux438ux44f-ux43fux440ux43eux43dux438ux446ux430ux435ux43cux43eux441ux442ux44c-ux441ux43aux43eux440ux43eux441ux442ux44c-ux438-ux442ux435ux43cux43fux435ux440ux430ux442ux443ux440ux430-ux443ux441ux43bux43eux432ux438ux44f-ux43eux441ux430ux436ux434ux435ux43dux438ux44f}

Центральная методологическая особенность рассматриваемого класса задач --- существование \textbf{замкнутого контура обратной связи}, принципиально отличающего их от простых ``однонаправленных'' (once-through) моделей переноса примеси. Схематически цепочка обратной связи выглядит следующим образом:

\[
\underbrace{T,\,p,\,\omega}_{\text{локальное состояние}} \;\Rightarrow\; \underbrace{f_s,\ \dot m^{dep}}_{\substack{\text{термодинамика +}\\ \text{кинетика осаждения}}} \;\Rightarrow\; \underbrace{\sigma \to \phi,\, k}_{\text{Козени–Карман}}
\]

\[
\Rightarrow\; \underbrace{\mathbf{u}_\alpha}_{\text{закон Дарси}} \;\Rightarrow\; \underbrace{T,\,p}_{\substack{\text{конвекция +}\\ \text{источник } Q_{dep}}} \;\Rightarrow\; \ldots
\]

Механизм обратной связи работает в обоих направлениях знака:

\begin{itemize}
\tightlist
\item
  \textbf{Отрицательная (стабилизирующая) обратная связь через скорость.} Снижение проницаемости в зоне интенсивного отложения приводит к локальному уменьшению скорости фильтрации \(\mathbf{u}\) (при заданном перепаде давления) либо к необходимости увеличения депрессии на пласт для поддержания заданного дебита. Уменьшение локальной скорости, в свою очередь, снижает конвективный подвод тепла в эту зону (уменьшается адвективный член в уравнении энергии), что при закачке холодной воды означает более медленное продвижение охлаждённого фронта вглубь пласта --- то есть само отложение частично ``экранирует'' зону дальнейшего охлаждения. Одновременно снижение скорости уменьшает коэффициент увлечения \(k_{ent}\) в кинетике осаждения (см. Раздел 3, п. 3.6), что при определённых условиях способствует ещё большему накоплению отложений в этой же зоне (уже положительная связь по массе).
\item
  \textbf{Положительная (усиливающая) обратная связь через локальный перепад давления.} Снижение проницаемости в некоторой зоне при сохранении общего дебита скважины требует увеличения локального градиента давления для проталкивания того же расхода флюида через суженные поровые каналы (согласно закону Дарси \(\mathbf{u}=- (k/\mu)\nabla p\)). Возросший градиент давления способен как повышать локальную скорость сдвига (что увеличивает механическое увлечение и может частично компенсировать дальнейшее осаждение), так и, в асфальтеновых системах, провоцировать дополнительное падение давления ниже точки начала осаждения (asphaltene onset pressure) именно в уже повреждённой зоне, что усиливает дальнейшее выпадение твёрдой фазы --- механизм прогрессирующего повреждения, детально анализируемый в модели прискважинного повреждения при асфальтеновом осаждении (Soulgani, Tohidi, Jamialahmadi и др., 2011).
\item
  \textbf{Локализация повреждения в прискважинной зоне и скин-фактор.} Поскольку скорость фильтрации и, соответственно, градиенты давления и температуры максимальны вблизи ствола скважины (радиальная сходимость линий тока), именно в этой зоне обратная связь ``отложение → снижение \(k\) → рост локального \(\nabla p\) и/или локальное охлаждение → дальнейшее отложение'' проявляется наиболее интенсивно, что и приводит к формированию сосредоточенного вблизи скважины \textbf{дополнительного фильтрационного сопротивления}, количественно описываемого в терминах \textbf{скин-фактора} \(S\). В рамках классической схемы Хокинса скин-фактор для тонкого кольцевого повреждённого слоя радиусом \(r_s\) вокруг скважины радиуса \(r_w\) выражается как
\end{itemize}

\[
S = \left(\frac{k}{k_s}-1\right)\ln\frac{r_s}{r_w},
\]

где \(k\) --- невозмущённая проницаемость пласта, \(k_s\) --- сниженная (повреждённая) проницаемость в кольцевой зоне отложений. При явном моделировании связанной неизотермической задачи величина \(k_s(r,t)\) не задаётся априорно, а вычисляется динамически как результат интегрирования по времени накопленной массы отложений \(\sigma(r,t)\) через модифицированную формулу Козени--Кармана (п. 4.2), так что скин-фактор оказывается не постоянным параметром, а \emph{растущей во времени и меняющейся в пространстве} величиной, явно зависящей от истории температурного и барического режима эксплуатации скважины. Такой динамический подход к формированию зоны пониженной проницаемости вблизи скважины и связанного с ней снижения продуктивности systematически рассматривается в базовой модели прогнозирования отложения парафина в пласте, где радиальная фильтрационная задача решается совместно с уравнением теплопереноса и кинетикой кристаллизации именно для околоскважинной зоны (Ring, Wattenbarger, Keating, 1994), а также в связанной тепло-массообменной модели отложения парафина при закачке холодной воды, где явно показано, что зона максимального охлаждения (и, следовательно, максимального отложения) продвигается вместе с фронтом закачиваемой воды, формируя нестационарную, движущуюся во времени зону повреждения, а не статичное кольцо вокруг скважины (Никифоров, Садовников, Никифоров, 2022).

\subsection{4.4. Замкнутая связанная система: сведение воедино}\label{ux437ux430ux43cux43aux43dux443ux442ux430ux44f-ux441ux432ux44fux437ux430ux43dux43dux430ux44f-ux441ux438ux441ux442ux435ux43cux430-ux441ux432ux435ux434ux435ux43dux438ux435-ux432ux43eux435ux434ux438ux43dux43e}

Итоговая связанная система, объединяющая Раздел 3 и Раздел 4, формально записывается как совместное решение:

\begin{enumerate}
\def\labelenumi{\arabic{enumi}.}
\tightlist
\item
  Уравнений неразрывности по компонентам с кинетическим стоком в отложения (п. 3.2, 3.6);
\item
  Многофазного закона Дарси с проницаемостью, явно зависящей от текущей концентрации отложений через модифицированную формулу Козени--Кармана (п. 4.2): \(k=k(\phi(\sigma), \tau(\sigma), S_v(\sigma))\);
\item
  Уравнения энергии с конвективным членом, зависящим от той же (уже повреждённой) скорости фильтрации, и источником скрытой теплоты кристаллизации/осаждения (п. 3.4--3.5);
\item
  Замыкающих термодинамических соотношений равновесия ``жидкость--твёрдое тело'', определяющих локальную движущую силу осаждения \(f_s(T,p,\omega)\) (п. 3.5).
\end{enumerate}

Эта система является существенно нелинейной и связанной ``в обе стороны'': температура и давление определяют интенсивность осаждения, осаждение определяет проницаемость, проницаемость определяет поле скоростей, поле скоростей определяет конвективный перенос тепла и массы --- то есть замыкает контур обратной связи на температуру и давление. Именно наличие этой петли обратной связи (а не просто последовательного, ``one-way'' расчёта сначала теплопереноса, а затем --- отдельно --- повреждения по готовому температурному полю) является методологически ключевой чертой физически корректных моделей данного класса, отличающей их от упрощённых инженерных оценок скин-фактора по фиксированной эмпирической корреляции. Экспериментальные исследования влияния асфальтеновых отложений на смачиваемость и проницаемость независимо подтверждают, что степень повреждения существенно зависит от локальной истории режима течения (скорости, температуры, перепада давления), а не только от суммарной массы осевшего материала, что и требует именно динамического, связанного по времени моделирования, а не статического пересчёта конечного состояния.

\subsection{4.5. Основные направления дальнейшего развития моделей}\label{ux43eux441ux43dux43eux432ux43dux44bux435-ux43dux430ux43fux440ux430ux432ux43bux435ux43dux438ux44f-ux434ux430ux43bux44cux43dux435ux439ux448ux435ux433ux43e-ux440ux430ux437ux432ux438ux442ux438ux44f-ux43cux43eux434ux435ux43bux435ux439}

Обзор рассмотренной литературы позволяет выделить нерешённые/частично решённые вопросы, актуальные для дальнейшего численного моделирования:

\begin{itemize}
\tightlist
\item
  согласование масштаба применимости модификаций Козени--Кармана (откалиброванных, как правило, на неконсолидированных или слабосцементированных песчаниках) с реальной структурой карбонатных и сильно сцементированных терригенных коллекторов;
\item
  явный учёт многокомпонентности отложений (одновременное присутствие парафина, асфальтена и смол с разной морфологией и разными механизмами прилипания) в единой эффективной корреляции ``проницаемость--пористость'', а не в виде суммы независимых вкладов;
\item
  корректный учёт локального термического неравновесия (LTNE) между твёрдым скелетом и флюидом в зонах резких температурных градиентов (фронт закачки холодной воды), где предположение о единой температуре (п. 3.7) может оказаться неточным;
\item
  валидация динамически вычисляемого скин-фактора на промысловых данных о снижении продуктивности скважин во времени, а не только на лабораторных кернах.
\end{itemize}

\newpage

\section{Раздел 5. Моделирование асфальтенов и смол}\label{ux440ux430ux437ux434ux435ux43b-5.-ux43cux43eux434ux435ux43bux438ux440ux43eux432ux430ux43dux438ux435-ux430ux441ux444ux430ux43bux44cux442ux435ux43dux43eux432-ux438-ux441ux43cux43eux43b}

\subsection{5.1. Физико-химические механизмы дестабилизации и агрегации асфальтенов и смол}\label{ux444ux438ux437ux438ux43aux43e-ux445ux438ux43cux438ux447ux435ux441ux43aux438ux435-ux43cux435ux445ux430ux43dux438ux437ux43cux44b-ux434ux435ux441ux442ux430ux431ux438ux43bux438ux437ux430ux446ux438ux438-ux438-ux430ux433ux440ux435ux433ux430ux446ux438ux438-ux430ux441ux444ux430ux43bux44cux442ux435ux43dux43eux432-ux438-ux441ux43cux43eux43b}

Асфальтены --- не индивидуальное химическое соединение, а операционно определяемый класс наиболее полярных и высокомолекулярных компонентов нефти, нерастворимых в лёгких \emph{н}-алканах (пентан, гептан) и растворимых в ароматических растворителях (толуол, бензол). Молекулярная архитектура асфальтенов на протяжении десятилетий оставалась предметом дискуссии между «архипелажной» моделью (несколько небольших полиароматических ядер, связанных алифатическими мостиками) и моделью «континентального» типа --- единое конденсированное полициклическое ароматическое ядро (7--10 колец), окружённое алифатическими цепями и функциональными группами. Прямые измерения молекулярной массы методами лазерной десорбционной масс-спектрометрии и флуоресцентной деполяризации показали, что доминирующей является модель одного конденсированного ароматического ядра на молекулу с молекулярной массой порядка 500--1000 Да и диаметром порядка 1.5 нм --- так называемая молекула Йена--Маллинса (Yen--Mullins) «Advances in Asphaltene Science and the Yen--Mullins Model» (Mullins et al., 2012). Дальнейшая работа с масс-спектрометрией на основе лазерной десорбции/ионизации подтвердила единую конденсированную полиароматическую архитектуру и позволила количественно оценить число ароматических колец, длину алифатических заместителей и степень агрегации в наноагрегаты (Pomerantz et al., 2015).

Ключевую роль в химической структуре асфальтеновых молекул играют гетероатомы --- сера, азот, кислород, а также металлопорфирины никеля и ванадия, которые встраиваются в полиароматические плоскости и удерживаются преимущественно в асфальтеновой фракции при разделении нефти. Ещё в классической работе Даннинга с соавторами было показано прямое включение порфириновых комплексов никеля и ванадия в структуру тяжёлых фракций нефти, что стало одним из первых указаний на металлоорганическую природу части асфальтеновых макромолекул (Dunning et al., 1960). При осаждении асфальтенов происходит избирательное обогащение осадка ванадий-, никель- и серосодержащими соединениями: экспериментально показано, что доля этих гетероатомных структур в осаждённой фазе растёт с увеличением степени осаждения, что подтверждает их локализацию преимущественно в наиболее конденсированных и полярных молекулах (Reynolds, 1994). Сера, присутствующая в форме тиофеновых и сульфидных группировок непосредственно в ароматических кластерах, модифицирует электронную плотность ядра и усиливает π--π-взаимодействия между соседними молекулами, что напрямую влияет на кинетику и термодинамику агрегации асфальтенов (Sodero et al., 2016).

Дестабилизация асфальтенов инициируется изменением термо-баро-композиционного состояния нефти при фильтрации в пласте и движении по стволу скважины. Наиболее распространённый механизм --- снижение давления ниже давления насыщения (или, для недонасыщенных нефтей, ниже так называемого давления начала осаждения асфальтенов, asphaltene onset pressure, AOP), при котором происходит перераспределение лёгких компонентов между жидкой и паровой фазами, ухудшающее термодинамическое сродство между асфальтеновыми макромолекулами и остальной нефтяной матрицей. При этом зависимость растворимости асфальтенов от давления немонотонна: она максимальна вблизи давления насыщения и снижается как при более высоких, так и при более низких давлениях, что и объясняет экспериментально наблюдаемый колоколообразный профиль осаждения по давлению. Помимо давления, дестабилизирующими факторами выступают изменение состава (смешение с лёгкими углеводородными или CO₂-содержащими агентами при закачке газа, смешение с растворителями при разбавлении тяжёлой нефти) и температура, влияющая как на плотность и энтропийный вклад в свободную энергию смешения, так и на кинетику агрегации через изменение вязкости среды и коэффициентов диффузии наноагрегатов.

\subsection{5.2. Коллоидные модели: модель Йена--Маллинса}\label{ux43aux43eux43bux43bux43eux438ux434ux43dux44bux435-ux43cux43eux434ux435ux43bux438-ux43cux43eux434ux435ux43bux44c-ux439ux435ux43dux430ux43cux430ux43bux43bux438ux43dux441ux430}

Согласно современной коллоидной картине, принятой в модели Йена--Маллинса, структурная иерархия асфальтенов в нефти включает несколько уровней организации: (i) отдельные молекулы асфальтенов, растворённые в мальтеновой фазе при низкой концентрации; (ii) наноагрегаты --- кластеры порядка 6--10 молекул с диаметром около 2 нм, образующиеся выше критической концентрации наноагрегации (CNAC) за счёт стэкинга ароматических плоскостей; (iii) кластеры наноагрегатов --- более крупные ассоциаты порядка 8 наноагрегатов диаметром порядка 5 нм, формирующиеся при более высокой концентрации асфальтенов (критическая концентрация кластерообразования, CCC) «Advances in Asphaltene Science and the Yen--Mullins Model» (Mullins et al., 2012). Именно кластеры наноагрегатов, а не отдельные молекулы или первичные наноагрегаты, считаются частицами, ответственными за нуклеацию макроскопической осаждённой фазы при дальнейшей дестабилизации. Независимая проверка этой иерархии на реальных пластовых флюидах вдоль вертикального профиля коллектора, где гравитационная сегрегация компонентов создаёт естественный градиент термодинамических условий, подтвердила постоянство молекулярной массы и массы наноагрегатов асфальтенов по всей толще пласта, что согласуется с моделью Йена--Маллинса и применимостью простой статистической термодинамики (уравнения Флори--Хаггинса в сочетании с баро-гравитационным распределением) для описания градиентов концентрации асфальтенов (Wu et al., 2014).

Коллоидная модель имеет решающее значение для интерпретации кинетики агрегации: рост частиц происходит по механизму диффузионно-лимитированной или реакционно-лимитированной агломерации (DLCA/RLCA), а не по классическому механизму нуклеации из истинного раствора, что определяет степенные законы роста радиуса частиц во времени и фрактальную размерность образующихся агрегатов. Взаимодействия асфальтеновых наноагрегатов на межфазных границах нефть/вода, где формируются устойчивые адсорбционные плёнки, дополнительно моделируются на основе решёточно-газовых (lattice-gas) представлений, воспроизводящих кооперативную адсорбцию в рамках концепций Йена--Маллинса (Darjani et al., 2022).

\subsection{5.3. Термодинамические модели: SAFT и PC-SAFT}\label{ux442ux435ux440ux43cux43eux434ux438ux43dux430ux43cux438ux447ux435ux441ux43aux438ux435-ux43cux43eux434ux435ux43bux438-saft-ux438-pc-saft}

Термодинамическое описание фазового поведения асфальтенов требует учёта сильной ассоциативной природы взаимодействий (водородные связи гетероатомных группировок, π--π-стэкинг ароматических плоскостей) и существенной полидисперсности по молекулярной массе и степени агрегации, что выходит за рамки применимости классических кубических уравнений состояния (Пенга--Робинсона, Суаве--Редлиха--Квонга) без искусственной подгонки параметров бинарного взаимодействия. Статистическая теория ассоциирующих флюидов (SAFT) и в особенности её версия с возмущённой цепью (Perturbed-Chain SAFT, PC-SAFT) моделируют молекулы как цепочки сегментов с потенциалами твёрдых сфер, дисперсионным взаимодействием и явным вкладом ассоциации, что позволяет физически обоснованно описывать асфальтены как крупные, слабо ассоциирующие или неассоциирующие полидисперсные макромолекулы.

Основополагающее применение PC-SAFT к фазовому поведению асфальтенов было выполнено группой Вентуры Гонзалеса и Варгаса, где асфальтены моделировались как сферические или короткоцепочечные сегменты с параметрами, извлекаемыми из данных по осаждению \emph{н}-алканами; впоследствии подход был переосмыслен с учётом полидисперсности молекулярной массы асфальтенов и more физически обоснованного выбора параметров сегмента и энергии дисперсии, что существенно улучшило предсказательную способность модели при экстраполяции на разные типы нефти и растворителей (Tavakkoli, Chen \& Vargas, 2016). Такой подход математически описывается через химический потенциал каждого псевдокомпонента \(i\), включающий вклады идеального газа, твёрдых цепей, дисперсии Ленарда-Джонса и, при необходимости, ассоциации:

\[
\frac{a^{res}}{RT} = \tilde a^{hc} + \tilde a^{disp} + \tilde a^{assoc},
\]

где \(\tilde a^{hc}\) --- вклад цепочки твёрдых сфер, \(\tilde a^{disp}\) --- дисперсионный член теории возмущений второго порядка, \(\tilde a^{assoc}\) --- вклад ассоциации по теории Вертхайма. Условие равновесия жидкость--жидкость (мальтеновая фаза / асфальтеновая плотная фаза) записывается через равенство летучестей:

\[
f_i^{L1}(T,P,\mathbf{x}^{L1}) = f_i^{L2}(T,P,\mathbf{x}^{L2}), \qquad i=1,\dots,N,
\]

а положение фазовой границы (AOP) находится как давление, при котором система впервые становится нестабильной по отношению к разделению на вторую жидкую фазу, что проверяется через анализ устойчивости (tangent-plane distance / стабильность по Михельсену).

Дальнейшее развитие PC-SAFT-подхода включает его встраивание в открытые платформы для композиционного моделирования пласта: в частности, реализация PC-SAFT в составе MRST Compositional позволила напрямую сопрягать расчёт осаждения асфальтенов с многофазным потоковым моделированием в пористой среде без обращения к внешним PVT-пакетам (Masoudi et al., 2021). Параллельно ведутся систематические работы по улучшению точности PC-SAFT при смешении нефти с различными растворителями и газами закачки, учитывающие переменные параметры бинарного взаимодействия \(k_{ij}(T)\) и альтернативные схемы характеризации асфальтеновой псевдокомпоненты (Kariman Moghaddam \& Jamshidi, 2022).

Помимо PC-SAFT, для описания асфальтеновой агрегации применяются и другие термодинамические подходы: модели ассоциации на основе теории Флори--Хаггинса с явным учётом энтальпии агрегации молекул в наноагрегаты и кластеры «Thermodynamic Modeling of Asphaltene Aggregation»; модель кубического уравнения состояния с ассоциативным членом (Cubic-Plus-Association, CPA), явно учитывающая водородное связывание асфальтенов с водной фазой в трёхфазных системах нефть--вода--асфальтен (Jia \& Okuno, 2018); а также подходы на основе термодинамики агрегации в рамках модели NRTL-SAC (Non-Random Two-Liquid Segment Activity Coefficient), которая рассматривает асфальтеновую агрегацию как последовательность равновесных стадий олигомеризации с константами равновесия, зависящими от температуры и состава растворителя (Wang et al., 2016; Islam, Hao \& Chen, 2020). Ранние феноменологические модели (модель Флори--Хаггинса--Скотта для мицеллообразования асфальтенов; скейлинговые уравнения) остаются актуальными как быстрые инженерные корреляции для оценки AOP на основе минимального набора PVT-данных (Jamialahmadi \& Ahmadi, 2003).

\subsection{5.4. Модели осаждения (кольматации) порового пространства и выноса частиц}\label{ux43cux43eux434ux435ux43bux438-ux43eux441ux430ux436ux434ux435ux43dux438ux44f-ux43aux43eux43bux44cux43cux430ux442ux430ux446ux438ux438-ux43fux43eux440ux43eux432ux43eux433ux43e-ux43fux440ux43eux441ux442ux440ux430ux43dux441ux442ux432ux430-ux438-ux432ux44bux43dux43eux441ux430-ux447ux430ux441ux442ux438ux446}

После термодинамической дестабилизации асфальтеновые кластеры наноагрегатов должны быть перенесены к твёрдой поверхности порового канала и закреплены на ней, чтобы вызвать снижение проницаемости --- этот процесс принципиально отличается от простого выпадения осадка в объёме и требует отдельного описания механизмов переноса, прилипания (adsorption/deposition) и последующего отрыва (entrainment). Экспериментальные исследования на стеклянных микромоделях с регулярной сеткой пор показали, что закупорка отдельных сужений (throats) происходит по механизму «снежного кома» (snow-ball effect): однажды образовавшаяся на стенке поры частица асфальтена служит центром захвата последующих частиц, что приводит к прогрессивно ускоряющемуся росту локальной пробки и резкому, нелинейному во времени падению локальной проницаемости (Onaka \& Sato, 2021).

Наиболее распространённый в литературе математический каркас для количественного описания осаждения и переноса частиц в пористой среде --- модель типа Иванса--Шехтера/Цивана (Civan) с раздельными уравнениями сохранения массы взвешенных частиц в потоке и массы осаждённых (иммобилизованных) частиц на скелете породы:

\[
\phi \frac{\partial c}{\partial t} + \nabla\cdot(\mathbf{u}c) - \nabla\cdot(D\nabla c) = -\frac{\partial \sigma}{\partial t},
\]

\[
\frac{\partial \sigma}{\partial t} = k_d\, c - k_e\, \sigma\, \|\mathbf{u}\|,
\]

где \(c\) --- объёмная концентрация взвешенных (мобильных) частиц асфальтена в несущей фазе, \(\sigma\) --- объёмная концентрация осевших (иммобилизованных) частиц на стенках пор, \(k_d\) --- коэффициент захвата (deposition rate constant), \(k_e\) --- коэффициент выноса/эрозии (entrainment rate constant), пропорциональный локальной скорости фильтрации (или гидродинамическому напряжению сдвига на стенке поры), \(\phi\) --- пористость, \(\mathbf{u}\) --- скорость Дарси. Стохастическая версия этой модели, учитывающая случайное распределение размеров поровых каналов и вероятностный характер захвата частиц данного размера порами данного диаметра, была развита для задач глубинной фильтрации (deep bed filtration) и повреждения призабойной зоны в контексте общей теории кольматации коллектора взвешенными частицами «A Stochastic Model for Deep Bed Filtration and Well Impairment». Снижение проницаемости и пористости при накоплении осадка описывается эмпирическими или полуэмпирическими соотношениями фильтрационного сопротивления, например степенными законами \(k/k_0 = f(\sigma/\phi_0)\) или моделями типа Козени--Кармана с эффективным уменьшением гидравлического радиуса пор.

Важным развитием классической схемы deposition--entrainment стала явная кинетическая модель осаждения асфальтенов, объединяющая термодинамику дестабилизации (доля выпавшей твёрдой фазы, рассчитываемая по модели типа Флори--Хаггинса или PC-SAFT) с кинетикой массопереноса частиц к стенке поры и последующей необратимой или обратимой фиксацией: «``A Rigorous Kinetic Model for Description of Asphaltene Deposition in Porous Media of Petroleum Reservoirs''» явно вводит зависящую от локальной скорости и концентрации кинетическую константу осаждения, что позволяет воспроизвести экспериментально наблюдаемые профили накопления осадка вдоль направления фильтрации в кернах (ZareNezhad \& Parsa, 2015). Похожая по структуре, но ориентированная на призабойную зону и ствол скважины модель динамического многофазного осаждения асфальтенов (Multiphase and Dynamic ADEPT, MAD-ADEPT) сопрягает уравнение баланса массы осевшей фазы с уравнением движения многофазного потока в трубе и явно учитывает как ламинарный, так и турбулентный режимы транспорта частиц к стенке через диффузионно-конвективный поток Шервуда (Naseri, Jamshidi \& Taghikhani, 2020). Предшествующий и широко цитируемый прототип такого рода моделей --- инструмент прогнозирования отложений асфальтенов (Asphaltene Deposition Tool, ADEPT), развитый группой Чапмена и Варгаса и впервые верифицированный на промысловых данных, где скорость осаждения частиц на стенке трубы связывается с локальным градиентом концентрации асфальтенов и коэффициентом массообмена, зависящим от числа Рейнольдса потока (Kurup et al., 2012). Систематическое сопоставление такого рода моделей осаждения в стволе скважины с трёхмерными расчётами методом вычислительной гидродинамики (CFD) для многофазного (нефть--газ) потока в скважине с закачкой газа показало, что радиальное распределение скорости сдвига и локальные зоны рециркуляции существенно модифицируют профиль осаждения по сравнению с одномерными инженерными корреляциями (Rajan Babu et al., 2019).

\subsection{5.5. Сопряжение с термо-баро-композиционными изменениями при фильтрации}\label{ux441ux43eux43fux440ux44fux436ux435ux43dux438ux435-ux441-ux442ux435ux440ux43cux43e-ux431ux430ux440ux43e-ux43aux43eux43cux43fux43eux437ux438ux446ux438ux43eux43dux43dux44bux43cux438-ux438ux437ux43cux435ux43dux435ux43dux438ux44fux43cux438-ux43fux440ux438-ux444ux438ux43bux44cux442ux440ux430ux446ux438ux438}

Полная математическая модель осаждения асфальтенов при неизотермической многокомпонентной фильтрации в пласте требует одновременного решения трёх взаимосвязанных подзадач: (1) фазового равновесия и доли выпавшей твёрдой асфальтеновой фазы как функции локального давления, температуры и состава (через SAFT/PC-SAFT или CPA-EOS); (2) переноса компонентов и фаз в пористой среде (уравнения сохранения массы для каждого компонента с учётом многофазного потока, аналогично композиционному пластовому моделированию); (3) кинетики захвата/выноса твёрдой асфальтеновой фазы на скелете породы и обратного влияния осадка на абсолютную и относительные проницаемости.

Ранняя и до сих пор широко цитируемая работа, объединившая термодинамическое моделирование осаждения асфальтенов (на основе модели твёрдого раствора, solid-solution model) с уравнениями многофазной многокомпонентной фильтрации в композиционном симуляторе, продемонстрировала, что учёт изменения относительных проницаемостей и абсолютной проницаемости вследствие осевшего асфальтена может кардинально (на порядки) изменить прогноз продуктивности скважины по сравнению с расчётом чистого термодинамического осаждения без кольматации «Modeling Asphaltene Precipitation in Reservoir Simulation» (Qin et al., 2000). Аналогичная по духу, но более ранняя по методологии работа явно связала уравнение фильтрации Дарси с локальным изменением проницаемости вследствие асфальтенового осадка через модифицированный закон Козени, вводя коэффициент кольматации, зависящий от локальной концентрации осевшей фазы (Monteagudo, Rajagopal \& Lage, 2002).

Температурная зависимость учитывается на двух уровнях: во-первых, через непосредственное влияние температуры на положение фазовой границы асфальтенового осаждения (растворимость асфальтенов, как правило, слабо, но не монотонно зависит от температуры --- при типичных пластовых условиях повышение температуры может как повышать, так и понижать растворимость в зависимости от преобладающего энтропийного или энтальпийного вклада в свободную энергию смешения); во-вторых, через сопряжение уравнения теплопереноса (конвективно-кондуктивного, с учётом теплоты фазового перехода при выпадении твёрдой фазы) с уравнениями композиционной фильтрации. Такая полная тепло-массообменная и фазовая связка типична для симуляторов термически-композиционного типа, разработанных изначально для методов увеличения нефтеотдачи с закачкой пара или горячей воды, и в последнее время адаптируемых для задач совместного моделирования осаждения асфальтенов и парафинов при неизотермической фильтрации (Khait \& Voskov, 2019).

Отдельного внимания заслуживает немонотонность профиля осаждения по давлению вдоль ствола скважины и в пласте: поскольку максимальная нестабильность асфальтенов достигается вблизи давления насыщения, а не при минимальном давлении, зона максимального осаждения часто локализуется не на забое, а на некотором удалении вверх по потоку (в пласте) или вверх по стволу скважины, что было явно показано и воспроизведено моделью с сопряжённым расчётом термодинамики и течения в композиционном коллекторно-скважинном симуляторе (Darabi et al., 2014). Это давление-зависимое, немонотонное распределение источника твёрдой фазы принципиально отличает задачу осаждения асфальтенов от более простых задач кольматации механическими примесями и требует явного разрешения профиля давления, температуры и состава по всей траектории фильтрации, а не только на условиях сепаратора.

\newpage

\section{Раздел 6. Численные методы и симуляторы}\label{ux440ux430ux437ux434ux435ux43b-6.-ux447ux438ux441ux43bux435ux43dux43dux44bux435-ux43cux435ux442ux43eux434ux44b-ux438-ux441ux438ux43cux443ux43bux44fux442ux43eux440ux44b}

\subsection{6.1. Численные схемы для связанных задач фильтрации, теплопереноса и фазовых переходов}\label{ux447ux438ux441ux43bux435ux43dux43dux44bux435-ux441ux445ux435ux43cux44b-ux434ux43bux44f-ux441ux432ux44fux437ux430ux43dux43dux44bux445-ux437ux430ux434ux430ux447-ux444ux438ux43bux44cux442ux440ux430ux446ux438ux438-ux442ux435ux43fux43bux43eux43fux435ux440ux435ux43dux43eux441ux430-ux438-ux444ux430ux437ux43eux432ux44bux445-ux43fux435ux440ux435ux445ux43eux434ux43eux432}

Математическая модель осаждения асфальтенов и парафинов при неизотермической фильтрации представляет собой систему нелинейных уравнений в частных производных параболического (диффузия, теплоперенос) и гиперболического (адвективный перенос компонентов) типа, дополненную алгебраическими соотношениями фазового равновесия и кинетическими уравнениями для осевшей фазы. Такая система решается, как правило, методом конечных объёмов (finite volume method, FVM) на структурированных или неструктурированных сетках --- этот подход доминирует в пластовом моделировании, поскольку обеспечивает точное локальное сохранение массы каждого компонента и естественно сочетается с двухточечной или многоточечной аппроксимацией потока (Two-Point / Multi-Point Flux Approximation) на сетках произвольной геометрии, включая ячейки с угловыми точками (corner-point grids), типичные для промысловых моделей.

Для компонентного (композиционного) моделирования типична IMPES-схема (Implicit Pressure, Explicit Saturation) или полностью неявная схема (fully implicit method, FIM): давление находится неявно из суммарного уравнения сохранения объёма, тогда как насыщенности и концентрации компонентов могут обновляться как явно, так и неявно в зависимости от требуемой устойчивости по числу Куранта. Один из современных подходов, направленный на снижение вычислительных затрат при сохранении устойчивости, --- полунеявная (semi-implicit) схема конечных объёмов для композиционного моделирования подземного течения, в которой перенос компонентов трактуется явно на основе TVD-ограничителей потока, а давление и общий поток --- неявно, что позволяет использовать существенно больший шаг по времени по сравнению с полностью явными схемами без потери устойчивости, характерной для полностью неявных методов (Echavarría-Montaña et al., 2021). Общая структура полностью неявной постановки для многофазной многокомпонентной системы записывается как система нелинейных уравнений сохранения массы компонента \(k\) в ячейке:

\[
\frac{\partial}{\partial t}\left(\phi \sum_{j} \rho_j x_{kj} S_j\right) + \nabla\cdot\left(\sum_j \rho_j x_{kj} \mathbf{u}_j\right) - \nabla\cdot\left(\phi \sum_j \rho_j S_j D_{kj}\nabla x_{kj}\right) = q_k - r_k^{dep},
\]

где \(S_j\) --- насыщенность фазы \(j\), \(x_{kj}\) --- мольная (массовая) доля компонента \(k\) в фазе \(j\), \(\mathbf{u}_j = -\dfrac{k\,k_{rj}}{\mu_j}(\nabla P_j - \rho_j g\nabla z)\) --- скорость Дарси фазы \(j\) с учётом относительной проницаемости \(k_{rj}\), \(q_k\) --- источниковый/стоковый член, а \(r_k^{dep}\) --- сток компонента \(k\) вследствие осаждения асфальтеновой/парафиновой твёрдой фазы, вычисляемый по кинетической модели раздела 5. Такая система дополняется уравнением теплопереноса

\[
\frac{\partial}{\partial t}\left(\phi\sum_j \rho_j U_j S_j + (1-\phi)\rho_r c_r T\right) + \nabla\cdot\left(\sum_j \rho_j H_j \mathbf{u}_j\right) - \nabla\cdot(\lambda_{eff}\nabla T) = q_H \pm \Delta H_{dep},
\]

где \(\Delta H_{dep}\) --- тепловой эффект фазового перехода при осаждении твёрдой фазы, \(\lambda_{eff}\) --- эффективная теплопроводность насыщенной пористой среды. Дискретизация по времени полностью неявных систем такого рода приводит к большим разреженным нелинейным алгебраическим системам, решаемым методом Ньютона с последующим решением линеаризованной системы GMRES-подобными krylov-методами с многоуровневым (constrained pressure residual, CPR) предобуславливанием --- это стандартный подход в промышленных композиционных симуляторах для задач такого масштаба.

Метод конечных элементов (FEM) реже применяется для чисто адвективно-доминирующих задач переноса компонентов из-за проблем с локальным сохранением массы и осцилляциями на резких фронтах, но остаётся востребованным для механически- и термически-сопряжённых задач (термоупругость породы, растрескивание при тепловом воздействии), а также в гибридных постановках смешанных конечных элементов (mixed finite element) для точного расчёта потока Дарси в сочетании с разрывным методом Галёркина (Discontinuous Galerkin) для адвекции компонентов, что обеспечивает как локальное сохранение массы, так и высокий порядок точности на негладких решениях.

Многомасштабные конечно-объёмные методы (multiscale finite volume, MSFV) были развиты специально для композиционных задач большой размерности, где явное разрешение всех гетерогенностей проницаемости на мелкой сетке вычислительно неподъёмно: локальные базисные функции строятся на грубой сетке и корректируются итерационно (алгоритм многомасштабного восстановления невязки), что позволяет достигать точности мелкосеточного решения при существенно меньших вычислительных затратах для композиционного многофазного потока (Hajibeygi \& Tchelepi, 2013). Для явного сопряжения термического и композиционного переноса при интенсивных тепловых методах разработана единая операторно-ориентированная (Operator-Based Linearization, OBL) архитектура симулятора, позволяющая гибко комбинировать термические и композиционные степени свободы без необходимости заново выводить якобиан при каждой смене физической модели, что особенно полезно при добавлении в общую систему уравнений дополнительного уравнения кинетики осаждения тяжёлых компонентов (Khait \& Voskov, 2019).

\subsection{6.2. Композиционное моделирование многофазных многокомпонентных систем: детальное и псевдокомпонентное представление тяжёлых фракций}\label{ux43aux43eux43cux43fux43eux437ux438ux446ux438ux43eux43dux43dux43eux435-ux43cux43eux434ux435ux43bux438ux440ux43eux432ux430ux43dux438ux435-ux43cux43dux43eux433ux43eux444ux430ux437ux43dux44bux445-ux43cux43dux43eux433ux43eux43aux43eux43cux43fux43eux43dux435ux43dux442ux43dux44bux445-ux441ux438ux441ux442ux435ux43c-ux434ux435ux442ux430ux43bux44cux43dux43eux435-ux438-ux43fux441ux435ux432ux434ux43eux43aux43eux43cux43fux43eux43dux435ux43dux442ux43dux43eux435-ux43fux440ux435ux434ux441ux442ux430ux432ux43bux435ux43dux438ux435-ux442ux44fux436ux451ux43bux44bux445-ux444ux440ux430ux43aux446ux438ux439}

Полное («детальное») представление нефтяной смеси в композиционном симуляторе потребовало бы решения уравнения состояния для десятков и сотен индивидуальных компонентов, что практически невозможно как из-за отсутствия экспериментальных данных для идентификации каждого компонента тяжёлого остатка, так и из-за неприемлемой вычислительной стоимости (число уравнений фазового равновесия растёт линейно с числом компонентов, а число итераций стабильности --- существенно быстрее). Поэтому стандартной практикой является псевдокомпонентное (lumped) представление: экспериментально измеренная кривая истинных точек кипения (True Boiling Point) или данные хромато-масс-спектрометрии группируются в ограниченное число (обычно 5--15) псевдокомпонентов, каждому из которых приписываются эффективные критические параметры (\(T_c\), \(P_c\), \(\omega\)) на основе корреляций (Ли--Кеслера, Твю и др.) и параметр бинарного взаимодействия с остальными компонентами. Основополагающая методология агрегации (lumping) компонентов сырой нефти для целей композиционного моделирования, включая выбор числа и границ псевдокомпонентов для минимизации потери точности фазовой диаграммы, была систематически изложена ещё в начале 1980-х годов и остаётся базой большинства современных схем характеризации (Hong, 1982).

Для тяжёлого «хвоста» (heptane-plus фракции, включающей смолы и асфальтены) псевдокомпонентное представление сталкивается с дополнительной сложностью: именно эта фракция ответственна за осаждение твёрдой фазы, а её термодинамическое поведение (в частности, положение AOP) крайне чувствительно к выбору числа под-фракций внутри C7+ и параметров уравнения состояния. Отсюда --- тенденция выделять асфальтеновую составляющую в отдельный, специально параметризованный псевдокомпонент (или несколько псевдокомпонентов разной молекулярной массы для отражения полидисперсности), термодинамические параметры которого настраиваются отдельно по данным осаждения \emph{н}-алканами, а не по общей корреляции для C7+. Именно такой подход --- трёхфазная многокомпонентная композиционная модель для осаждения асфальтенов на основе CPA-EOS с явным выделением асфальтеновой псевдокомпоненты и явным учётом равновесия между водной, нефтяной и асфальтеновой твёрдой/жидкой фазами --- был реализован для месторождений с высоким содержанием асфальтенов (Nasrabadi, Moortgat \& Firoozabadi, 2013). Аналогичная по духу, но методологически отдельная реализация ввела в открытый симулятор явный термодинамический блок расчёта осаждения асфальтенов на основе теории регулярных растворов, сопряжённый с многофазным композиционным потоком в трёхмерной гетерогенной модели пласта (Gonzalez, Barrufet \& Nasrabadi, 2014).

Сравнение детального и псевдокомпонентного подходов, а также сопоставление разных термодинамических моделей осаждения асфальтенов (модель твёрдого раствора, CPA-EOS) в рамках трёх независимых промышленных и исследовательских симуляторов (UTCOMP, CMG/GEM, ECLIPSE) на едином наборе входных PVT-данных показало существенные расхождения в предсказанном положении верхней и нижней границ AOP и в абсолютном количестве осаждённой твёрдой фазы между симуляторами при формально одинаковых входных данных, что подчёркивает сильную зависимость результата не только от термодинамической модели, но и от деталей численной реализации фазового равновесия и итерационной схемы стабильности (Al-Qasim \& AlDawsari, 2017). Это системная проблема отрасли: сопоставление чисто композиционного и чёрно-нефтяного (black-oil) подходов к описанию многофазного потока показывает, что чёрно-нефтяная аппроксимация, не различающая индивидуальные компоненты, в принципе неприменима к задачам осаждения асфальтенов/парафинов, так как сама природа явления требует детального знания локального состава нефти, и лишь композиционный подход обеспечивает физически корректную основу (Coats et al., 1995).

Для описания необратимого фазового перехода твёрдой фазы (в отличие от обратимого равновесия жидкость--жидкость/пар) вводится дополнительная неизвестная --- концентрация осевшего компонента \(\sigma\), для которой в общую систему уравнений сохранения массы добавляется кинетическое уравнение источника (см. раздел 5.4). Такая расширенная формулировка была одной из первых, где явно связаны композиционное фазовое равновесие асфальтенов и локальное изменение проницаемости внутри стандартного трёхфазного (нефть--вода--газ) композиционного симулятора, что позволило воспроизвести не только распределение осаждённой фазы, но и обратное влияние на профиль давления и дебит скважины (Kohse \& Nghiem, 2004). Расширение той же логики на связанную систему «скважина--пласт» с явным разрешением как радиального профиля в стволе, так и трёхмерного распределения в ближней зоне пласта потребовало дополнительной итерационной связки между гидравлической моделью многофазного течения в трубе и композиционным симулятором пласта на каждом шаге по времени (Darabi et al., 2014).

Для парафиновых отложений композиционная связка чаще реализуется через более простую, чем для асфальтенов, термодинамику (модель твёрдого раствора для \emph{н}-парафинов на основе идеальной или регулярной жидкой фазы и модели множественной твёрдой фазы для описания политипии кристаллов парафина), но с существенно более выраженной ролью молекулярной диффузии как основного механизма переноса к охлаждённой стенке трубы/скважины. Классическая математическая модель отложения парафина в трубопроводе, связывающая радиальный градиент температуры, локальное равновесие растворимости твёрдого парафина в нефти по закону Аррениуса-подобной зависимости от температуры и диффузионный поток парафина к стенке (закон Фика с коэффициентом диффузии, зависящим от температуры по корреляции Вилке--Чанга), заложила основу для большинства последующих инженерных моделей толщины отложений (Svendsen, 1993). Численная реализация аналогичной физики применительно к пласту (а не трубопроводу) --- с учётом неизотермической фильтрации при закачке холодной воды в добывающую скважину, добывающую парафинистую нефть, --- потребовала решения связанной системы уравнений двухфазной фильтрации, конвективно-кондуктивного теплопереноса и локального равновесия растворимости парафина методом конечных разностей на прямоугольной сетке, что позволило воспроизвести пространственно-временную динамику зоны кольматации вблизи нагнетательной скважины (Nie \& Yang, 2014).

\subsection{6.3. Сравнительный обзор симуляторов}\label{ux441ux440ux430ux432ux43dux438ux442ux435ux43bux44cux43dux44bux439-ux43eux431ux437ux43eux440-ux441ux438ux43cux443ux43bux44fux442ux43eux440ux43eux432}

Ниже приведена сравнительная таблица основных исследовательских и промышленных инструментов, применяемых для моделирования осаждения тяжёлых компонентов нефти (асфальтенов, парафинов) при фильтрации и течении в стволе скважины/трубопроводе.

{\def\LTcaptype{none} % do not increment counter
\begin{xltabular}{\linewidth}{@{}
  >{\raggedright\arraybackslash}X
  >{\raggedright\arraybackslash}X
  >{\raggedright\arraybackslash}X
  >{\raggedright\arraybackslash}X
  >{\raggedright\arraybackslash}X@{}}
\toprule\noalign{}
Симулятор / модель & Тип термодинамической модели & Что учитывает & Область применения & Ограничения \\
\midrule\noalign{}
\endhead
\bottomrule\noalign{}
\endlastfoot
\textbf{ADEPT} (Asphaltene Deposition Tool, Rice Univ./Chapman--Vargas) & Модель твёрдого раствора / регулярного раствора для доли выпавшей фазы + отдельная кинетика захвата на стенке & 1D перенос вдоль скважины/трубы, массообмен через число Шервуда, калибровка по промысловым данным (Kurup et al., 2012) & Прогноз профиля отложений в стволе скважины и НКТ & Одномерность, эмпирические корреляции коэффициента массообмена, не учитывает детальную гидродинамику \\
\textbf{MAD-ADEPT} (многофазное динамическое расширение ADEPT) & То же + явный учёт многофазного (нефть--газ) потока & Динамика осаждения при переменном режиме потока, ламинарный/турбулентный перенос частиц (Naseri, Jamshidi \& Taghikhani, 2020) & Стволы скважин с газлифтом и переменной обводнённостью & Требует калибровки эмпирических кинетических констант под конкретное месторождение \\
\textbf{UTCOMP} (Center for Petroleum and Geosystems Engineering, UT Austin) & Модель твёрдого раствора / CPA-EOS (в расширенных версиях) & Полностью композиционный многофазный поток в пласте с опцией осаждения асфальтенов (Al-Qasim \& AlDawsari, 2017) & Исследовательское композиционное моделирование пласта & Открытый исследовательский код, ограниченная промышленная поддержка и GUI \\
\textbf{CMG/GEM} (Computer Modelling Group) & Модель твёрдого раствора / кубическое EOS с настраиваемым осаждением & Полный композиционный симулятор пласта промышленного уровня, включая опцию осаждения асфальтенов и кольматации (Al-Qasim \& AlDawsari, 2017) & Промышленное пластовое моделирование (в т.ч. закачка газа, WAG) & Закрытый код; ограниченная гибкость встраивания нестандартных термодинамических моделей (напр. PC-SAFT) без API \\
\textbf{Schlumberger ECLIPSE (E300/E100 + PVTi)} & Модель твёрдого раствора для осаждения асфальтенов & Промышленный композиционный/чёрно-нефтяной симулятор с модулями PVT и опцией осаждения асфальтенов (Al-Qasim \& AlDawsari, 2017) & Отраслевой стандарт пластового моделирования & Аналогично GEM --- закрытая архитектура; расхождения с другими симуляторами при одинаковых входных PVT-данных \\
\textbf{MRST Compositional + PC-SAFT} (открытая академическая платформа, SINTEF/UiB) & PC-SAFT (полностью встроенный расчёт фазового равновесия) & Прямое сопряжение PC-SAFT с многофазным потоком в открытом коде на MATLAB (Masoudi et al., 2021) & Исследовательские задачи, тестирование новых термодинамических моделей & Ограниченная производительность на промысловых по масштабу сетках, MATLAB-based \\
\textbf{CPA-EOS композиционная модель} (Nasrabadi et al.) & Cubic-Plus-Association EOS, трёхфазная (нефть--вода--асфальтен) & Явный учёт водородного связывания асфальтенов с водной фазой, трёхфазное равновесие (Nasrabadi, Moortgat \& Firoozabadi, 2013) & Исследование влияния заводнения/закачки воды на осаждение асфальтенов & Исследовательская реализация, не промышленный продукт \\
\textbf{Симулятор Kohse--Nghiem (CMG-based research prototype)} & Модель твёрдого раствора + явная кинетика осаждения/кольматации & Сопряжение композиционного равновесия, кинетики осаждения и изменения проницаемости (Kohse \& Nghiem, 2004) & Пласт, скважина & Упрощённая (не мульти-масштабная) кинетика захвата частиц \\
\textbf{DARTS} (Delft Advanced Research Terra Simulator, Operator-Based Linearization) & Гибкая архитектура --- совместима с произвольной EOS (кубической, PC-SAFT) через операторную форму & Совместное термическое и композиционное моделирование многофазного потока, гибкое добавление физики (в т.ч. кинетики осаждения) без переписывания якобиана (Khait \& Voskov, 2019) & Исследовательские задачи термических и комбинированных процессов & Относительно новая платформа, ограниченная промышленная валидация для асфальтеновых/парафиновых задач \\
\textbf{Многомасштабный композиционный конечно-объёмный метод (MSFV)} & Совместим с любой EOS нижнего уровня (кубической, ассоциативной) & Ускоренное композиционное моделирование на сильно гетерогенных по проницаемости пластах (Hajibeygi \& Tchelepi, 2013) & Крупномасштабные модели пласта с детальной геологической сеткой & Требует дополнительной итерационной коррекции для сохранения точности при сильных гетерогенностях \\
\textbf{Модель Svendsen (парафин, трубопровод)} & Термодинамика твёрдого раствора \emph{н}-парафинов + диффузионный перенос по Фику & Радиальный градиент температуры, молекулярная диффузия к стенке, рост слоя отложений (Svendsen, 1993) & Трубопроводы, стволы скважин & Не учитывает сдвиговое дисперсионное осаждение и последующую эрозию слоя \\
\textbf{Модель Nie--Yang (парафин, пласт, закачка холодной воды)} & Локальное равновесие растворимости парафина + конечно-разностная схема & Неизотермическая двухфазная фильтрация с кольматацией парафином вблизи нагнетательной скважины (Nie \& Yang, 2014) & Заводнение месторождений с высоким содержанием парафина & Конечно-разностная схема на регулярной сетке, ограниченная геометрия \\
\end{xltabular}
}

Сопоставление приведённых инструментов показывает общую тенденцию: узкоспециализированные модели типа ADEPT и MAD-ADEPT обеспечивают быстрый инженерный расчёт профиля отложений в стволе скважины ценой одномерности и эмпиричности кинетических констант, тогда как промышленные композиционные симуляторы (GEM, ECLIPSE, UTCOMP) обеспечивают полную трёхмерную многофазную связку с пластом, но заметно расходятся между собой в количественном прогнозе AOP и массы осадка на идентичных входных данных (Al-Qasim \& AlDawsari, 2017) --- это прямое следствие различий в реализации термодинамики твёрдой фазы и алгоритмов проверки фазовой стабильности. Открытые исследовательские платформы (MRST, DARTS) занимают промежуточную нишу, позволяя гибко тестировать новые термодинамические модели (PC-SAFT, CPA) в связке с многофазным потоком, но пока уступают промышленным пакетам по производительности на сетках промыслового масштаба.

\subsection{6.4. Выводы и открытые проблемы}\label{ux432ux44bux432ux43eux434ux44b-ux438-ux43eux442ux43aux440ux44bux442ux44bux435-ux43fux440ux43eux431ux43bux435ux43cux44b}

Анализ современного состояния численного моделирования показывает, что наиболее слабым звеном остаётся не столько дискретизация уравнений переноса (для которой существует зрелый арсенал конечно-объёмных и конечно-элементных схем), сколько согласованность между термодинамической моделью фазового поведения тяжёлых компонентов (SAFT/PC-SAFT, CPA, модель твёрдого раствора) и кинетической моделью захвата/выноса частиц на масштабе пор, которая в большинстве работ остаётся полуэмпирической и требует калибровки под конкретный тип нефти и коллектора. Комплексный обзор состояния проблемы осаждения асфальтенов в пластах прямо указывает на нехватку универсальных, физически обоснованных корреляций для кинетических констант захвата/выноса, переносимых между месторождениями без повторной калибровки, как на одну из ключевых нерешённых задач области (Alimohammadi, Zendehboudi \& James, 2019). Дальнейшее развитие в этом направлении связано с интеграцией более точных наноуровневых и порово-масштабных (pore-scale) моделей захвата частиц (на основе прямого численного моделирования течения в реалистичной геометрии пор, например по данным микро-КТ) с макроскопическими композиционными симуляторами через методы апскейлинга, а также с расширением операторно-ориентированных архитектур симуляторов (DARTS и аналоги) для бесшовного встраивания произвольных термодинамических моделей асфальтенов без необходимости переписывания решателя.

\newpage

\section{Раздел 7. Экспериментальная валидация и промысловые кейсы}\label{ux440ux430ux437ux434ux435ux43b-7.-ux44dux43aux441ux43fux435ux440ux438ux43cux435ux43dux442ux430ux43bux44cux43dux430ux44f-ux432ux430ux43bux438ux434ux430ux446ux438ux44f-ux438-ux43fux440ux43eux43cux44bux441ux43bux43eux432ux44bux435-ux43aux435ux439ux441ux44b}

\subsection{7.1. Введение}\label{ux432ux432ux435ux434ux435ux43dux438ux435-1}

Любая математическая модель осаждения парафинов, асфальтенов и смол (ПАС) в конечном счёте оценивается по её способности воспроизводить наблюдаемые лабораторные и промысловые данные. Поскольку процесс осаждения тяжёлых компонентов нефти при неизотермической фильтрации управляется совокупностью термодинамических (фазовое равновесие твёрдое--жидкость), кинетических (скорость кристаллизации/агрегации) и гидродинамических (диффузия, сдвиговая дисперсия, старение отложений) факторов, экспериментальная база данных играет двойную роль: она поставляет входные параметры моделей (температура начала кристаллизации парафина WAT, давление начала осаждения асфальтенов АОП, компонентный состав по SARA) и одновременно служит эталоном для проверки прогнозов. В данном разделе систематизированы основные экспериментальные методики, рассмотрены опубликованные сопоставления «модель--эксперимент» и «модель--промысел», а также проанализированы типичные источники расхождений.

\subsection{7.2. Экспериментальные методики определения WAT, SARA-состава и кинетики отложения}\label{ux44dux43aux441ux43fux435ux440ux438ux43cux435ux43dux442ux430ux43bux44cux43dux44bux435-ux43cux435ux442ux43eux434ux438ux43aux438-ux43eux43fux440ux435ux434ux435ux43bux435ux43dux438ux44f-wat-sara-ux441ux43eux441ux442ux430ux432ux430-ux438-ux43aux438ux43dux435ux442ux438ux43aux438-ux43eux442ux43bux43eux436ux435ux43dux438ux44f}

\subsubsection{7.2.1. Определение температуры начала кристаллизации парафина (WAT/WDT)}\label{ux43eux43fux440ux435ux434ux435ux43bux435ux43dux438ux435-ux442ux435ux43cux43fux435ux440ux430ux442ux443ux440ux44b-ux43dux430ux447ux430ux43bux430-ux43aux440ux438ux441ux442ux430ux43bux43bux438ux437ux430ux446ux438ux438-ux43fux430ux440ux430ux444ux438ux43dux430-watwdt}

Температура появления твёрдой фазы парафинов (Wax Appearance Temperature, WAT) и обратная величина --- температура полного растворения (Wax Disappearance Temperature, WDT) --- являются ключевыми калибровочными параметрами для термодинамических блоков моделей осаждения. Помимо классических визуальных и калориметрических методов (дифференциальная сканирующая калориметрия, ДСК), для количественной оценки доли твёрдой фазы парафина всё активнее используется ядерный магнитный резонанс (ЯМР). Новые методики ЯМР позволяют оценивать эволюцию осаждения парафина в системах сырой нефти количественно и с высоким временным разрешением (Energy \& Fuels, DOI 10.1021/acs.energyfuels.2c03309), что даёт кинетическую кривую доли выпавшей твёрдой фазы, пригодную для прямой калибровки кинетических констант в моделях типа Бурже или модели молекулярной диффузии Хансена. Протонный ЯМР также применяется для прогноза содержания парафинов в нефтях и нефтяных фракциях как более быстрая альтернатива гравиметрическим методам (Petroleum Science and Technology), а низкополевая ЯМР-релаксометрия --- для комплексной характеризации состава сырой нефти без предварительной пробоподготовки (Magn. Reson. Chem., DOI 10.1002/mrc.70096).

\subsubsection{7.2.2. SARA-анализ (Saturates--Aromatics--Resins--Asphaltenes)}\label{sara-ux430ux43dux430ux43bux438ux437-saturatesaromaticsresinsasphaltenes}

Компонентный состав нефти по фракциям насыщенных, ароматических углеводородов, смол и асфальтенов (SARA) остаётся базовым инструментом характеризации сырья для калибровки термодинамических моделей осаждения (как парафинового, так и асфальтенового блоков), поскольку соотношение смол и асфальтенов напрямую определяет коллоидную устойчивость асфальтеновых наноагрегатов. Классический метод жидкостной хроматографии по ASTM/IP при этом обладает известными недостатками --- большой длительностью анализа, сложностью системы и трудностями с воспроизводимостью (Bissada et al., ACS Omega, DOI 10.1021/acsomega.5c00439). Для их преодоления разработаны автоматизированные схемы на основе высокоэффективной жидкостной хроматографии (HPLC) с раздельным сбором фракций. В частности, метод μSARA-HPLC использует последовательность из трёх колонок: колонку Polaris Si-SA MetaGuard для осаждения асфальтенов, колонку ZORBAX SB-CN для выделения смол и колонку ZORBAX RX-SIL для ароматических соединений, при этом границы элюирования фракций калибруются с помощью модельных соединений и верифицируются газовой хроматографией с масс-спектрометрией (ГХ-МС) для насыщенных и ароматических соединений и масс-спектрометрией ионно-циклотронного резонанса с преобразованием Фурье (ФТ-ИЦР-МС) для смол и асфальтенов (там же). Ключевое эксплуатационное преимущество --- существенное сокращение времени анализа: система работает при расходе 1.5 мл/мин с использованием н-пентана, дихлорметана, метанола и изопропанола, каждый прогон завершается за 35 минут, что значительно быстрее классических методов (там же). Аналитические характеристики метода также существенны для оценки погрешности входных данных моделей: двойное детектирование с испарительным детектором светорассеяния (ELSD) и диодно-матричным детектором (DAD) обеспечивает исключительную чувствительность для всех фракций, достигая пределов обнаружения до 0.05 масс.\% и среднего показателя извлечения 99.52\% (там же). Воспроизводимость проверена на широком спектре нефтей и битумов: она подтверждена на нефтях и битумах с различной плотностью API, включая сверхтяжёлые и сверхлёгкие образцы, с использованием эталонной нефти NSO-1 на протяжении четырёх месяцев, что привело к отклонению менее 1\% (там же). Это задаёт практический ориентир: даже при использовании наиболее совершенных лабораторных протоколов неопределённость содержания асфальтенов/смол в составе входных данных модели не опускается ниже долей процента, а на менее стандартизированных лабораторных базах расхождения между лабораториями по классическим ASTM-методам могут составлять несколько процентов по массе, что напрямую транслируется в сдвиг расчётного давления начала осаждения асфальтенов (АОП) на десятки атмосфер при использовании термодинамических моделей на основе теории регулярных растворов или уравнений состояния с асфальтеновой псевдокомпонентой.

\subsubsection{7.2.3. Лабораторные методики оценки скорости отложения: cold finger и flow loop}\label{ux43bux430ux431ux43eux440ux430ux442ux43eux440ux43dux44bux435-ux43cux435ux442ux43eux434ux438ux43aux438-ux43eux446ux435ux43dux43aux438-ux441ux43aux43eux440ux43eux441ux442ux438-ux43eux442ux43bux43eux436ux435ux43dux438ux44f-cold-finger-ux438-flow-loop}

Для количественной оценки кинетики роста отложений парафина на охлаждаемой поверхности в промысловой практике и в академических исследованиях исторически используются два взаимодополняющих типа установок.

\textbf{Метод «холодного пальца» (cold finger test)} представляет собой лабораторную ячейку, в которой охлаждаемый стержень погружён в перемешиваемый объём нефти, поддерживаемой при температуре выше WAT; парафин осаждается на холодной поверхности за счёт градиента температуры и соответствующего градиента растворимости. Метод широко применяется для быстрого скрининга ингибиторов депарафинизации и для оценки относительной скорости отложения при различных термических режимах и режимах сдвига.

\textbf{Метод петли потока (flow loop)} воспроизводит реальные гидродинамические условия трубопровода --- турбулентное течение, сдвиговые напряжения на стенке, конвективный теплообмен --- и считается более представительным для масштабирования на промысловые условия, хотя и требует значительно больших объёмов образца и более сложного оборудования. Сопоставительные исследования показывают, что результаты этих двух методов не всегда совпадают количественно из-за различий в гидродинамике и локальном сдвиге на охлаждаемой стенке, что делает выбор корректного метода масштабирования (up-scaling) от лабораторной ячейки к трубопроводу диаметром в десятки сантиметров одной из ключевых нерешённых практических задач валидации.

\subsection{7.3. Сопоставление модельных прогнозов с лабораторными и промысловыми данными}\label{ux441ux43eux43fux43eux441ux442ux430ux432ux43bux435ux43dux438ux435-ux43cux43eux434ux435ux43bux44cux43dux44bux445-ux43fux440ux43eux433ux43dux43eux437ux43eux432-ux441-ux43bux430ux431ux43eux440ux430ux442ux43eux440ux43dux44bux43cux438-ux438-ux43fux440ux43eux43cux44bux441ux43bux43eux432ux44bux43cux438-ux434ux430ux43dux43dux44bux43cux438}

\subsubsection{7.3.1. Валидация на масштабе трубопровода (парафиновые отложения)}\label{ux432ux430ux43bux438ux434ux430ux446ux438ux44f-ux43dux430-ux43cux430ux441ux448ux442ux430ux431ux435-ux442ux440ux443ux431ux43eux43fux440ux43eux432ux43eux434ux430-ux43fux430ux440ux430ux444ux438ux43dux43eux432ux44bux435-ux43eux442ux43bux43eux436ux435ux43dux438ux44f}

Для случая осаждения парафина в трубопроводах наиболее полные валидационные исследования опираются на комбинацию лабораторных петель потока, промысловых замеров толщины отложений (например, по данным очистки скребком/поршнем --- «пиггования») и математических моделей молекулярной диффузии. Модель, объединяющая одновременно эмульсию вода-в-нефти, кинетику осаждения парафина, молекулярную диффузию и сдвиговую дисперсию, была верифицирована на данных полевого масштаба (Ochieng et al., J. Applied Mathematics, 2022, DOI 10.1155/2022/2845221): формирование твёрдых кристаллов парафина, которые смыкаются и образуют гелеобразный слой на внутренней стенке трубопровода сырой нефти, влияет на транспортировку парафинистой сырой нефти, при этом целью работы было разработать математическую модель, которая включает эмульсии вода-в-нефти, кинетику осаждения парафина, молекулярную диффузию и сдвиговую дисперсию для точного прогнозирования как скорости роста отложений парафина, так и старения отложений. Численно задача решалась методом связанных нелинейных дифференциальных уравнений в частных производных, дискретизированных по времени полунеявной схемой второго порядка на основе методов Адамса-Башфорта и Крэнка-Николсон, а по пространству --- бивариантным спектральным коллокационным методом на основе узлов Чебышева-Гаусса-Лобатто, реализованным в MATLAB. Результаты моделирования согласуются с ожидаемыми качественными закономерностями: толщина отложений прямо пропорциональна числу Рейнольдса и обратно пропорциональна числу Грасгофа по массе, числу Шмидта и числу Вебера, а старение отложений интенсивно на ранних стадиях осаждения парафина, после чего стабилизируется на определённом значении с течением времени.

Отдельного внимания заслуживает эмпирическая модель, разработанная и провалидированная непосредственно на промысловых данных (Sustainability, 2023, 15, 9616, DOI 10.3390/su15129616): создана новая эмпирическая модель формирования отложений, основанная на результатах массива лабораторных исследований, теоретических данных и технологических расчётов, учитывающая условия течения нефти, данные лабораторных исследований и обводнённость продукции. Промышленная апробация показала высокую степень соответствия: опыт промышленной эксплуатации методики показал высокую сходимость расчётного и фактического профилей отложений, причём количественно на основе сравнения расчётного и фактического профилей отложений сделан вывод, что среднеквадратичное отклонение по максимальной толщине отложений парафина составляет 6.0\%, а по глубине с наибольшей толщиной парафина --- 3.5\%. Такой уровень точности (единицы процентов по толщине и по глубине залегания максимума) представляет собой ориентир, к которому стремятся более детализированные композиционные модели, и одновременно демонстрирует, что даже относительно простая эмпирическая модель, откалиброванная на большом массиве промысловых данных одного месторождения, может обеспечивать приемлемую для инженерных решений (планирование скребковой очистки, тепловой изоляции, глубины греющего кабеля) точность.

\subsubsection{7.3.2. Валидация на масштабе скважины (асфальтеновые отложения)}\label{ux432ux430ux43bux438ux434ux430ux446ux438ux44f-ux43dux430-ux43cux430ux441ux448ux442ux430ux431ux435-ux441ux43aux432ux430ux436ux438ux43dux44b-ux430ux441ux444ux430ux43bux44cux442ux435ux43dux43eux432ux44bux435-ux43eux442ux43bux43eux436ux435ux43dux438ux44f}

Для асфальтеновых отложений в стволе скважины классическим и часто цитируемым эталонным набором промысловых данных остаются кейсы месторождения Marrat (Кувейт), где были задокументированы падение забойного и устьевого давления, а также кавернометрия (каверномер/шаблонирование труб), фиксирующая профиль сужения проходного сечения НКТ по глубине из-за асфальтеновых отложений. В сравнительном исследовании нескольких термодинамических и кинетических моделей осаждения асфальтенов на этом кейсе было показано, что различные модели (в том числе модели Фридлендера-Джонстона, Била, Эскобедо-Мансури, Кливера-Йейтса и Джамиалахмади) дают качественно схожий, но количественно расходящийся профиль отложений по глубине скважины, причём расхождение между моделями и промысловыми данными каверномера растёт по мере приближения к забою скважины, где велики отклонения давления от давления начала осаждения асфальтенов (АОП) и наиболее выражены эффекты турбулентного массопереноса.

Повторный анализ того же промыслового кейса (Kargarpour, Dandekar, J. Petrol. Explor. Prod. Technol., 2016, DOI 10.1007/s13202-015-0221-7) предложил принципиально иную интерпретацию наблюдаемого падения устьевого давления в скважинах месторождения Marrat. Авторы указывают, что снижение забойного и устьевого давления традиционно интерпретировалось напрямую как следствие роста слоя асфальтеновых отложений на стенке НКТ, однако, проведя аналогию с явлением скопления жидкости в стволе газовой скважины (liquid loading) и с типичной кривой производительности НКТ, показали, что часть наблюдаемого падения давления объясняется гидравликой двухфазного потока (эффект скольжения фаз, слип-скорость) в присутствии асфальтеновых микрочастиц, оставшихся в стволе, а не только уменьшением проходного сечения трубы. В рамках этой переинтерпретации процесс отложения асфальтена в НКТ описывается как динамический процесс «осаждения и среза» («seating and carving»), при котором по достижении критической толщины слоя усиливается механический срез («cutting») отложений потоком, что и обусловливает дальнейшее падение давления. Такой пример показывает, что расхождение между промысловыми наблюдениями и модельным профилем толщины отложений может частично объясняться не ошибкой кинетической модели осаждения асфальтена как таковой, а неполнотой описания сопутствующей многофазной гидродинамики потока в стволе скважины, используемой при интерпретации косвенных промысловых данных (динамика давления) в терминах толщины твёрдого слоя.

\subsubsection{7.3.3. Полевые кейсы: осложняющие факторы}\label{ux43fux43eux43bux435ux432ux44bux435-ux43aux435ux439ux441ux44b-ux43eux441ux43bux43eux436ux43dux44fux44eux449ux438ux435-ux444ux430ux43aux442ux43eux440ux44b}

Помимо чисто модельных расхождений, промысловые данные вносят собственную неопределённость: реальный дебит, давление и температурный профиль по стволу скважины меняются во времени вследствие истощения пласта, обводнения продукции и изменения газового фактора, тогда как большинство детальных композиционных моделей осаждения решаются в квазистационарной или медленно-нестационарной постановке с фиксированным профилем PVT-параметров. Кроме того, каверномер и другие геофизические методы оценки толщины отложений в стволе скважины сами обладают инструментальной погрешностью и не позволяют напрямую разделить вклад асфальтеновых, парафиновых и неорганических (солевых) отложений там, где они сосуществуют, что вносит систематическую неопределённость в референсные данные, используемые для валидации.

\subsection{7.4. Типичные расхождения и источники неопределённости}\label{ux442ux438ux43fux438ux447ux43dux44bux435-ux440ux430ux441ux445ux43eux436ux434ux435ux43dux438ux44f-ux438-ux438ux441ux442ux43eux447ux43dux438ux43aux438-ux43dux435ux43eux43fux440ux435ux434ux435ux43bux451ux43dux43dux43eux441ux442ux438}

Систематизируя опубликованные сопоставления модель--эксперимент/промысел, можно выделить несколько устойчиво повторяющихся категорий источников расхождения.

\textbf{1. Погрешность характеризации компонентного состава.} Как показано выше, даже наиболее точные современные протоколы SARA-анализа обеспечивают предел обнаружения порядка \(0.05\%\) масс. и воспроизводимость на уровне \textless1\% при многократных измерениях эталонной нефти в одной лаборатории, однако межлабораторная воспроизводимость классических ASTM-методов существенно хуже. Поскольку термодинамические модели осаждения асфальтенов (модель регулярных растворов Флори-Хаггинса, кубические уравнения состояния с асфальтеновой псевдофазой, модель PC-SAFT) чрезвычайно чувствительны к доле и молекулярной массе асфальтеновой фракции, ошибка входных данных в единицы процентов транслируется в сдвиг расчётного давления начала осаждения (АОП) на существенную величину, что напрямую влияет на предсказанную глубину начала отложений в стволе скважины.

\textbf{2. Неопределённость кинетических параметров.} В отличие от термодинамического равновесия (WAT, АОП), кинетические константы скорости кристаллизации парафина, коэффициенты молекулярной диффузии в модели Хансена-Бурже, а также параметры моделей прилипания/эрозии асфальтеновых частиц на стенке (используемые, например, в модели Ескобедо-Мансури) обычно определяются путём обратной калибровки (history matching) на ограниченном наборе экспериментальных кривых накопления массы отложений, что делает их плохо экстраполируемыми на условия, отличные от калибровочных (другой диапазон чисел Рейнольдса, температурного градиента, обводнённости). Именно этим объясняется наблюдаемое в цитированных работах систематическое расхождение между моделями при их применении к одному и тому же промысловому кейсу.

\textbf{3. Различие лабораторных методик и масштабный эффект (scale-up).} Как отмечено выше, cold finger и flow loop дают не всегда согласующиеся количественные оценки скорости отложения парафина из-за различий в локальной гидродинамике и сдвиговых напряжениях на охлаждаемой поверхности, а перенос лабораторных корреляций на трубопровод промыслового диаметра требует введения дополнительных эмпирических поправочных коэффициентов, физическая обоснованность которых ограничена.

\textbf{4. Нестационарность и многофазность реальных промысловых условий.} Изменение дебита, обводнённости и газового фактора во времени, а также многофазный характер потока (нефть-газ-вода) в значительной части реальных скважин и трубопроводов создают условия, для которых большинство однофазных или квазистационарных моделей требуют дополнительной валидации; в частности, для многофазных систем взаимодействие капель воды с растущим слоем отложений и изменение эффективного сдвигового напряжения на стенке --- предмет отдельного моделирования, включённого в упомянутую выше модель эмульсии вода-в-нефти.

\textbf{5. Неполнота учёта сопутствующих механизмов.} Ряд обзорных работ отмечает, что классические механизмы (молекулярная диффузия и «старение» отложений за счёт диффузии более тяжёлых парафинов вглубь геля) не всегда достаточны для объяснения наблюдаемой кинетики: молекулярная диффузия и старение традиционно считаются основными механизмами осаждения парафина, однако всё больше исследований доказывают, что эти механизмы по отдельности не могут полностью объяснить явление осаждения парафина (Energy \& Fuels, DOI 10.1021/acs.energyfuels.3c04687), что стимулирует включение дополнительных механизмов (сдвиговая дисперсия, дисперсия за счёт броуновского движения кристаллов, гравитационное осаждение) в модели нового поколения, обсуждаемые в разделе 8.

\subsection{7.5. Промежуточные выводы}\label{ux43fux440ux43eux43cux435ux436ux443ux442ux43eux447ux43dux44bux435-ux432ux44bux432ux43eux434ux44b}

Экспериментальная валидация моделей осаждения ПАС остаётся неотъемлемой, но методически неоднородной частью исследовательского цикла: лабораторные методы (cold finger, flow loop, ДСК, ЯМР, HPLC-SARA) обеспечивают входные параметры и калибровочные кривые, тогда как промысловые кейсы (профили давления, кавернометрия, данные пиггования) служат конечным критерием практической применимости модели. Опубликованные сопоставления демонстрируют, что при качественной калибровке современные модели способны воспроизводить промысловые профили отложений с отклонением порядка нескольких процентов по ключевым метрикам (максимальная толщина, глубина максимума), однако такая точность достижима только при доступности детальных и надёжных входных данных, что на практике оказывается основным лимитирующим фактором точности прогноза, а не структура самой математической модели.

\newpage

\section{Раздел 8. Современные направления и заключение по трендам}\label{ux440ux430ux437ux434ux435ux43b-8.-ux441ux43eux432ux440ux435ux43cux435ux43dux43dux44bux435-ux43dux430ux43fux440ux430ux432ux43bux435ux43dux438ux44f-ux438-ux437ux430ux43aux43bux44eux447ux435ux43dux438ux435-ux43fux43e-ux442ux440ux435ux43dux434ux430ux43c}

\subsection{8.1. Введение}\label{ux432ux432ux435ux434ux435ux43dux438ux435-2}

За последние 5--7 лет (2019--2026 гг.) направление математического моделирования осаждения тяжёлых компонентов нефти претерпело заметную трансформацию: наряду с развитием классических композиционных и термодинамико-кинетических моделей на первый план вышли методы машинного обучения (ML), а также гибридные физически-информированные подходы (physics-informed machine learning, PIML), объединяющие статистическую мощность data-driven моделей с физическими ограничениями, заложенными в уравнения сохранения массы, импульса и энергии. Ниже систематизированы основные направления развития области и сформулированы открытые проблемы.

\subsection{8.2. Продвинутое композиционное моделирование с детальным химическим составом}\label{ux43fux440ux43eux434ux432ux438ux43dux443ux442ux43eux435-ux43aux43eux43cux43fux43eux437ux438ux446ux438ux43eux43dux43dux43eux435-ux43cux43eux434ux435ux43bux438ux440ux43eux432ux430ux43dux438ux435-ux441-ux434ux435ux442ux430ux43bux44cux43dux44bux43c-ux445ux438ux43cux438ux447ux435ux441ux43aux438ux43c-ux441ux43eux441ux442ux430ux432ux43eux43c}

Традиционные модели осаждения парафина (модель молекулярной диффузии Бурже, модель Хансена) и асфальтена (термодинамика регулярных растворов, модель Флори-Хаггинса-Захариа) исторически оперировали укрупнёнными псевдокомпонентами. Современная тенденция состоит в переходе к явному учёту детального углеводородного состава --- вплоть до индивидуальных парафиновых фракций C1--C30+ и распределения молекулярной массы асфальтенов, получаемого методами высокоразрешающей масс-спектрометрии (ФТ-ИЦР-МС) в сочетании с автоматизированным SARA-фракционированием (см. раздел 7.2.2). Такая детализация состава требуется как входные данные для моделей нового поколения, включая гибридные ML-модели, которые explicitly используют компонентный состав как признаковое пространство, а не только интегральные характеристики (плотность, вязкость, WAT).

Одновременно продолжается развитие механистических композиционных моделей многофазного и многомеханизменного переноса. В частности, модель, объединяющая одновременно кинетику осаждения парафина, молекулярную диффузию, сдвиговую дисперсию и эволюцию эмульсии вода-в-нефти (Ochieng et al., 2022, DOI 10.1155/2022/2845221), а также более поздняя работа той же группы, посвящённая численному исследованию осаждения парафина при многофазном течении с учётом тепло- и массообмена (Ochieng et al., Mathematical Problems in Engineering, 2023, DOI 10.1155/2023/1173505), иллюстрируют тенденцию к объединению в единой системе уравнений термодинамики фазового равновесия, кинетики кристаллизации и полноценной гидродинамики многофазного потока, решаемой методами высокого порядка точности (спектральные коллокационные методы на сетке Чебышева-Гаусса-Лобатто вместо конечно-разностных схем первого порядка, использовавшихся в более ранних работах).

Для асфальтенов аналогичная тенденция проявляется в развитии моделей осаждения частиц из турбулентного потока в стволе скважины и трубопроводах с явным учётом профиля концентрации частиц по сечению трубы, турбулентной диффузии и вероятности прилипания к стенке; расширяются также термодинамические модели, связывающие профиль давления вдоль ствола скважины напрямую с термодинамикой многокомпонентного расширения нефти (PVT-моделирование) для расчёта профиля пересыщения по асфальтену без упрощающего предположения о постоянстве давления начала осаждения по всей длине колонны.

\subsection{8.3. Машинное обучение для прогноза WAT, АОП и скорости отложений}\label{ux43cux430ux448ux438ux43dux43dux43eux435-ux43eux431ux443ux447ux435ux43dux438ux435-ux434ux43bux44f-ux43fux440ux43eux433ux43dux43eux437ux430-wat-ux430ux43eux43f-ux438-ux441ux43aux43eux440ux43eux441ux442ux438-ux43eux442ux43bux43eux436ux435ux43dux438ux439}

\subsubsection{8.3.1. WAT и температура кристаллизации парафина}\label{wat-ux438-ux442ux435ux43cux43fux435ux440ux430ux442ux443ux440ux430-ux43aux440ux438ux441ux442ux430ux43bux43bux438ux437ux430ux446ux438ux438-ux43fux430ux440ux430ux444ux438ux43dux430}

Значительный объём публикаций последних лет посвящён прямому прогнозированию WAT/WDT методами ML на основе легко измеримых характеристик нефти (плотность, вязкость, содержание парафина, SARA-состав) без необходимости трудоёмких лабораторных измерений для каждого нового образца. Опубликованы модели на основе радиально-базисных нейронных сетей, оптимизированных генетическим алгоритмом и алгоритмом искусственной пчелиной колонии (RBFNN-GA/ABC), которые достигают коэффициента детерминации порядка \(R^2 \approx 0.97\) на тестовой выборке (Energy \& Fuels, 2019, DOI 10.1021/acs.energyfuels.9b03296). Более поздние работы (2024--2025 гг.) переходят к ансамблевым методам --- градиентному бустингу и стекинг-ансамблям, а также сравнивают их с классическими корреляциями для прогноза WAT тяжёлых нефтей, отмечая улучшение точности за счёт учёта нелинейных взаимодействий между параметрами состава, которые плохо описываются полиномиальными корреляциями (ScienceDirect, 2025, ``Toward accurate modeling of wax appearance temperature (WAT) using intelligent ensemble techniques''; Fuel, 2024, ``New intelligent models for predicting wax appearance temperature'').

\subsubsection{8.3.2. Скорость и профиль отложения парафина в трубопроводах}\label{ux441ux43aux43eux440ux43eux441ux442ux44c-ux438-ux43fux440ux43eux444ux438ux43bux44c-ux43eux442ux43bux43eux436ux435ux43dux438ux44f-ux43fux430ux440ux430ux444ux438ux43dux430-ux432-ux442ux440ux443ux431ux43eux43fux440ux43eux432ux43eux434ux430ux445}

Отдельное направление --- прогнозирование не только термодинамического порога (WAT), но и кинетической величины --- скорости накопления массы отложений или толщины слоя во времени, что представляет собой более сложную регрессионную задачу с учётом временной зависимости. Показано, что многослойный персептрон, обучаемый алгоритмом Левенберга-Марквардта с байесовской регуляризацией (MLP-LMA/BR), на выборке из 88 образцов достигает коэффициента детерминации \(R^2 = 0.9971\) при прогнозе скорости осаждения парафина (South African Journal of Chemical Engineering, 2021, DOI из ScienceDirect S2405656121000560, ``Predicting wax deposition using robust machine learning techniques''). Комитетные (ансамблевые) модели, комбинирующие несколько интеллектуальных алгоритмов через взвешенное голосование, также применялись для повышения устойчивости прогноза скорости осаждения парафина по сравнению с единичными моделями (ResearchGate, ``Combining of intelligent models through committee machine for estimation of wax deposition''). Для оперативного прогнозирования скорости отложения в промысловых трубопроводах предложена модель на основе нейронной сети Элмана (рекуррентная архитектура с контекстным слоем, учитывающая временную динамику процесса), параметры которой оптимизируются усовершенствованным алгоритмом поиска рептилий (Reptile Search Algorithm) --- подход, отражающий более широкий тренд на использование метаэвристических алгоритмов оптимизации гиперпараметров нейросетевых моделей взамен градиентного обучения при малых обучающих выборках, типичных для промысловых данных (ScienceDirect, ``Prediction Model of Wax Deposition Rate in Waxy Crude Oil Pipelines by Elman Neural Network Based on Improved Reptile Search'').

Наиболее методологически развитая из недавних работ --- исследование, напрямую связывающее детальный компонентный состав нефти (доли фракций C1--C30+) со скоростью осаждения парафина через сравнение случайного леса, градиентного бустинга (XGBoost) и физически-информированной нейронной сети (PINN) на выборке из 88 образцов (Ahmadi, Physics of Fluids, 2025, том 37, вып. 7, статья 072132, DOI 10.1063/5.0276315). В этой работе градиентный бустинг демонстрирует наилучшую точность прогноза среди чисто статистических моделей, тогда как физически-информированная нейронная сеть, включающая физические ограничения на диффузионный поток парафина в структуру функции потерь, обеспечивает более устойчивую экстраполяцию за пределы диапазона обучающих данных --- что является ключевым практическим преимуществом PIML-подходов перед чисто статистическими моделями машинного обучения, не гарантирующими физически согласованное поведение вне обучающей выборки.

\subsubsection{8.3.3. Прогноз давления начала осаждения асфальтенов (АОП) и риска отложений}\label{ux43fux440ux43eux433ux43dux43eux437-ux434ux430ux432ux43bux435ux43dux438ux44f-ux43dux430ux447ux430ux43bux430-ux43eux441ux430ux436ux434ux435ux43dux438ux44f-ux430ux441ux444ux430ux43bux44cux442ux435ux43dux43eux432-ux430ux43eux43f-ux438-ux440ux438ux441ux43aux430-ux43eux442ux43bux43eux436ux435ux43dux438ux439}

Для асфальтенов основной фокус ML-моделей последних лет смещён на прогноз давления начала осаждения (Asphaltene Onset Pressure, АОП) и интегрального показателя риска повреждения продуктивности скважины асфальтеновыми отложениями. Разработаны интегрированные модели машинного обучения, одновременно классифицирующие риск асфальтенового повреждения и количественно прогнозирующие АОП по данным PVT и составу нефти (Energy \& Fuels, 2023, DOI 10.1021/acs.energyfuels.2c03319, ``Integrated Machine Learning Model for Predicting Asphaltene Damage Risk and the Asphaltene Onset Pressure''). Ансамблевые методы (бэггинг, бустинг, стекинг) также применяются для повышения устойчивости прогноза АОП по сравнению с единичными моделями регрессии (ScienceDirect, ``Robust asphaltene onset pressure prediction using ensemble learning''), а гибридные архитектуры искусственных нейронных сетей, объединённые с оптимизационными метаэвристиками, используются для прогноза устойчивости асфальтенов к осаждению непосредственно по составу и условиям PVT без необходимости проведения дорогостоящих лабораторных тестов титрования n-алканами (ResearchGate/Springer, ``Asphaltene Stability Prediction Using Hybrid Artificial Neural Network Modeling Approach''). Сопоставительные исследования нескольких алгоритмов ML для прогноза асфальтенового осаждения с новыми схемами кросс-валидации (Earth Science Informatics, 2023, DOI 10.1007/s12145-023-01075-8) указывают на систематическую проблему области: при малых и неоднородных по происхождению обучающих выборках (объединение данных разных месторождений и лабораторий) стандартная случайная кросс-валидация склонна давать завышенную оценку обобщающей способности модели, что требует более консервативных схем валидации (group k-fold по месторождению/лаборатории, blocked validation).

\subsection{8.4. Гибридные data-driven + физические (physics-informed) подходы}\label{ux433ux438ux431ux440ux438ux434ux43dux44bux435-data-driven-ux444ux438ux437ux438ux447ux435ux441ux43aux438ux435-physics-informed-ux43fux43eux434ux445ux43eux434ux44b}

Ключевой методологический тренд последних 3--4 лет --- переход от «чёрного ящика» ML к физически-информированным архитектурам, в которых уравнения сохранения (диффузионный перенос парафина, кинетика кристаллизации, тепловой баланс) включаются непосредственно в функцию потерь нейронной сети (physics-informed neural networks, PINN) или используются для генерации синтетических данных, дополняющих ограниченные экспериментальные и промысловые выборки. Как отмечено выше, для задачи прогноза скорости осаждения парафина по составу PINN-архитектура демонстрирует более физически согласованное поведение при экстраполяции по сравнению с чисто статистическими моделями (Ahmadi, 2025, DOI 10.1063/5.0276315), что типично для PIML-подходов в целом: снижение потребности в больших объёмах данных за счёт встраивания априорного физического знания и повышение доверия инженеров-нефтяников к прогнозам модели вне диапазона калибровочных данных.

Более широкая тенденция применения PINN к смежным задачам разработки и симуляции нефтяных пластов также демонстрирует, что метод разложения области (domain decomposition) в сочетании с PINN позволяет решать задачи гидродинамики пласта при разреженных промысловых данных, что напрямую применимо и к задаче неизотермической фильтрации с осаждением тяжёлых компонентов, где количество доступных для калибровки промысловых замеров (давление, температура, состав) как правило ограничено единицами-десятками скважин на месторождение. Существование отдельного недавнего обзора, посвящённого именно применениям физически-информированного машинного обучения в подземной (subsurface) инженерии в целом (Geoenergy Science and Engineering, 2025, статья 213713, DOI 10.1016/j.geoen.2025.213713), само по себе показательно: оформление PIML-методологии в самостоятельное систематизируемое направление подземной инженерии произошло именно в последние 1--2 года, что синхронно с появлением первых PIML-приложений к задаче осаждения парафина, обсуждённых выше.

Отдельное практическое применение гибридных подходов --- комплексный обзор состояния flow assurance-решений, включающий как классические экспериментальные платформы (flow loop), так и AI-driven инструменты для более широкого класса многофазных транспортных систем (в том числе для новых сред --- водорода, CO\(_2\), аммиака), что свидетельствует о переносе методологии, отработанной на задачах осаждения парафина и асфальтена, на смежные области энергоперехода (Arabian Journal for Science and Engineering, 2025, DOI 10.1007/s13369-025-10758-x).

\subsection{8.5. Открытые проблемы и перспективы развития области}\label{ux43eux442ux43aux440ux44bux442ux44bux435-ux43fux440ux43eux431ux43bux435ux43cux44b-ux438-ux43fux435ux440ux441ux43fux435ux43aux442ux438ux432ux44b-ux440ux430ux437ux432ux438ux442ux438ux44f-ux43eux431ux43bux430ux441ux442ux438}

\textbf{1. Дефицит и неоднородность промысловых данных для обучения ML-моделей.} Большинство опубликованных ML-моделей обучены на выборках порядка \(10^1\)--\(10^2\) образцов, собранных из разнородных литературных источников с различающимися лабораторными протоколами (см. раздел 7.4), что ограничивает как точность, так и обобщающую способность моделей за пределы диапазона обучающих данных. Развитие стандартизированных, открытых, достаточно больших баз данных SARA-состава, WAT и АОП, аналогичных используемым в других областях материаловедения, остаётся насущной инфраструктурной задачей области.

\textbf{2. Ограниченная физическая интерпретируемость чисто статистических ML-моделей.} Несмотря на высокие метрики точности (\(R^2 > 0.97\)) на тестовых выборках, чисто статистические модели (случайный лес, градиентный бустинг, обычные нейронные сети) не гарантируют физически согласованного поведения при экстраполяции за пределы диапазона обучающих данных --- проблема, которую решают физически-информированные архитектуры (PINN), однако их применение к задаче осаждения ПАС в пористой среде при неизотермической фильтрации пока представлено единичными работами и требует дальнейшего развития, в частности встраивания уравнений двухфазной фильтрации Дарси с учётом снижения проницаемости из-за отложений непосредственно в архитектуру физически-информированной модели.

\textbf{3. Масштабирование от лабораторных экспериментов к промысловым условиям.} Как показано в разделе 7, количественное согласование данных cold finger, flow loop и промысловых замеров остаётся нетривиальной задачей из-за различий в локальной гидродинамике; относительно мало работ систематически используют ML для построения именно моделей масштабирования (scale-up), а не прямого прогноза целевой переменной по составу.

\textbf{4. Интеграция композиционного моделирования осаждения с полномасштабным гидродинамическим симулятором пласта.} Для задачи неизотермической фильтрации в пористой среде (в отличие от трубопроводного транспорта) остаётся открытым вопрос эффективной с вычислительной точки зрения интеграции детальной кинетики осаждения ПАС и связанного с этим снижения проницаемости/пористости в промышленные симуляторы пласта (black-oil, композиционные), где полный набор кинетических уравнений для каждого узла сетки может быть вычислительно неподъёмным; здесь перспективным направлением видится использование обученных суррогатных ML/PINN-моделей вместо прямого решения детальной кинетики в каждой ячейке сетки на каждом временном шаге, по аналогии с уже применяемым подходом ускорения физически-информированными нейросетевыми суррогатами при разреженных данных в реcервуарной симуляции.

\textbf{5. Неопределённость и её количественная оценка (uncertainty quantification).} Действующие ML- и гибридные модели в подавляющем большинстве публикаций дают точечный прогноз без байесовской оценки доверительных интервалов; более широкое внедрение гауссовских процессов и байесовских нейронных сетей, естественным образом дающих оценку неопределённости прогноза, могло бы напрямую использоваться при планировании интервалов между операциями депарафинизации/удаления асфальтеновых отложений с учётом риска, а не только средней ожидаемой скорости отложения.

\subsection{8.6. Заключение по трендам}\label{ux437ux430ux43aux43bux44eux447ux435ux43dux438ux435-ux43fux43e-ux442ux440ux435ux43dux434ux430ux43c}

Обзор литературы за 2019--2026 гг. показывает отчётливое смещение исследовательского фокуса от чисто механистических композиционных моделей к гибридным подходам, сочетающим детальную физико-химическую характеризацию состава нефти (μSARA-HPLC, ФТ-ИЦР-МС, ЯМР) с методами машинного обучения для прогноза ключевых параметров (WAT, АОП, скорость отложения). При этом наиболее перспективным направлением ближайших лет представляется не замена механистических моделей чисто статистическими, а их синтез в форме физически-информированных нейронных сетей и суррогатных моделей, обученных на результатах детальных механистических симуляций и ограниченных промысловых данных одновременно. Ключевым сдерживающим фактором остаётся не столько выбор архитектуры модели, сколько качество, объём и стандартизация экспериментальных и промысловых данных, используемых для калибровки и валидации --- вывод, симметричный тому, что было получено при анализе классических механистических моделей в разделе 7.

\newpage

\section{Заключение}\label{ux437ux430ux43aux43bux44eux447ux435ux43dux438ux435}

\subsection{Сопоставление подходов к моделированию}\label{ux441ux43eux43fux43eux441ux442ux430ux432ux43bux435ux43dux438ux435-ux43fux43eux434ux445ux43eux434ux43eux432-ux43a-ux43cux43eux434ux435ux43bux438ux440ux43eux432ux430ux43dux438ux44e}

Проведённый анализ литературы позволяет выделить три уровня описания задачи осаждения тяжёлых компонентов нефти при неизотермической фильтрации, каждый из которых отвечает за свой аспект физики процесса и обладает собственным балансом точности и вычислительной стоимости.

\textbf{Термодинамический уровень} (раздел 1, раздел 5.3) отвечает на вопрос «сколько твёрдой фазы существует в равновесии при данных \(P\), \(T\) и составе» и представлен спектром моделей возрастающей сложности: от эмпирических корреляций WAT и скейлинговых уравнений для давления начала осаждения асфальтенов (быстрые, но плохо экстраполируемые за пределы калибровочного диапазона) через модель твёрдого раствора и модель множественной твёрдой фазы для парафинов и модель регулярных растворов/мицеллизации для асфальтенов (физически более обоснованные, но требующие параметризации по SARA-составу) до статистической теории ассоциирующих флюидов PC-SAFT и кубически-ассоциативных уравнений состояния CPA (наиболее физически строгие, явно учитывающие полидисперсность и ассоциативные взаимодействия, но вычислительно наиболее затратные и требующие наибольшего числа настраиваемых параметров сегмента). Принципиальное различие между парафиновой и асфальтеновой термодинамикой --- в самой природе твёрдой фазы: кристаллизация парафинов есть переход из истинного молекулярного раствора, тогда как дестабилизация асфальтенов есть разрушение уже существующей коллоидной структуры (наноагрегатов и их кластеров по модели Йена--Маллинса), что и объясняет немонотонную, колоколообразную зависимость растворимости асфальтенов от давления --- свойство, не имеющее аналога в термодинамике парафинов.

\textbf{Кинетический уровень} (раздел 2, раздел 5.4) отвечает на вопрос «как быстро и где именно термодинамически неустойчивая твёрдая фаза закрепляется на скелете породы или стенке канала» и представлен моделями диффузионного переноса (закон Фика для растворённых молекул парафина, модель Хансена), сдвиговой дисперсии и броуновской/гравитационной седиментации для взвешенных частиц, а также парными моделями захвата--выноса (deposition--entrainment) типа Иванса--Шехтера/Цивана для описания накопления и последующей эрозии осевшего материала. В отличие от термодинамического равновесия, кинетические константы этого уровня в подавляющем большинстве публикаций определяются полуэмпирически, путём обратной калибровки (history matching) по ограниченному набору лабораторных или промысловых данных, что ограничивает их переносимость между месторождениями и типами нефти --- этот вывод последовательно подтверждается и в разделе 6.4 (численные модели), и в разделе 7.4 (экспериментальная валидация).

\textbf{Континуальный (пластовый) уровень} (разделы 3--4, раздел 6) объединяет термодинамический и кинетический уровни в единую сопряжённую систему уравнений сохранения массы, импульса (закон Дарси) и энергии, дополненную моделью изменения проницаемости и пористости (Козени--Карман и модификации) под действием накопленного осадка. Именно на этом уровне проявляется ключевая физическая особенность задачи --- двусторонняя обратная связь между осаждением и течением: осадок снижает проницаемость → снижение проницаемости меняет локальное поле скоростей и температур → изменённые условия меняют локальную скорость дальнейшего осаждения. Численная реализация этой связанной системы (раздел 6) в подавляющем большинстве случаев опирается на метод конечных объёмов с полностью неявной или IMPES-подобной схемой по времени, при этом сравнение независимых промышленных симуляторов (UTCOMP, CMG/GEM, ECLIPSE) на идентичных входных PVT-данных показывает существенные расхождения в прогнозе давления начала осаждения асфальтенов и массы осадка, что указывает на чувствительность результата не только к выбору термодинамической модели, но и к деталям численной реализации проверки фазовой стабильности.

Таблица ниже суммирует эту трёхуровневую систематизацию.

{\def\LTcaptype{none} % do not increment counter
\begin{xltabular}{\linewidth}{@{}
  >{\raggedright\arraybackslash}X
  >{\raggedright\arraybackslash}X
  >{\raggedright\arraybackslash}X
  >{\raggedright\arraybackslash}X
  >{\raggedright\arraybackslash}X@{}}
\toprule\noalign{}
Уровень описания & Ключевой вопрос & Основные модели (парафины) & Основные модели (асфальтены/смолы) & Основной источник неопределённости \\
\midrule\noalign{}
\endhead
\bottomrule\noalign{}
\endlastfoot
Термодинамический & Сколько твёрдой фазы в равновесии? & Модель твёрдого раствора, multi-solid model, Флори--Хаггинс & Модель Йена--Маллинса, регулярные растворы, PC-SAFT, CPA-EOS & Компонентная характеризация (SARA, SCN), параметры сегмента SAFT \\
Кинетический & Как быстро и где закрепляется осадок? & Диффузия Фика, модель Хансена/Бурже, сдвиговая дисперсия & Диффузионно/реакционно-лимитированная агрегация (DLCA/RLCA), модели deposition--entrainment & Полуэмпирические константы, откалиброванные по ограниченным данным \\
Континуальный (пластовый) & Как осаждение влияет на течение и наоборот? & Связанные Дарси+энергия+кинетика, Козени--Карман & То же + явная связь с P--T--составом потока & Численная реализация фазовой стабильности, масштаб сетки, upscaling \\
\end{xltabular}
}

\subsection{Общие выводы}\label{ux43eux431ux449ux438ux435-ux432ux44bux432ux43eux434ux44b}

Во-первых, несмотря на структурные различия между парафинами и асфальтенами, оба класса задач подчиняются единой логической схеме «термодинамика равновесия → кинетика переноса и захвата → континуальная связь с фильтрацией и теплопереносом → обратное влияние на условия дестабилизации», и наиболее продуктивные модели последних десятилетий (ADEPT и его многофазное расширение MAD-ADEPT, композиционные симуляторы с явной асфальтеновой псевдокомпонентой, операторно-ориентированные архитектуры типа DARTS) явно реализуют все три уровня в связанном виде, а не рассматривают их изолированно.

Во-вторых, систематическое сопоставление модельных прогнозов с лабораторными (cold finger, flow loop, ЯМР, HPLC-SARA) и промысловыми (кавернометрия, история давления, данные пиггования) данными, обсуждённое в разделе 7, показывает, что при качественной калибровке современные модели способны воспроизводить наблюдаемые профили отложений с отклонением порядка нескольких процентов по ключевым метрикам. Однако такая точность достижима только при доступности надёжных входных данных о компонентном составе и кинетических параметрах --- то есть лимитирующим фактором точности прогноза на практике чаще оказывается качество входных данных, а не структура самой математической модели или численная схема её решения.

В-третьих, обзор современных тенденций (раздел 8) фиксирует отчётливый сдвиг от чисто механистических моделей к гибридным подходам, сочетающим детальную физико-химическую характеризацию состава (высокоразрешающая масс-спектрометрия, автоматизированный SARA-анализ) с методами машинного обучения для прогноза интегральных параметров (WAT, АОП, скорость отложения). При этом наиболее физически обоснованным направлением развития представляются не чисто статистические ML-модели, а физически-информированные архитектуры (PINN) и суррогатные модели, обученные на результатах детальных механистических симуляций и ограниченных промысловых данных одновременно --- поскольку именно они сохраняют устойчивость экстраполяции за пределы диапазона обучающих данных, критически важную для промысловых приложений.

\subsection{Открытые проблемы}\label{ux43eux442ux43aux440ux44bux442ux44bux435-ux43fux440ux43eux431ux43bux435ux43cux44b}

Систематизация материала разделов 6--8 позволяет сформулировать следующий консолидированный список открытых проблем области, актуальных на момент подготовки настоящего обзора (2026 г.):

\begin{enumerate}
\def\labelenumi{\arabic{enumi}.}
\item
  \textbf{Отсутствие универсальных, физически обоснованных (не откалиброванных заново под каждое месторождение) корреляций для кинетических констант захвата и выноса частиц} --- как для парафинов, так и в особенности для асфальтенов, где коэффициенты деполимеризации/агрегации наноагрегатов и константы прилипания к стенке поры остаются существенно эмпирическими.
\item
  \textbf{Вычислительная стоимость интеграции детальной кинетики осаждения в промышленные композиционные симуляторы пласта} --- при сеточных моделях промыслового масштаба явное решение кинетического уравнения в каждой ячейке на каждом временном шаге может быть неподъёмным, что мотивирует разработку суррогатных ML/PINN-моделей и апскейлинговых процедур.
\item
  \textbf{Масштабирование от лабораторных экспериментов (cold finger, flow loop) к промысловым условиям} --- количественное согласование результатов этих методов друг с другом и с промысловыми наблюдениями остаётся нетривиальной задачей из-за различий в локальной гидродинамике и сдвиговых напряжениях.
\item
  \textbf{Дефицит стандартизированных, достаточно больших и мультилабораторных открытых баз данных} по SARA-составу, WAT и AOP, ограничивающий как точность классических корреляций, так и обобщающую способность ML-моделей.
\item
  \textbf{Количественная оценка неопределённости прогноза} (uncertainty quantification) --- подавляющее большинство публикуемых моделей, как механистических, так и ML-based, дают точечный прогноз без байесовских доверительных интервалов, что ограничивает их использование при риск-ориентированном планировании профилактических мероприятий.
\item
  \textbf{Разделение вкладов термодинамической дестабилизации, кинетики захвата/выноса и сопутствующей многофазной гидродинамики} при интерпретации косвенных промысловых данных (динамика давления, кавернометрия) --- как показывает переинтерпретация классического кейса месторождения Marrat, наблюдаемые тренды не всегда однозначно указывают на рост слоя отложений и могут отражать иные гидродинамические эффекты.
\end{enumerate}

Совокупность этих проблем указывает на то, что дальнейшее развитие математического моделирования осаждения парафинов, асфальтенов и смол при неизотермической фильтрации будет определяться не столько появлением принципиально новых термодинамических теорий или численных схем --- этот инструментарий в значительной мере сформирован, --- сколько прогрессом в интеграции разноуровневых моделей друг с другом, качеством и объёмом данных для их калибровки и валидации, а также способностью гибридных физически-информированных методов обеспечить надёжную экстраполяцию за пределы условий, для которых модели были откалиброваны.

\newpage

\section{Список литературы}\label{ux441ux43fux438ux441ux43eux43a-ux43bux438ux442ux435ux440ux430ux442ux443ux440ux44b}

\begin{enumerate}
\def\labelenumi{\arabic{enumi}.}
\tightlist
\item
  Anjushri S. Kurup, Jianxin Wang, Hariprasad J. Subramani, Jill Buckley, Jefferson L. Creek, Walter G. Chapman. (2012). Revisiting Asphaltene Deposition Tool (ADEPT): Field Application. Energy \& Fuels, 26, 5702-5710. DOI: 10.1021/ef300714p.
\item
  Md Rashedul Islam, Yifan Hao, Chau-Chyun Chen. (2020). Aggregation thermodynamics of asphaltenes: Prediction of asphaltene precipitation in petroleum fluids with NRTL-SAC. Fluid Phase Equilibria, 520, 112655. DOI: 10.1016/j.fluid.2020.112655.
\item
  Ahmadi, M. (2025). From composition to deposition: A machine learning framework for wax deposition prediction in petroleum fluids using compositional properties. Physics of Fluids, 37(7), 072132. DOI: 10.1063/5.0276315.
\item
  Alhosani, Ahmed, Ravichandran, Sriram, Daraboina, Nagu. (2020). Review of Asphaltene Deposition Modeling in Oil and Gas Production. Energy \& Fuels, 35(2), 965-986. DOI: 10.1021/acs.energyfuels.0c02981.
\item
  Alimohammadi, Sepideh, Zendehboudi, Sohrab, James, Lesley. (2019). A comprehensive review of asphaltene deposition in petroleum reservoirs: Theory, challenges, and tips. Fuel, 252, 753-791. DOI: 10.1016/j.fuel.2019.03.016.
\item
  (2024). Robust asphaltene onset pressure prediction using ensemble learning. Geoenergy Science and Engineering. DOI: 10.1016/j.geoen.2024.212857.
\item
  (2023). Integrated Machine Learning Model for Predicting Asphaltene Damage Risk and the Asphaltene Onset Pressure. Energy \& Fuels. DOI: 10.1021/acs.energyfuels.2c03319.
\item
  Hamed Darabi, Mahdy Shirdel, M. Hosein Kalaei, Kamy Sepehrnoori. (2014). Aspects of Modeling Asphaltene Deposition in a Compositional Coupled Wellbore/Reservoir Simulator. SPE Improved Oil Recovery Symposium. DOI: 10.2118/169121-ms.
\item
  Asphaltene deposition, plugging. (2025). Asphaltene deposition and plugging. DOI: 10.2118/pw0043.
\item
  Ahmed Alhosani, Nagu Daraboina. (2020). Modeling of asphaltene deposition during oil/gas flow in wellbore. Fuel, 280, 118617. DOI: 10.1016/j.fuel.2020.118617.
\item
  John G. Reynolds. (1994). Chapter 10 Effects of Asphaltene Precipitation on the Size of Vanadium-, Nickel-, And Sulfur-Containing Compounds in Heavy Crude Oils and Residua. Developments in Petroleum Science, 233-248. DOI: 10.1016/s0376-7361(09)70257-2.
\item
  (2024). Asphaltene Stability Prediction Using Hybrid Artificial Neural Network Modeling Approach. Springer Nature.
\item
  Banki, Reza, Hoteit, Hussein, Firoozabadi, Abbas. (2008). Mathematical formulation and numerical modeling of wax deposition in pipelines from enthalpy-porosity approach and irreversible thermodynamics. International Journal of Heat and Mass Transfer, 51(13-14), 3387-3398. DOI: 10.1016/j.ijheatmasstransfer.2007.11.012.
\item
  Borodin, S. L., Musakaev, N. G., Belskikh, Denis S. (2022). Mathematical Modeling of a Non-Isothermal Flow in a Porous Medium Considering Gas Hydrate Decomposition: A Review. Mathematics, 10(24), 4674. DOI: 10.3390/math10244674.
\item
  Brown, T. S., Niesen, Vicki G., Erickson, Dale. (1993). Measurement and Prediction of the Kinetics of Paraffin Deposition. SPE Annual Technical Conference and Exhibition. DOI: 10.2118/26548-ms.
\item
  Buenrostro-González, Eduardo, Lira-Galeana, C., Gil‐Villegas, Alejandro, Wu, Jianzhong. (2004). Asphaltene precipitation in crude oils: Theory and experiments. AIChE Journal, 50, 2552-2570. DOI: 10.1002/aic.10243.
\item
  Burger, E. D., Perkins, T.K., Striegler, J.H. (1981). Studies of Wax Deposition in the Trans Alaska Pipeline. Journal of Petroleum Technology, 33, 1075-1086. DOI: 10.2118/8788-pa.
\item
  Chapman, Walter G., Gubbins, Keith E., Jackson, George, Radosz, Maciej. (1989). SAFT: Equation-of-state solution model for associating fluids. Fluid Phase Equilibria, 52, 31-38. DOI: 10.1016/0378-3812(89)80308-5.
\item
  Civan, Faruk. (2015). Modified Formulations of Particle Deposition and Removal Kinetics in Saturated Porous Media. Transport in Porous Media, 111(2), 381-410. DOI: 10.1007/s11242-015-0600-z.
\item
  (2018). Effect of the Flow Field on the Wax Deposition and Performance of Wax Inhibitors: Cold Finger and Flow Loop Testing. Energy \& Fuels. DOI: 10.1021/acs.energyfuels.7b00253.
\item
  (2009). Solids Deposition during ``Cold Flow'\,' of Wax-Solvent Mixtures in a Flow-loop Apparatus with Heat Transfer. Energy \& Fuels. DOI: 10.1021/ef900224r.
\item
  (2018). Combining of intelligent models through committee machine for estimation of wax deposition. Petroleum.
\item
  Abdulaziz Al-Qasim, Mudhish A AlDawsari. (2017). Comparison Study of Asphaltene Precipitation Models Using UTCOMP, CMG/GEM and ECLIPSE Simulators. SPE Oil and Gas India Conference and Exhibition. DOI: 10.2118/185370-ms.
\item
  K. H. Coats, K. H. Coats, L. K. Thomas, R. G. Pierson. (1995). Compositional and Black Oil Reservoir Simulation. SPE Reservoir Simulation Symposium. DOI: 10.2118/29111-ms.
\item
  H. Hajibeygi, H. A. Tchelepi. (2013). Compositional Multiscale Finite-Volume Formulation. SPE Reservoir Simulation Symposium. DOI: 10.2118/163664-ms.
\item
  K. Gonzalez, M. A. Barrufet, H. Nasrabadi. (2014). Development of a Compositional Reservoir Simulator Including Asphaltene Precipitation from a Thermodynamic Consistent Model. SPE Latin America and Caribbean Petroleum Engineering Conference. DOI: 10.2118/169401-ms.
\item
  Coutinho, João A. P., Stenby, Erling H. (1996). Predictive Local Composition Models for Solid/Liquid Equilibrium in n-Alkane Systems: Wilson Equation for Multicomponent Systems. Industrial \& Engineering Chemistry Research, 35, 918-925. DOI: 10.1021/ie950447u.
\item
  Wenlong Jia, Ryosuke Okuno. (2018). Modeling of asphaltene and water associations in petroleum reservoir fluids using cubic‐plus‐association EOS. AIChE Journal, 64, 3429-3442. DOI: 10.1002/aic.16191.
\item
  Hadi Nasrabadi, Joachim Moortgat, Abbas Firoozabadi. (2013). A New Three-Phase Multicomponent Compositional Model for Asphaltene Precipitation Using CPA-EOS. SPE Reservoir Simulation Symposium. DOI: 10.2118/163587-ms.
\item
  P. Bedrikovetsky, A. Santos, A. Siqueira, A. L. Souza, F. Shecaira. (2003). A Stochastic Model for Deep Bed Filtration and Well Impairment. SPE European Formation Damage Conference. DOI: 10.2118/82230-ms.
\item
  (2024). Prediction Model of Wax Deposition Rate in Waxy Crude Oil Pipelines by Elman Neural Network Based on Improved Reptile Search Algorithm. Journal of Petroleum Science and Engineering. DOI: 10.1016/j.petrol.2024.212763.
\item
  Feng, Xiao, Zeng, Jianhui, Ma, Yong, Jia, Kaiyu, Qiao, Juncheng, Zhang, Yongchao, Feng, Sen.~(2017). Asphaltene Deposition Preference and Permeability Reduction Mechanisms in Oil Reservoirs: Evidence from Combining X-ray Microtomography with Fluorescence Microscopy. Energy \& Fuels, 31(10), 10467-10478. DOI: 10.1021/acs.energyfuels.7b01389.
\item
  Sebastián Echavarría-Montaña, Steven Velásquez, Nicolás Bueno, Juan David Valencia, Hillmert~Alexander Solano, Juan Manuel Mejía. (2021). Semi-Implicit Finite Volume Procedure for Compositional Subsurface Flow Simulation in Highly Anisotropic Porous Media. Fluids, 6, 341. DOI: 10.3390/fluids6100341.
\item
  Xiangjun Qin, Peng Wang, Kamy Sepehrnoori, Gary A. Pope. (2000). Modeling Asphaltene Precipitation in Reservoir Simulation. Industrial \& Engineering Chemistry Research, 39, 2644-2654. DOI: 10.1021/ie990781g.
\item
  (2025). A Comprehensive Review of Flow Assurance in the Energy Transition: Flow Loop Platforms and AI-Driven Solutions for Hydrogen, CO2, and Ammonia Transport. Arabian Journal for Science and Engineering. DOI: 10.1007/s13369-025-10758-x.
\item
  Groß, Joachim, Sadowski, Gabriele. (2001). Perturbed-Chain SAFT: An Equation of State Based on a Perturbation Theory for Chain Molecules. Industrial \& Engineering Chemistry Research, 40, 1244-1260. DOI: 10.1021/ie0003887.
\item
  Gruesbeck, C., Collins, R. E. (1982). Entrainment and Deposition of Fine Particles in Porous Media. Society of Petroleum Engineers Journal, 22(6), 847-856. DOI: 10.2118/8430-PA.
\item
  Han, Jiang-Xia, Xue, Liang, Wei, Yu-Shu, others. (2023). Physics-informed neural network-based petroleum reservoir simulation with sparse data using domain decomposition. Petroleum Science, 20(6), 3450-3460. DOI: 10.1016/j.petsci.2023.10.019.
\item
  Hansen, Jens Henrik, Fredenslund, Aa., Pedersen, Karen Schou, Rønningsen, Hans Petter. (1988). A thermodynamic model for predicting wax formation in crude oils. AIChE Journal, 34, 1937-1942. DOI: 10.1002/aic.690341202.
\item
  Huang, Zhenyu, Lee, H. S., Senra, Michael, Fogler, H. Scott. (2010). A fundamental model of wax deposition in subsea oil pipelines. AIChE Journal, 57, 2955-2964. DOI: 10.1002/aic.12517.
\item
  Kargarpour, M., Dandekar, A. Y. (2016). Analysis of asphaltene deposition in Marrat oil well string: a new approach. Journal of Petroleum Exploration and Production Technology, 6, 845-856. DOI: 10.1007/s13202-015-0221-7.
\item
  Katz, David L., Firoozabadi, Abbas. (1978). Predicting Phase Behavior of Condensate/Crude-Oil Systems Using Methane Interaction Coefficients. Journal of Petroleum Technology, 30, 1649-1655. DOI: 10.2118/6721-pa.
\item
  B. ZareNezhad, H. Parsa. (2015). A Rigorous Kinetic Model for Description of Asphaltene Deposition in Porous Media of Petroleum Reservoirs. Petroleum Science and Technology, 33, 1853-1860. DOI: 10.1080/10916466.2015.1107846.
\item
  Kor, Pooria, Kharrat, Riyaz, Ayoubi, Ali. (2017). Comparison and evaluation of several models in prediction of asphaltene deposition profile along an oil well: a case study. Journal of Petroleum Exploration and Production Technology, 7, 497-510. DOI: 10.1007/s13202-016-0269-z.
\item
  Krauss, Eva D., Mays, D. C. (2013). Modification of the Kozeny-Carman Equation to Quantify Formation Damage by Fines in Clean Unconsolidated Porous Media. SPE Reservoir Evaluation \& Engineering, 17(4), 466-472. DOI: 10.2118/165148-PA.
\item
  Shaghayegh Darjani, Archana Jagadisan, Joel Koplik, Sanjoy Banerjee. (2022). Lattice-Gas Model for Asphaltene Interactions Observed at Interfaces in the Context of the Yen-Mullins Model. Energy \& Fuels, 36, 8778-8785. DOI: 10.1021/acs.energyfuels.2c01413.
\item
  Leontaritis, Kosta J., Mansoori, G. Ali. (1987). Asphaltene Flocculation During Oil Production and Processing: A Thermodynamic Colloidal Model. SPE International Symposium on Oilfield Chemistry. DOI: 10.2118/16258-ms.
\item
  Leontaritis, Kosta J., Mansoori, G. Ali. (1988). Asphaltene deposition: a survey of field experiences and research approaches. Journal of Petroleum Science and Engineering, 1(3), 229-239. DOI: 10.1016/0920-4105(88)90013-7.
\item
  Lira-Galeana, C., Firoozabadi, Abbas, Prausnitz, John M. (1996). Thermodynamics of wax precipitation in petroleum mixtures. AIChE Journal, 42, 239-248. DOI: 10.1002/aic.690420120.
\item
  (2025). Crude Oil Analysis by Low-Field NMR Relaxometry. Magnetic Resonance in Chemistry. DOI: 10.1002/mrc.70096.
\item
  K. C. Hong. (1982). Lumped-Component Characterization of Crude Oils for Compositional Simulation. SPE Enhanced Oil Recovery Symposium. DOI: 10.2118/10691-ms.
\item
  Saman Naseri, Saeid Jamshidi, Vahid Taghikhani. (2020). A new multiphase and dynamic asphaltene deposition tool (MAD-ADEPT) to predict the deposition of asphaltene particles on tubing wall. Journal of Petroleum Science and Engineering, 195, 107553. DOI: 10.1016/j.petrol.2020.107553.
\item
  (2023). Comparative study of machine learning algorithms in predicting asphaltene precipitation with a novel validation technique. Earth Science Informatics. DOI: 10.1007/s12145-023-01075-8.
\item
  Bruce F. Kohse, Long X. Nghiem. (2004). Modelling Asphaltene Precipitation and Deposition in a Compositional Reservoir Simulator. SPE/DOE Symposium on Improved Oil Recovery. DOI: 10.2118/89437-ms.
\item
  Monteagudo, Jorge E., Rajagopal, Krishnaswamy, Lage, Paulo L. C. (2002). Simulating oil flow in porous media under asphaltene deposition. Chemical Engineering Science, 57(3), 323-337. DOI: 10.1016/S0009-2509(01)00407-9.
\item
  Qinghao Wu, Douglas J. Seifert, Andrew E. Pomerantz, Oliver C. Mullins, Richard N. Zare. (2014). Constant Asphaltene Molecular and Nanoaggregate Mass in a Gravitationally Segregated Reservoir. Energy \& Fuels, 28, 3010-3015. DOI: 10.1021/ef500281s.
\item
  Andrew E. Pomerantz, Qinghao Wu, Oliver C. Mullins, Richard N. Zare. (2015). Laser-Based Mass Spectrometric Assessment of Asphaltene Molecular Weight, Molecular Architecture, and Nanoaggregate Number. Energy \& Fuels, 29, 2833-2842. DOI: 10.1021/ef5020764.
\item
  Oliver C. Mullins, Hassan Sabbah, Joëlle Eyssautier, Andrew E. Pomerantz, Loïc Barré, A. Ballard Andrews, Yosadara Ruiz-Morales, Farshid Mostowfi, Richard McFarlane, Lamia Goual, Richard Lepkowicz, Thomas Cooper, Jhony Orbulescu, Roger M. Leblanc, John Edwards, Richard N. Zare. (2012). Advances in Asphaltene Science and the Yen--Mullins Model. Energy \& Fuels, 26, 3986-4003. DOI: 10.1021/ef300185p.
\item
  Mullins, Oliver C., Seifert, Douglas J., Zuo, Julian Y., Zeybek, Murat. (2012). Clusters of Asphaltene Nanoaggregates Observed in Oilfield Reservoirs. Energy \& Fuels, 27, 1752-1761. DOI: 10.1021/ef301338q.
\item
  Nichita, Dan Vladimir, Goual, Lamia, Firoozabadi, Abbas. (2001). Wax Precipitation in Gas Condensate Mixtures. SPE Production \& Facilities, 16, 250-259. DOI: 10.2118/74686-pa.
\item
  Nikiforov, A. I., Sadovnikov, R. V., Nikiforov, G. A. (2022). Modeling of Paraffin Deposition During Oil Production by Cold Water Injection. Lobachevskii Journal of Mathematics, 43(5), 1178-1183. DOI: 10.1134/S199508022208025X.
\item
  (2023). Novel Nuclear Magnetic Resonance Techniques To Assess the Wax Precipitation Evolution in Crude Oil Systems. Energy \& Fuels. DOI: 10.1021/acs.energyfuels.2c03309.
\item
  Ochieng, R., others. (2022). Mathematical Modeling of Wax Deposition in Field-Scale Crude Oil Pipeline Systems. Journal of Applied Mathematics, 2022, 2845221. DOI: 10.1155/2022/2845221.
\item
  Ochieng, R., others. (2023). Numerical Study of Wax Deposition from Multiphase Flow in Oil Pipelines with Heat and Mass Transfer. Mathematical Problems in Engineering, 2023, 1173505. DOI: 10.1155/2023/1173505.
\item
  Pan, Huanquan, Firoozabadi, Abbas, Fotland, Per. (1997). Pressure and Composition Effect on Wax Precipitation: Experimental Data and Model Results. SPE Production \& Facilities, 12, 250-258. DOI: 10.2118/36740-pa.
\item
  B. Soltani, S. Esteghamat, R. Kharrat. (2017). Asphaltene Precipitation Modeling by Using of PC-SAFT Equation of State. Proceedings. DOI: 10.3997/2214-4609.201701511.
\item
  M. Masoudi, S. Parvin, R. Miri, S. Kord, H. Hellevang. (2021). Implementation of PC-SAFT Equation of State Into MRST Compositional for Modelling of Asphaltene Precipitation. 82nd EAGE Annual Conference \& Exhibition, 1-5. DOI: 10.3997/2214-4609.202011432.
\item
  Ali Kariman Moghaddam, Saeid Jamshidi. (2022). Performance evaluation and improvement of PC-SAFT equation of state for the asphaltene precipitation modeling during mixing with various fluid types. Fluid Phase Equilibria, 554, 113340. DOI: 10.1016/j.fluid.2021.113340.
\item
  Mohammad Tavakkoli, Andrew Chen, Francisco M. Vargas. (2016). Rethinking the modeling approach for asphaltene precipitation using the PC-SAFT Equation of State. Fluid Phase Equilibria, 416, 120-129. DOI: 10.1016/j.fluid.2015.11.003.
\item
  Pedersen, Karen Schou. (1995). Prediction of Cloud Point Temperatures and Amount of Wax Precipitation. SPE Production \& Facilities, 10, 46-49. DOI: 10.2118/27629-pa.
\item
  (2025). Review of physics-informed machine learning (PIML) methods applications in subsurface engineering. Geoenergy Science and Engineering, 250, 213713. DOI: 10.1016/j.geoen.2025.213713.
\item
  Yutaka Onaka, Kozo Sato. (2021). Dynamics of pore-throat plugging and snow-ball effect by asphaltene deposition in porous media micromodels. Journal of Petroleum Science and Engineering, 207, 109176. DOI: 10.1016/j.petrol.2021.109176.
\item
  H. N. Dunning, J. W. Moore, Herman Bieber, R. B. Williams. (1960). Porphyrin, Nickel, Vanadium, and Nitrogen in Petroleum. Journal of Chemical \& Engineering Data, 5, 546-549. DOI: 10.1021/je60008a036.
\item
  (2018). Prediction of wax content in crude oil and petroleum fraction by proton NMR. Petroleum Science and Technology. DOI: 10.1080/10916466.2018.1536713.
\item
  (2019). Modeling Wax Disappearance Temperature Using Advanced Intelligent Frameworks. Energy \& Fuels. DOI: 10.1021/acs.energyfuels.9b03296.
\item
  Riazi, Mohammad R., Daubert, Thomas E. (1987). Characterization parameters for petroleum fractions. Industrial \& Engineering Chemistry Research, 26, 755-759. DOI: 10.1021/ie00064a023.
\item
  Ring, Jasper N., Wattenbarger, R. A., Keating, James F., Peddibhotla, Sriram. (1994). Simulation of Paraffin Deposition in Reservoirs. SPE Production \& Facilities, 9(1), 36-42. DOI: 10.2118/24069-PA.
\item
  (2021). Predicting wax deposition using robust machine learning techniques. South African Journal of Chemical Engineering. DOI: 10.1016/j.sajce.2021.02.003.
\item
  ACS Omega authors. (2025). ACS Omega. DOI: 10.1021/acsomega.5c00439.
\item
  M. Jamialahmadi, K. Ahmadi. (2003). A New Scaling Equation for Modeling of Asphaltene Precipitation. Nigeria Annual International Conference and Exhibition. DOI: 10.2118/85673-ms.
\item
  Singh, Probjot, Venkatesan, R., Fogler, H. Scott, Nagarajan, Nagi. (2000). Formation and aging of incipient thin film wax‐oil gels. AIChE Journal, 46, 1059-1074. DOI: 10.1002/aic.690460517.
\item
  Singh, Probjot, Venkatesan, R., Fogler, H. Scott, Nagarajan, Nagi. (2001). Morphological evolution of thick wax deposits during aging. AIChE Journal, 47, 6-18. DOI: 10.1002/aic.690470103.
\item
  Soulgani, Bahram Soltani, Tohidi, Bahman, Jamialahmadi, M., Rashtchian, Davood. (2011). Modeling Formation Damage due to Asphaltene Deposition in the Porous Media. Energy \& Fuels, 25(2), 753-761. DOI: 10.1021/ef101195a.
\item
  Ana C. R. Sodero, Hugo Santos Silva, Patricia Guevara Level, Brice Bouyssiere, Jean-Pierre Korb, Hervé Carrier, Ahmad Alfarra, Didier Bégué, Isabelle Baraille. (2016). Investigation of the Effect of Sulfur Heteroatom on Asphaltene Aggregation. Energy \& Fuels, 30, 4758-4766. DOI: 10.1021/acs.energyfuels.6b00757.
\item
  Narmadha Rajan Babu, Pei-Hsuan Lin, Mohammed I. L. Abutaqiya, Caleb J. Sisco, Jianxin Wang, Francisco M. Vargas. (2019). Systematic Investigation of Asphaltene Deposition in the Wellbore and Near-Wellbore Region of a Deepwater Oil Reservoir under Gas Injection. Part 2: Computational Fluid Dynamics Modeling of Asphaltene Deposition. Energy \& Fuels, 33, 3645-3661. DOI: 10.1021/acs.energyfuels.8b03239.
\item
  Taheri-Shakib, Jaber, Shekarifard, Ali, Naderi, Hassan. (2018). Experimental investigation of the asphaltene deposition in porous media: Accounting for the microwave and ultrasonic effects. Journal of Petroleum Science and Engineering, 163, 453-462. DOI: 10.1016/j.petrol.2018.01.017.
\item
  Mark Khait, Denis Voskov. (2019). Integrated Framework for Modelling of Thermal-Compositional Multiphase Flow in Porous Media. SPE Reservoir Simulation Conference. DOI: 10.2118/193932-ms.
\item
  E. Rogel. (2004). Thermodynamic Modeling of Asphaltene Aggregation. Langmuir, 20, 1003-1012. DOI: 10.1021/la035339o.
\item
  Yurko Duda, C. Lira-Galeana. (2006). Thermodynamics of asphaltene structure and aggregation. Fluid Phase Equilibria, 241, 257-267. DOI: 10.1016/j.fluid.2005.12.043.
\item
  Victorov, Alexey I., Firoozabadi, Abbas. (1996). Thermodynamic micellizatin model of asphaltene precipitation from petroleum fluids. AIChE Journal, 42, 1753-1764. DOI: 10.1002/aic.690420626.
\item
  Wang, Meng, Hao, Yifan, Islam, Md Rashedul, Chen, Chau‐Chyun. (2016). Aggregation thermodynamics for asphaltene precipitation. AIChE Journal, 62, 1254-1264. DOI: 10.1002/aic.15173.
\item
  (2025). Toward accurate modeling of wax appearance temperature (WAT) using intelligent ensemble techniques: Application to heavy crude oils. Journal of Molecular Liquids. DOI: 10.1016/j.molliq.2025.128500.
\item
  (2024). New intelligent models for predicting wax appearance temperature using experimental data - Flow assurance implications. Fuel. DOI: 10.1016/j.fuel.2024.133423.
\item
  Xiangrong Nie, Shenglai Yang. (2014). Numerical simulation of paraffin wax deposition in waxy crude oil reservoir with cold water flooding. Geosystem Engineering, 17, 219-225. DOI: 10.1080/12269328.2014.959623.
\item
  John A. Svendsen. (1993). Mathematical modeling of wax deposition in oil pipeline systems. AIChE Journal, 39, 1377-1388. DOI: 10.1002/aic.690390815.
\item
  (2024). Wax Deposition during the Transportation of Waxy Crude Oil: Mechanisms, Influencing Factors, Modeling, and Outlook. Energy \& Fuels. DOI: 10.1021/acs.energyfuels.3c04687.
\item
  (2022). Crude oil wax: A review on formation, experimentation, prediction, and remediation techniques. Petroleum Science, 19, 2214-2237. DOI: 10.1016/j.petsci.2022.08.008.
\item
  (2020). Rheology of waxy crude oils in relation to restart of gelled pipelines. Chemical Engineering Science. DOI: 10.1016/j.ces.2019.115270.
\item
  (2008). Time-Dependent Rheology of a Model Waxy Crude Oil with Relevance to Gelled Pipeline Restart. Energy \& Fuels. DOI: 10.1021/ef800628g.
\item
  (2023). Development of a New Model for the Formation of Wax Deposits through the Passage of Crude Oil within the Well. Sustainability, 15(12), 9616. DOI: 10.3390/su15129616.
\item
  Weingarten, J. S., Euchner, Julie. (1988). Methods for Predicting Wax Precipitation and Deposition. SPE Production Engineering, 3, 121-126. DOI: 10.2118/15654-pa.
\item
  Whitson, Curtis H. (1983). Characterizing Hydrocarbon Plus Fractions. Society of Petroleum Engineers Journal, 23, 683-694. DOI: 10.2118/12233-pa.
\item
  Won, K.W. (1986). Thermodynamics for solid solution-liquid-vapor equilibria: wax phase formation from heavy hydrocarbon mixtures. Fluid Phase Equilibria, 30, 265-279. DOI: 10.1016/0378-3812(86)80061-9.
\item
  Wu, Jianzhong, Prausnitz, John M., Firoozabadi, Abbas. (1998). Molecular‐thermodynamic framework for asphaltene‐oil equilibria. AIChE Journal, 44, 1188-1199. DOI: 10.1002/aic.690440516.
\item
  Wu, Jianzhong, Prausnitz, John M., Firoozabadi, Abbas. (2000). Molecular thermodynamics of asphaltene precipitation in reservoir fluids. AIChE Journal, 46, 197-209. DOI: 10.1002/aic.690460120.
\end{enumerate}

\end{document}
